\begin{titlepage}\pagecolor{navy}\color{white}\vspace*{1.1in}{\Huge\bfseries AI Engineering\\Interview Atlas\par}\vspace{.4in}{\Large Systems, theory, and 1,000+ questions\par}\vfill{\large Research snapshot: 2 September 2026\par}\vspace{.15in}{\normalsize Primary papers, official documentation, specifications, and standards\par}\vspace{.7in}{\color{cyan}\rule{\textwidth}{3pt}}\end{titlepage}\nopagecolor\color{black}
\tableofcontents\newpage
\section{How to use this handbook}
This handbook contains 192 topic notes, 1920 generated practice questions, 50 core interview questions, and 12 progressive system-design challenges. Study mechanisms first, state tradeoffs second, then describe measurement, failure handling, and rollback.
\subsection{Curriculum architecture}
\begin{center}\begin{tikzpicture}[node distance=9mm and 12mm,box/.style={draw=cyan,fill=paper,rounded corners,align=center,minimum width=27mm,minimum height=9mm},arr/.style={-{Stealth},thick,draw=blue}]\node[box](f){Foundations};\node[box,right=of f](r){Retrieval\\and RAG};\node[box,right=of r](a){Agents\\and control};\node[box,below=of f](t){Transformers};\node[box,right=of t](s){Serving\\and LLMOps};\node[box,right=of s](o){Observability\\and safety};\draw[arr](f)--(r);\draw[arr](r)--(a);\draw[arr](f)--(t);\draw[arr](t)--(s);\draw[arr](s)--(o);\draw[arr](a)--(o);\end{tikzpicture}\end{center}
\section{Retrieval and Search Theory}
\subsection{BM25}
A probabilistic lexical ranker that rewards term frequency while saturating repeated terms and normalizing document length.
\tradeoff{It is fast, interpretable, and exact for rare tokens but cannot directly match paraphrases.}
\pitfall{Treating it as obsolete removes a strong baseline and hurts identifier-heavy queries.}
\subsection{TF-IDF}
A sparse weighting scheme that combines within-document frequency with inverse collection frequency.
\tradeoff{It is transparent and cheap but lacks BM25's term-frequency saturation and tuned length normalization.}
\pitfall{Comparing raw cosine scores across changing corpora without rebuilding IDF causes drift.}
\subsection{Dense retrieval}
A bi-encoder maps queries and passages into a shared embedding space for nearest-neighbor search.
\tradeoff{Semantic recall improves, while exact names, numbers, and out-of-domain language may regress.}
\pitfall{Evaluating only on in-domain examples hides lexical misses and embedding drift.}
\subsection{Sparse neural retrieval}
Learned sparse models retain token-aligned inverted-index retrieval while expanding or reweighting terms.
\tradeoff{They bridge lexical and semantic matching but increase index size and model cost.}
\pitfall{Unbounded expansion creates latency and storage blowups.}
\subsection{Late interaction}
ColBERT-style systems preserve token-level embeddings and aggregate fine-grained similarities at query time.
\tradeoff{Accuracy often beats single-vector retrieval, with a larger index and heavier scoring path.}
\pitfall{Treating a late-interaction index like one-vector-per-document underestimates capacity needs.}
\subsection{Cross-encoder reranking}
A model jointly encodes query and candidate text to assign a relevance score.
\tradeoff{It gives strong precision on a small candidate set but is too expensive for full-corpus retrieval.}
\pitfall{Reranking too many long passages dominates p95 latency.}
\subsection{Hybrid retrieval}
Dense and lexical signals are retrieved together and fused, often with weighted scores or reciprocal-rank fusion.
\tradeoff{It improves robustness across semantic and exact-match queries but requires score calibration and duplicate handling.}
\pitfall{Adding incomparable raw scores without normalization makes one retriever dominate.}
\subsection{Reciprocal rank fusion}
RRF combines ranked lists using reciprocal rank positions rather than raw score scales.
\tradeoff{It is robust without score calibration but ignores the magnitude and confidence of scores.}
\pitfall{Using an overly small rank constant can overreward noisy top results.}
\subsection{ANN recall}
Approximate nearest-neighbor recall measures how many true nearest items are recovered by an index configuration.
\tradeoff{Higher recall usually costs memory, build time, or query latency.}
\pitfall{Reporting application answer quality without measuring index recall hides retrieval infrastructure defects.}
\subsection{Retrieval metrics}
Recall@k, precision@k, MRR, MAP, and nDCG measure different ranking behaviors.
\tradeoff{A metric must match the product objective and judgment granularity.}
\pitfall{Optimizing MRR for a task that consumes many passages can worsen context coverage.}
\subsection{Embedding geometry}
Similarity functions, normalization, anisotropy, dimensionality, and training objectives shape neighborhood quality.
\tradeoff{Cosine, dot product, and Euclidean distance can be equivalent only under specific normalization assumptions.}
\pitfall{Mixing an index metric with a differently trained embedding objective silently harms ranking.}
\subsection{Query expansion}
Expansion adds synonyms, generated subqueries, or pseudo-relevance terms before retrieval.
\tradeoff{Recall can improve at the expense of latency and topic drift.}
\pitfall{Letting an LLM expand unconstrained queries can inject unsupported intent.}
\section{Vector Databases and Search Engines}
\subsection{pgvector}
A PostgreSQL extension that stores vectors beside relational data and supports exact, HNSW, and IVFFlat search.
\tradeoff{It simplifies transactions and joins, while specialized stores may scale independent vector workloads more easily.}
\pitfall{Adding an ANN index without tuning filtering and planner behavior can reduce recall or miss the index.}
\subsection{Pinecone}
A managed vector service with metadata filtering, namespaces, and dense-sparse hybrid retrieval.
\tradeoff{Operations are outsourced, trading infrastructure control and portability for managed scale.}
\pitfall{Poor namespace and metadata design creates noisy tenants and expensive fan-out queries.}
\subsection{Qdrant}
An open-source vector engine centered on HNSW, payload filtering, quantization, and distributed collections.
\tradeoff{It offers operational control and rich filtering but adds a stateful distributed service to run.}
\pitfall{Applying post-filtering after ANN can return too few eligible results.}
\subsection{Weaviate}
A vector database with collections, hybrid search, filters, modules, and multi-tenancy features.
\tradeoff{An integrated schema accelerates product work but may couple ingestion and inference choices.}
\pitfall{Automatic vectorization without explicit model version metadata makes reindexing unsafe.}
\subsection{Milvus and Zilliz}
Milvus is a distributed vector database; Zilliz provides its managed service with multiple index choices.
\tradeoff{It targets large-scale separation of compute and storage, with more operational concepts than embedded options.}
\pitfall{Selecting an index by popularity instead of dataset size and latency-recall tests wastes resources.}
\subsection{Elasticsearch vector search}
Elasticsearch combines inverted indexes, filters, aggregations, and approximate kNN in one search platform.
\tradeoff{It is compelling for existing Elastic estates but vector memory and shard design require care.}
\pitfall{Oversharding fragments candidate pools and increases coordination overhead.}
\subsection{OpenSearch vector search}
OpenSearch supports lexical and vector workloads with pluggable kNN engines and search pipelines.
\tradeoff{Open governance and integrated search are attractive, while engine-specific behavior complicates tuning.}
\pitfall{Assuming identical parameters across FAISS, Lucene, and NMSLIB backends produces misleading tests.}
\subsection{Redis vector search}
Redis Query Engine adds vector indexes and filters to an in-memory data platform.
\tradeoff{Very low latency is possible, but RAM cost and persistence design matter at scale.}
\pitfall{Using Redis only as a vector database without capacity planning can evict critical data.}
\subsection{Chroma}
Chroma is a developer-oriented embedding store commonly used for local prototypes and small applications.
\tradeoff{It favors simplicity and iteration speed over complex distributed operations.}
\pitfall{Promoting a local prototype topology directly to production skips backup, tenancy, and load tests.}
\subsection{LanceDB}
LanceDB uses a columnar, versioned data format and supports vector, full-text, and multimodal data locally or in cloud storage.
\tradeoff{It fits data-lake and embedded workflows but differs operationally from a server-first database.}
\pitfall{Ignoring compaction and version retention can degrade read performance and inflate storage.}
\subsection{Vespa}
Vespa combines structured filtering, lexical search, tensors, ranking functions, and online serving.
\tradeoff{It enables sophisticated multi-stage ranking but has a steeper schema and operations learning curve.}
\pitfall{Putting expensive ranking expressions in the first phase can collapse throughput.}
\subsection{HNSW}
A navigable multilayer proximity graph controlled by construction degree and search breadth.
\tradeoff{It offers strong latency-recall behavior with significant memory and build cost.}
\pitfall{Deleting or heavily filtering data without maintenance degrades graph quality.}
\subsection{IVF and PQ}
Inverted files narrow the search to clusters; product quantization compresses vectors for cheaper approximate scoring.
\tradeoff{Memory and latency fall, while training quality and quantization error reduce recall.}
\pitfall{Training centroids or codebooks on an unrepresentative sample bakes bias into the index.}
\subsection{Metadata filtering}
Attribute constraints restrict eligible vectors before, during, or after ANN search.
\tradeoff{Early filtering preserves relevance but can make graph traversal sparse; post-filtering may underfill results.}
\pitfall{Failing to test selective filters separately from unfiltered queries hides severe recall loss.}
\section{RAG Architecture and Evaluation}
\subsection{Document ingestion}
Ingestion parses, normalizes, classifies, deduplicates, and records provenance before indexing.
\tradeoff{Rich preprocessing improves retrieval but increases pipeline cost and reprocessing complexity.}
\pitfall{Losing stable document and chunk identifiers makes updates create duplicates.}
\subsection{Chunking}
Chunking chooses semantic units, size, overlap, and metadata boundaries for retrieval and context assembly.
\tradeoff{Smaller chunks improve precision; larger chunks preserve context and reduce fragmentation.}
\pitfall{Fixed-size splitting across tables, code, or headings destroys meaning.}
\subsection{Parent-child retrieval}
Small child chunks retrieve precisely while larger parent sections are supplied to the model.
\tradeoff{It separates retrieval granularity from generation context at extra storage and mapping cost.}
\pitfall{Returning many overlapping parents wastes context and amplifies one source.}
\subsection{Contextual compression}
A post-retrieval step extracts or rewrites only query-relevant evidence from candidates.
\tradeoff{It saves context tokens but can delete qualifiers or introduce summarization errors.}
\pitfall{Using the same weak model for compression and answering creates correlated failures.}
\subsection{Query routing}
A router chooses retrievers, indexes, tools, or no-retrieval paths based on intent and risk.
\tradeoff{Specialization improves quality and cost but adds classification errors and operational branches.}
\pitfall{A route without confidence or fallback turns ambiguous queries into hard failures.}
\subsection{Multi-query retrieval}
The system generates several query formulations and merges their result sets.
\tradeoff{Recall rises at the cost of additional searches, duplicates, and possible intent drift.}
\pitfall{Treating all generated queries as equally trustworthy invites off-topic evidence.}
\subsection{HyDE}
Hypothetical Document Embeddings generate a likely answer passage and embed it as the retrieval query.
\tradeoff{It can bridge terse questions to document style but may anchor retrieval on invented details.}
\pitfall{Passing the hypothetical text to generation as evidence confuses a retrieval aid with a source.}
\subsection{Reranking pipeline}
A staged ranker retrieves broadly, applies filters, reranks candidates, and packs final context.
\tradeoff{More stages improve control but add latency, failure modes, and tuning surfaces.}
\pitfall{Evaluating only final answers makes it hard to locate which stage regressed.}
\subsection{Context packing}
Packing selects, orders, deduplicates, and labels evidence within a token budget.
\tradeoff{Diversity and source coverage compete with local relevance and continuity.}
\pitfall{Greedy top-score packing often repeats one document and crowds out complementary facts.}
\subsection{Grounded generation}
Prompts and decoding require answers to use cited evidence and abstain when evidence is insufficient.
\tradeoff{Strong grounding can reduce creativity and answer rate while improving auditability.}
\pitfall{A citation marker is not proof that the cited text entails the claim.}
\subsection{RAG evaluation}
Retrieval, context, generation, citation, latency, and cost require separate offline and online measures.
\tradeoff{Automated judges accelerate iteration but need human calibration and bias checks.}
\pitfall{Collapsing all dimensions into one score hides compensating failures.}
\subsection{Faithfulness}
Faithfulness asks whether claims in an answer are supported by the supplied context.
\tradeoff{It is narrower and more testable than global factuality but depends on context quality.}
\pitfall{Marking unsupported common knowledge as grounded inflates results.}
\subsection{Answer relevance}
Answer relevance measures whether a response directly addresses the user question without unnecessary content.
\tradeoff{Concision may improve relevance while losing useful caveats.}
\pitfall{A fluent on-topic answer can still be factually wrong, so relevance cannot stand alone.}
\subsection{Synthetic test generation}
Models generate questions, expected evidence, and variants from a source corpus to expand evaluation coverage.
\tradeoff{Coverage grows cheaply, while generator bias and leakage can make the benchmark too easy.}
\pitfall{Testing with the same model family that authored the cases can overstate quality.}
\subsection{Online RAG monitoring}
Production monitoring tracks retrieval gaps, unsupported claims, drift, feedback, latency, and cost by slice.
\tradeoff{Rich telemetry improves diagnosis but raises privacy and storage obligations.}
\pitfall{Logging full prompts and context without redaction can expose sensitive data.}
\section{Knowledge Graphs and GraphRAG}
\subsection{Property graphs}
Nodes, typed relationships, and properties represent entities and their connections for traversal and pattern matching.
\tradeoff{They are intuitive for operational queries but semantics depend on disciplined labels and schema.}
\pitfall{Using untyped generic edges turns the graph into an expensive document store.}
\subsection{RDF and ontologies}
RDF triples plus vocabularies such as RDFS or OWL provide globally identified semantics and inference rules.
\tradeoff{Formal interoperability grows, with more modeling and reasoning complexity.}
\pitfall{Confusing open-world semantics with database constraints produces invalid assumptions.}
\subsection{Cypher}
Cypher expresses graph patterns, predicates, paths, updates, and aggregation for property graphs.
\tradeoff{Declarative queries are readable, while unbounded path patterns can be expensive.}
\pitfall{Omitting labels or relationship types can trigger broad scans.}
\subsection{Knowledge graph construction}
Construction extracts entities and relations, resolves identities, applies schema, validates facts, and records provenance.
\tradeoff{Automation scales coverage but introduces extraction and entity-resolution errors.}
\pitfall{Allowing generated edges without source spans prevents later verification.}
\subsection{Entity resolution}
Entity resolution links mentions and records that refer to the same real-world entity.
\tradeoff{Aggressive merging improves connectivity but risks catastrophic false joins.}
\pitfall{Matching on names alone fails for aliases, homonyms, and temporal identity changes.}
\subsection{GraphRAG local search}
Local search expands from query-relevant entities through nearby relationships and associated text.
\tradeoff{It supports entity-centric multi-hop questions but depends on extraction and linking quality.}
\pitfall{Unbounded neighborhood expansion floods context with high-degree hubs.}
\subsection{GraphRAG global search}
Global search summarizes graph communities and combines relevant community reports for corpus-level questions.
\tradeoff{It addresses whole-corpus themes but has expensive indexing and lossy summaries.}
\pitfall{Expecting global search to answer precise record lookups wastes cost and may lose detail.}
\subsection{KAG}
Knowledge Augmented Generation combines schema-aware graphs, chunk mutual indexing, alignment, and logical-form-guided reasoning.
\tradeoff{It improves domain logic and multi-hop control at substantial modeling and maintenance cost.}
\pitfall{Adopting KAG without stable domain schema simply relocates ambiguity.}
\subsection{Graph quality evaluation}
Evaluation checks entity precision, relation correctness, completeness, consistency, provenance, and downstream QA value.
\tradeoff{Intrinsic graph metrics catch construction defects but may not predict application utility.}
\pitfall{Reporting only downstream answer scores can hide a brittle graph with unverifiable facts.}
\subsection{Temporal knowledge graphs}
Facts carry validity intervals, event time, and version history so queries can reason about change.
\tradeoff{Temporal accuracy improves, while updates, conflicts, and query logic become harder.}
\pitfall{Overwriting old facts destroys the ability to answer as-of questions.}
\section{Agents, MCP, and Control}
\subsection{Agent loop}
An agent repeatedly observes state, chooses an action, invokes a tool, and incorporates the result until stopping.
\tradeoff{Flexibility handles open-ended tasks, but loop length and behavior are less predictable than workflows.}
\pitfall{Omitting hard budgets and stop conditions creates runaway calls.}
\subsection{ReAct}
ReAct interleaves reasoning with actions and observations to update plans using external evidence.
\tradeoff{Tool feedback can correct model knowledge, while longer trajectories add latency and error opportunities.}
\pitfall{Treating every thought as reliable state allows early mistakes to compound.}
\subsection{Planning}
Planning decomposes goals into ordered, conditional, or revisable steps before or during execution.
\tradeoff{Upfront plans improve structure but become stale in uncertain environments.}
\pitfall{Forcing a long immutable plan prevents recovery after new observations.}
\subsection{Reflection}
Reflection critiques outcomes or trajectories and stores actionable lessons for another attempt.
\tradeoff{It can improve recovery without weight updates, while self-critique may repeat the same bias.}
\pitfall{Saving vague reflections adds context without changing behavior.}
\subsection{Tool calling}
Models emit structured arguments for declared functions that a host validates and executes.
\tradeoff{It narrows the action interface but does not make proposed actions safe or correct.}
\pitfall{Executing model-produced arguments without schema and authorization checks invites damage.}
\subsection{Model Context Protocol}
MCP standardizes client-server discovery and invocation of tools, resources, prompts, and related capabilities.
\tradeoff{Interoperability improves, while trust boundaries, consent, and capability scoping remain host responsibilities.}
\pitfall{Treating a discovered MCP server as trusted creates prompt-injection and privilege risks.}
\subsection{Agent memory}
Working, episodic, semantic, and procedural memories preserve relevant state beyond a single model call.
\tradeoff{Memory supports continuity but retrieval, deletion, privacy, and stale facts become product concerns.}
\pitfall{Writing every interaction to long-term memory creates noise and privacy exposure.}
\subsection{Human in the loop}
High-risk actions pause for review, edit, approval, rejection, or direct human response.
\tradeoff{Oversight reduces irreversible mistakes but adds latency and reviewer load.}
\pitfall{Asking for approval after an action has already caused side effects is theater.}
\subsection{State machines}
Explicit states and transitions constrain allowed agent behavior and make recovery inspectable.
\tradeoff{Predictability and testability improve at the cost of less open-ended autonomy.}
\pitfall{Encoding side effects inside transition guards breaks replay safety.}
\subsection{Idempotent tools}
An idempotent tool can be retried with the same operation key without duplicating the intended effect.
\tradeoff{Reliable retries require deduplication state and carefully defined semantics.}
\pitfall{Assuming HTTP method alone guarantees business idempotency causes duplicate orders or messages.}
\subsection{Agent handoffs}
A handoff transfers a task and structured context to a specialized agent or deterministic service.
\tradeoff{Specialization reduces prompt overload but adds routing and ownership ambiguity.}
\pitfall{Passing the entire transcript leaks irrelevant data and confuses the recipient.}
\subsection{Autonomy levels}
Autonomy is graduated from suggest-only through approval-gated execution to bounded autonomous operation.
\tradeoff{More autonomy reduces human effort while increasing required assurance and containment.}
\pitfall{Assigning one global autonomy level ignores action-specific risk.}
\subsection{Context engineering}
Context engineering selects instructions, memory, evidence, tools, schemas, and state for each model turn.
\tradeoff{Targeted context improves reliability and cost, but selection logic becomes core application code.}
\pitfall{Appending everything causes attention dilution and instruction conflicts.}
\subsection{Sandbox runs}
A sandbox isolates untrusted code with resource limits, scoped files, network policy, timeouts, and audit logs.
\tradeoff{Isolation enables capability while imposing startup and compatibility overhead.}
\pitfall{A container without kernel, credential, and network boundaries is not a sufficient sandbox.}
\section{Orchestration Frameworks}
\subsection{LangGraph}
LangGraph models long-running agent workflows as stateful graphs with checkpoints, interrupts, streaming, and durable execution.
\tradeoff{Fine-grained control aids production reliability but exposes more state and graph design decisions.}
\pitfall{Non-idempotent node side effects make checkpoint replay unsafe.}
\subsection{LlamaIndex workflows}
Event-driven workflows coordinate retrieval, tools, agents, and structured steps in the LlamaIndex ecosystem.
\tradeoff{Strong data and indexing integrations speed RAG work but can couple application flow to framework abstractions.}
\pitfall{Hiding all retrieval behind defaults makes evaluation and migration difficult.}
\subsection{Haystack pipelines}
Haystack connects typed components into retrieval, generation, routing, and agent pipelines.
\tradeoff{Explicit components support testing, while custom graphs still require lifecycle discipline.}
\pitfall{Allowing incompatible document schemas across components creates runtime surprises.}
\subsection{PydanticAI}
PydanticAI uses typed dependencies, tools, results, and validation to build model-agnostic Python agents.
\tradeoff{Type safety improves integration but cannot validate semantic truth.}
\pitfall{Confusing schema conformance with a correct answer produces false confidence.}
\subsection{Semantic Kernel}
Semantic Kernel offers planners, agents, plugins, memory connectors, and multi-language enterprise integration.
\tradeoff{Its broad integration surface suits platform teams but introduces framework concepts and version coupling.}
\pitfall{Registering overly broad plugins grants the model unnecessary authority.}
\subsection{AutoGen}
AutoGen structures multi-agent conversations and tool-using participants for collaborative task execution.
\tradeoff{Role decomposition enables complex workflows but multiplies calls, coordination failures, and evaluation cost.}
\pitfall{Adding agents without independent capabilities creates expensive echo chambers.}
\subsection{CrewAI}
CrewAI organizes role-based agents, tasks, crews, and flows for sequential or coordinated work.
\tradeoff{It gives accessible orchestration but still needs explicit state, permissions, and observability.}
\pitfall{Persona descriptions cannot substitute for tool-level access control.}
\subsection{DSPy}
DSPy treats prompts and modules as optimizable programs evaluated against examples and metrics.
\tradeoff{Systematic optimization can beat manual prompting, while it depends on representative data and stable metrics.}
\pitfall{Optimizing against a leaky judge overfits wording instead of capability.}
\subsection{Temporal}
Temporal persists workflow history and replays deterministic workflow code around durable activities.
\tradeoff{It provides robust retries and long-running execution but requires deterministic workflow discipline.}
\pitfall{Reading wall-clock time or random state directly in replayed code causes nondeterminism.}
\subsection{Airflow, Dagster, and Prefect}
General-purpose data orchestrators schedule DAGs, assets, or flows with retries, lineage, and operational UIs.
\tradeoff{They excel at batch pipelines but are not automatically low-latency interactive agent runtimes.}
\pitfall{Running per-token or per-chat-turn work as heavy batch tasks adds unacceptable overhead.}
\section{Observability and Monitoring}
\subsection{OpenTelemetry}
OpenTelemetry provides vendor-neutral APIs, SDKs, collectors, and semantic conventions for traces, metrics, logs, and profiles.
\tradeoff{Portability improves, while instrumentation quality and backend cost remain engineering work.}
\pitfall{High-cardinality prompt or user attributes can overwhelm telemetry systems.}
\subsection{Distributed tracing}
Traces connect request spans across gateways, retrievers, models, tools, queues, and databases.
\tradeoff{They reveal critical paths but sampling can hide rare quality failures.}
\pitfall{Reusing trace identifiers across independent user requests breaks isolation.}
\subsection{LLM tracing}
LLM traces capture prompts, outputs, model settings, token usage, latency, tool calls, retrieval, and evaluations.
\tradeoff{Debugging improves, while content capture raises privacy and retention risk.}
\pitfall{Logging secrets or personal data by default creates a second sensitive datastore.}
\subsection{Metrics and SLOs}
Counters, histograms, gauges, service-level indicators, and objectives quantify reliability and performance.
\tradeoff{Aggregate signals scale well but cannot explain individual semantic failures.}
\pitfall{Averaging latency hides long-tail p95 and p99 user pain.}
\subsection{Structured logging}
Structured events use stable fields, severity, timestamps, correlation IDs, and redaction.
\tradeoff{Search and aggregation improve, while schemas must evolve consistently.}
\pitfall{Placing unbounded model output in one log field causes cost and ingestion failures.}
\subsection{LangSmith}
LangSmith provides tracing, datasets, evaluation, prompt management, and deployment-oriented tooling for LLM applications.
\tradeoff{Tight LangChain integration accelerates work but requires portability planning.}
\pitfall{Depending only on vendor UI annotations leaves evaluation logic unreproducible.}
\subsection{Langfuse}
Langfuse is an open-source platform for LLM traces, sessions, prompts, scores, datasets, and cost analytics.
\tradeoff{Self-hosting offers control but shifts scaling, upgrades, and security to the team.}
\pitfall{Capturing every token forever creates expensive sensitive retention.}
\subsection{Arize Phoenix}
Phoenix provides open-source tracing and evaluation for LLM, retrieval, and agent applications with OpenTelemetry support.
\tradeoff{It unifies qualitative traces and quantitative evals, while production operations still need design.}
\pitfall{Judge scores without slice analysis can mask failures for important cohorts.}
\subsection{MLflow tracing}
MLflow connects traces, evaluation datasets, scorers, experiments, and production monitoring.
\tradeoff{A common ML platform improves lineage but may be heavier than a focused tracing tool.}
\pitfall{Mixing incomparable scorer versions in one trend chart produces false regressions.}
\subsection{W\&B Weave}
Weave tracks calls, objects, datasets, evaluations, and model application experiments.
\tradeoff{Rich experiment lineage helps collaboration but introduces another governance surface.}
\pitfall{Saving mutable datasets without immutable versions undermines reproducibility.}
\subsection{Prompt and model drift}
Drift monitoring detects changes in inputs, retrieved evidence, outputs, model versions, and user behavior.
\tradeoff{Alerts catch regressions but thresholds must distinguish seasonality from harm.}
\pitfall{Measuring embedding distribution shift alone does not prove answer quality changed.}
\subsection{Cost observability}
Cost telemetry attributes input, output, cache, retrieval, tool, and infrastructure spend to requests and tenants.
\tradeoff{Unit economics become actionable, while provider pricing and caching rules evolve.}
\pitfall{Tracking only average cost hides pathological long-running agents.}
\subsection{Evaluation in production}
Online evaluators combine sampled judges, deterministic checks, human feedback, and business outcomes.
\tradeoff{Live coverage improves realism but evaluations must not expose or delay users.}
\pitfall{Treating implicit clicks as unambiguous quality labels creates bias.}
\section{Serving, Deployment, and LLMOps}
\subsection{vLLM}
vLLM is a high-throughput inference server using PagedAttention, continuous batching, parallelism, and OpenAI-compatible APIs.
\tradeoff{Throughput and memory utilization improve, while model and kernel compatibility must be verified.}
\pitfall{Benchmarking only one concurrency level hides latency-throughput collapse points.}
\subsection{SGLang}
SGLang combines a structured generation runtime with optimized model serving, prefix reuse, batching, and parallelism.
\tradeoff{Tight runtime-serving integration can accelerate agentic workloads but changes application abstractions.}
\pitfall{Assuming every structured program benefits equally from cache reuse leads to poor capacity models.}
\subsection{Text Generation Inference}
TGI is a Hugging Face server for transformer inference with streaming, batching, quantization, and distributed support.
\tradeoff{Ecosystem integration is strong, while supported optimization paths vary by model and hardware.}
\pitfall{Enabling quantization without quality and kernel tests can reduce both accuracy and speed.}
\subsection{TensorRT-LLM}
TensorRT-LLM compiles and serves NVIDIA-optimized LLM graphs and kernels across GPUs.
\tradeoff{Maximum NVIDIA performance comes with build complexity and hardware coupling.}
\pitfall{Rebuilding engines for every shape without a deployment cache slows releases.}
\subsection{NVIDIA Triton}
Triton serves multiple model frameworks with dynamic batching, ensembles, metrics, and model repositories.
\tradeoff{It standardizes diverse serving, though LLM-specific scheduling may require specialized backends.}
\pitfall{Using dynamic batching without queue-delay limits breaks interactive latency.}
\subsection{llama.cpp}
llama.cpp runs quantized models efficiently on CPUs and varied local accelerators through a portable C++ stack.
\tradeoff{Edge privacy and low dependency count trade off against datacenter throughput and model size.}
\pitfall{Selecting a quantization by file size alone ignores accuracy and hardware kernels.}
\subsection{KServe}
KServe provides Kubernetes-native inference services, autoscaling, revisions, traffic splitting, and model runtimes.
\tradeoff{Platform consistency improves, with Kubernetes operational overhead and cold-start concerns.}
\pitfall{Scaling only on CPU misses token queue pressure in GPU workloads.}
\subsection{Ray Serve}
Ray Serve composes Python model deployments and scales replicas across a Ray cluster.
\tradeoff{Flexible distributed Python pipelines fit multi-model systems but require cluster and object-store discipline.}
\pitfall{Passing large tensors repeatedly through serialization paths wastes memory and latency.}
\subsection{BentoML}
BentoML packages model services, dependencies, APIs, and deployment configuration with adaptive batching.
\tradeoff{Developer workflow is streamlined, while platform integration and runtime tuning still matter.}
\pitfall{Treating packaging success as load readiness skips concurrency tests.}
\subsection{Kubernetes and GPUs}
Kubernetes schedules containers and device resources while operators add GPU sharing, queues, and autoscaling.
\tradeoff{A common control plane helps governance but default scheduling is not token-aware.}
\pitfall{Requesting whole GPUs for tiny models can strand capacity.}
\subsection{LLM gateways}
Gateways normalize provider APIs and add routing, budgets, rate limits, fallbacks, caching, and audit policy.
\tradeoff{Central control improves governance but becomes a high-impact availability and security dependency.}
\pitfall{Automatic fallback across non-equivalent models can silently change behavior.}
\subsection{Canary and shadow releases}
Canary traffic exposes a new version to a small cohort; shadowing duplicates requests without serving the new output.
\tradeoff{Risk falls, while cost, comparison logic, and privacy obligations grow.}
\pitfall{Comparing canaries with different traffic mixes confounds the result.}
\subsection{Model registry}
A registry versions model artifacts, lineage, signatures, evaluation evidence, aliases, and promotion state.
\tradeoff{Governance and rollback improve but approval gates can slow iteration.}
\pitfall{Using mutable stage names without immutable digests makes rollbacks uncertain.}
\subsection{Autoscaling for LLMs}
LLM autoscaling should consider queue depth, time-to-first-token, tokens per second, memory, and startup time.
\tradeoff{Extra replicas reduce latency but GPUs are costly and slow to warm.}
\pitfall{CPU-based autoscaling reacts late or incorrectly to GPU saturation.}
\section{Data and ML Lifecycle}
\subsection{DVC}
DVC versions large data and model pointers with Git and defines reproducible pipeline stages and experiments.
\tradeoff{Familiar Git workflows improve reproducibility, while remote storage and cache discipline are required.}
\pitfall{Committing only code while silently replacing remote data breaks lineage.}
\subsection{lakeFS}
lakeFS adds Git-like branches, commits, merges, and hooks over object-storage data lakes.
\tradeoff{Atomic data-lake changes enable testing, with server and metadata operations to manage.}
\pitfall{Treating object versioning alone as branch-level data semantics is insufficient.}
\subsection{Delta Lake and Iceberg}
Open table formats add schemas, snapshots, partition evolution, and transactional metadata over object storage.
\tradeoff{Reliable analytics improve, while compaction and catalog operations become essential.}
\pitfall{Excess small files destroy scan performance despite correct table semantics.}
\subsection{Feature stores}
Feature stores define reusable features and connect historical offline data with low-latency online serving.
\tradeoff{Training-serving consistency improves but adds infrastructure and ownership requirements.}
\pitfall{Recomputing features differently online creates skew.}
\subsection{Point-in-time correctness}
Historical joins select only feature values known at each label timestamp.
\tradeoff{Leakage is prevented at the cost of temporal keys and more expensive joins.}
\pitfall{Joining the latest feature value into old examples creates unrealistically strong training results.}
\subsection{Experiment tracking}
Runs capture code, parameters, data, metrics, artifacts, environment, and parent relationships.
\tradeoff{Comparison and reproducibility improve but inconsistent naming creates a junk drawer.}
\pitfall{Logging a metric without dataset and code versions makes it uninterpretable.}
\subsection{Data contracts}
Contracts specify schema, semantics, freshness, ownership, quality, and compatibility expectations between producers and consumers.
\tradeoff{Clear boundaries reduce silent breakage but require governance and enforcement.}
\pitfall{Validating shape without semantic definitions misses unit and meaning changes.}
\subsection{Data quality testing}
Tests cover schema, nulls, ranges, uniqueness, distribution, freshness, and cross-table invariants.
\tradeoff{Early detection prevents downstream damage, while brittle thresholds create alert fatigue.}
\pitfall{Checking only aggregate distributions misses failures in important slices.}
\subsection{Data lineage}
Lineage links sources, transformations, datasets, features, prompts, models, evaluations, and deployments.
\tradeoff{Impact analysis and auditability improve, while automated lineage can miss dynamic behavior.}
\pitfall{Lineage without immutable identifiers gives a false sense of reproducibility.}
\subsection{Reproducible environments}
Lockfiles, containers, hardware metadata, seeds, and deterministic settings constrain run variability.
\tradeoff{Reproducibility improves at the expense of build maintenance and sometimes performance.}
\pitfall{A random seed alone cannot guarantee deterministic GPU execution.}
\section{Transformer Theory and Training}
\subsection{Tokenization}
Tokenizers map text or bytes into model vocabulary IDs using algorithms such as BPE, WordPiece, or unigram models.
\tradeoff{Larger vocabularies shorten sequences but increase embedding parameters and sparse coverage.}
\pitfall{Comparing context lengths across tokenizers as if tokens were equal misstates capacity and cost.}
\subsection{Embedding layer}
Learned token vectors map discrete IDs into continuous model space and are often tied to output weights.
\tradeoff{Weight tying saves parameters but constrains input-output geometry.}
\pitfall{Forgetting vocabulary changes invalidates checkpoint embedding shapes.}
\subsection{Self-attention math}
Attention computes scaled query-key scores, applies masking and softmax, then mixes value vectors.
\tradeoff{Every token can interact directly, while standard sequence cost grows quadratically.}
\pitfall{Omitting the square-root head-dimension scale saturates softmax gradients.}
\subsection{Multi-head attention}
Multiple heads project attention into separate subspaces before concatenation and output projection.
\tradeoff{Diverse interactions improve capacity but add projections and KV memory.}
\pitfall{Assuming heads are independently interpretable can overstate mechanistic conclusions.}
\subsection{Positional encoding}
Position signals break permutation invariance using absolute, relative, rotary, or bias-based schemes.
\tradeoff{Relative methods extrapolate differently, while implementation and scaling choices affect long context.}
\pitfall{Extending context without validating the position scheme can collapse quality.}
\subsection{RoPE}
Rotary position embeddings rotate query and key coordinates so dot products encode relative offsets.
\tradeoff{RoPE integrates cleanly with attention but long-context scaling needs careful frequency treatment.}
\pitfall{Naively increasing the maximum length does not preserve trained frequency behavior.}
\subsection{Feed-forward networks}
Transformer MLP blocks expand hidden dimensions, apply nonlinear gates, and project back per token.
\tradeoff{They hold substantial parameters and compute while offering parallel token processing.}
\pitfall{Ignoring activation memory underestimates training capacity.}
\subsection{Layer normalization and RMSNorm}
Normalization stabilizes activations; RMSNorm removes mean centering and normalizes by root-mean-square magnitude.
\tradeoff{Simpler normalization can be faster but architecture placement changes optimization dynamics.}
\pitfall{Moving pre-norm to post-norm without retuning can destabilize deep training.}
\subsection{Residual connections}
Residual paths add transformed activations back to the stream to improve gradient flow and feature reuse.
\tradeoff{Deep optimization becomes feasible, while residual scale can grow and require controls.}
\pitfall{In-place mutation of residual tensors can break autograd or checkpointing.}
\subsection{Mixture of Experts}
MoE routes each token to a subset of expert networks, increasing parameters without proportional per-token compute.
\tradeoff{Capacity scales efficiently but routing balance, communication, and memory are hard.}
\pitfall{Expert collapse and uneven routing create stragglers and weak specialization.}
\subsection{Causal language modeling}
Next-token training factorizes sequence likelihood from left to right under a causal mask.
\tradeoff{It provides a simple scalable objective but does not directly optimize instruction usefulness or truth.}
\pitfall{Interpreting low perplexity as safe task performance is invalid.}
\subsection{Cross-entropy and perplexity}
Cross-entropy is negative log likelihood; perplexity exponentiates average token loss.
\tradeoff{It is stable for model comparison on the same tokenization and data, not across arbitrary setups.}
\pitfall{Comparing perplexity across vocabularies or tokenizers is misleading.}
\subsection{SFT}
Supervised fine-tuning teaches instruction-response behavior from curated demonstrations.
\tradeoff{It is direct and stable but can overfit style and inherit annotation errors.}
\pitfall{Training on conflicting instructions without provenance makes behavior unstable.}
\subsection{RLHF and PPO}
A preference model supplies rewards while PPO constrains policy updates relative to a reference model.
\tradeoff{It can optimize human preferences but adds reward hacking and training complexity.}
\pitfall{A biased reward model can produce confidently optimized failure modes.}
\subsection{DPO}
Direct Preference Optimization fits preferred over rejected responses through a classification-style objective relative to a reference policy.
\tradeoff{It simplifies preference training but still depends on pair quality and distribution coverage.}
\pitfall{Treating implicit user choices as clean preferences introduces confounding.}
\subsection{LoRA}
LoRA trains low-rank update matrices while freezing base weights for parameter-efficient adaptation.
\tradeoff{Memory and storage fall, while rank and target-module choices limit expressiveness.}
\pitfall{Merging incompatible adapters without evaluation creates interference.}
\subsection{QLoRA}
QLoRA backpropagates through a frozen quantized base model into LoRA adapters using memory-saving quantization techniques.
\tradeoff{Fine-tuning large models becomes cheaper, while kernels and quantization settings affect stability.}
\pitfall{Confusing training quantization with final serving quantization produces wrong expectations.}
\subsection{Knowledge distillation}
A student learns from teacher probabilities, generated data, rationales, or intermediate representations.
\tradeoff{Serving cost falls but rare capabilities and calibration may be lost.}
\pitfall{Distilling only easy teacher outputs creates a deceptively strong narrow model.}
\section{Inference Optimization}
\subsection{KV caching}
Autoregressive decoding stores prior keys and values so each new token avoids recomputing attention over the entire prefix.
\tradeoff{Decode speed improves while cache memory grows with layers, sequence, batch, and KV heads.}
\pitfall{Capacity planning only model weights causes out-of-memory failures at long context.}
\subsection{PagedAttention}
PagedAttention manages KV cache in non-contiguous blocks using paging-inspired allocation to reduce fragmentation and enable sharing.
\tradeoff{Utilization and batching improve, while block tables and kernels add runtime complexity.}
\pitfall{Oversized blocks waste tail space; tiny blocks add metadata overhead.}
\subsection{Continuous batching}
Requests join and leave a running decode batch at iteration boundaries rather than waiting for a fixed batch.
\tradeoff{Throughput and utilization rise, while scheduling fairness and latency isolation become harder.}
\pitfall{Optimizing total tokens per second alone can starve short requests.}
\subsection{Prefix caching}
Shared prompt prefixes reuse precomputed KV blocks across requests with matching tokens and model state.
\tradeoff{Repeated system prompts become cheaper, but hit rate, invalidation, and tenant isolation matter.}
\pitfall{Cache keys that omit adapters or model revision return incorrect state.}
\subsection{FlashAttention}
FlashAttention tiles exact attention to reduce high-bandwidth-memory traffic and materialization of the score matrix.
\tradeoff{Wall-clock speed and memory improve without approximate attention, but kernels depend on shapes and hardware.}
\pitfall{Calling it subquadratic compute confuses IO optimization with arithmetic complexity.}
\subsection{GQA and MQA}
Grouped-query and multi-query attention share key-value heads across query heads to shrink KV cache and memory bandwidth.
\tradeoff{Decode efficiency improves, with possible quality loss relative to full multi-head KV.}
\pitfall{Converting weights naively after training does not reproduce a model trained for shared KV heads.}
\subsection{Speculative decoding}
A cheaper draft proposes tokens and a target model verifies them while preserving the target distribution.
\tradeoff{Latency falls when acceptance is high, but drafting and verification overhead can dominate.}
\pitfall{Measuring draft speed without acceptance rate gives meaningless projections.}
\subsection{Quantization}
Quantization maps weights or activations to lower precision using scales, groups, calibration, or quantization-aware training.
\tradeoff{Memory and bandwidth fall, while accuracy and kernel support can degrade.}
\pitfall{A smaller model file may run slower if hardware lacks optimized kernels.}
\subsection{Tensor parallelism}
Tensor parallelism shards matrix operations across devices and communicates partial results within layers.
\tradeoff{It enables large models and lower per-request latency but requires high-bandwidth interconnects.}
\pitfall{Spanning slow network links can make communication dominate compute.}
\subsection{Pipeline parallelism}
Pipeline parallelism assigns layer stages to devices and streams microbatches through them.
\tradeoff{It scales model depth but creates bubbles and per-request latency.}
\pitfall{Too few microbatches leave stages idle.}
\subsection{Data parallel inference}
Replicas serve independent requests and scale throughput with routing and load balancing.
\tradeoff{It is operationally simple but duplicates model weights on every replica.}
\pitfall{Random routing ignores prefix-cache locality and uneven sequence lengths.}
\subsection{Expert parallelism}
MoE experts are distributed across devices and tokens are routed with all-to-all communication.
\tradeoff{Parameter capacity scales, while communication and load balance determine performance.}
\pitfall{Average expert load hides hot experts that gate step time.}
\subsection{Dynamic batching}
A server waits briefly to combine compatible requests into a batch.
\tradeoff{GPU efficiency improves while queue delay harms time-to-first-token.}
\pitfall{One global maximum delay cannot serve both interactive and batch workloads well.}
\subsection{Streaming generation}
Servers emit tokens incrementally with backpressure, disconnect handling, and final usage accounting.
\tradeoff{Perceived latency improves, but partial outputs complicate moderation and retries.}
\pitfall{Retrying after visible partial output can duplicate text or side effects.}
\section{Backend, APIs, and Microservices}
\subsection{FastAPI}
FastAPI builds typed ASGI APIs from Python hints with validation, dependency injection, and OpenAPI generation.
\tradeoff{Development is fast, while blocking work must be isolated from the event loop.}
\pitfall{Declaring a blocking database or model call inside async code stalls all requests.}
\subsection{Starlette}
Starlette is a lightweight ASGI toolkit providing routing, middleware, WebSockets, background tasks, and test support.
\tradeoff{It offers control with fewer batteries than a full API framework.}
\pitfall{Reimplementing validation and dependency patterns inconsistently increases defects.}
\subsection{Litestar}
Litestar is an ASGI framework with typing, dependency injection, OpenAPI, plugins, and data-transfer objects.
\tradeoff{Rich framework features reduce boilerplate but introduce conventions and migration cost.}
\pitfall{Choosing by benchmark alone ignores ecosystem and team familiarity.}
\subsection{Flask}
Flask is a minimal WSGI web framework with a large extension ecosystem.
\tradeoff{Simplicity suits synchronous services, while async-first streaming needs different architecture.}
\pitfall{Assuming an async view converts a WSGI deployment into high-concurrency ASGI is wrong.}
\subsection{Django and DRF}
Django supplies an ORM, admin, authentication, migrations, and conventions; DRF adds APIs.
\tradeoff{Batteries accelerate business systems but can be heavy for a narrow inference gateway.}
\pitfall{Performing N+1 ORM queries around an expensive model magnifies tail latency.}
\subsection{REST API design}
REST uses resource-oriented URLs, HTTP methods, status codes, caching, content negotiation, and idempotency semantics.
\tradeoff{Broad interoperability is strong, while chatty workflows and schema evolution need care.}
\pitfall{Returning 200 for every failure breaks clients and observability.}
\subsection{gRPC}
gRPC uses protobuf contracts and HTTP/2 for efficient unary and streaming RPCs with generated clients.
\tradeoff{Typed internal services perform well, while browser and human debugging are less direct.}
\pitfall{Changing field numbers or incompatible meanings breaks wire compatibility.}
\subsection{WebSockets and SSE}
WebSockets provide bidirectional frames; server-sent events provide simpler one-way HTTP streams.
\tradeoff{SSE fits token streaming and reconnect semantics; WebSockets fit interactive duplex protocols.}
\pitfall{Using WebSockets when one-way streaming is enough adds state and proxy complexity.}
\subsection{API idempotency}
Idempotency keys let retried mutation requests resolve to the same business operation and stored outcome.
\tradeoff{Duplicate effects are prevented at the cost of durable deduplication state.}
\pitfall{Expiring keys before upstream retry windows can reintroduce duplicates.}
\subsection{Rate limiting}
Token buckets, leaky buckets, and concurrency limits protect shared resources by identity and cost.
\tradeoff{Stability and fairness improve but strict limits can reject valuable bursts.}
\pitfall{Limiting requests instead of tokens lets a few long prompts monopolize capacity.}
\subsection{Circuit breakers}
Circuit breakers stop calls to a failing dependency, probe recovery, and combine with timeouts and fallbacks.
\tradeoff{Cascading failure is reduced, while bad thresholds can create synchronized outages.}
\pitfall{A breaker without bounded timeouts reacts only after resources are already exhausted.}
\subsection{Microservice boundaries}
Boundaries should follow ownership, data, scaling, failure isolation, and change cadence rather than nouns alone.
\tradeoff{Independent deployment helps teams but adds networks, consistency, and operations.}
\pitfall{Splitting too early creates a distributed monolith.}
\subsection{Event-driven architecture}
Producers publish immutable events to brokers while consumers process independently with delivery semantics.
\tradeoff{Decoupling and replay improve, while ordering, duplication, schemas, and lag must be managed.}
\pitfall{Assuming exactly-once business effects from at-least-once delivery causes duplicates.}
\subsection{Object-oriented design}
Encapsulation, composition, interfaces, and explicit invariants organize stateful behavior.
\tradeoff{Abstraction improves changeability, while deep inheritance hides coupling.}
\pitfall{Modeling every data record as a behavior-heavy class adds ceremony without invariants.}
\section{Python Engineering}
\subsection{Python threading and the GIL}
CPython threads share memory; the GIL limits simultaneous Python bytecode but releases around many blocking operations.
\tradeoff{Threads suit I/O-bound work and shared state, not pure Python CPU scaling.}
\pitfall{Claiming threads are always useless ignores network and native-extension concurrency.}
\subsection{Multiprocessing}
Separate processes bypass the GIL and isolate memory while communicating through serialization or shared memory.
\tradeoff{CPU parallelism improves at startup, memory, and coordination cost.}
\pitfall{Passing large models to spawned workers duplicates memory unexpectedly.}
\subsection{Asyncio}
Asyncio cooperatively schedules coroutines on an event loop around awaitable I/O.
\tradeoff{High connection concurrency is efficient, while blocking functions freeze the loop.}
\pitfall{Creating coroutines without awaiting them silently drops work.}
\subsection{Concurrent futures}
ThreadPoolExecutor and ProcessPoolExecutor expose a common future-based API for concurrent callables.
\tradeoff{Integration is simple, but cancellation and process serialization have limits.}
\pitfall{Submitting unbounded tasks creates memory pressure and queue latency.}
\subsection{Decorators}
Decorators wrap or replace functions or classes to add behavior while preserving a callable interface.
\tradeoff{Cross-cutting concerns become reusable, but hidden control flow complicates debugging.}
\pitfall{Omitting functools.wraps breaks metadata and framework introspection.}
\subsection{Context managers}
Context managers guarantee paired setup and cleanup through with, \_\_enter\_\_/\_\_exit\_\_, or contextlib.
\tradeoff{Resource safety improves and exceptions are localized.}
\pitfall{Swallowing exceptions unintentionally in \_\_exit\_\_ hides failures.}
\subsection{Generators}
Generators suspend and resume execution around yield, enabling lazy iteration and streaming pipelines.
\tradeoff{Memory use falls, while one-shot iteration and cleanup require care.}
\pitfall{Reusing an exhausted generator returns no data without warning.}
\subsection{Descriptors and properties}
Descriptors define attribute access; properties provide managed getter, setter, and deleter behavior.
\tradeoff{Validation and lazy access fit attribute syntax but may hide expensive operations.}
\pitfall{Performing network I/O in a property surprises callers and tooling.}
\subsection{Dataclasses and Pydantic}
Dataclasses reduce class boilerplate; Pydantic adds parsing, validation, serialization, and schema generation.
\tradeoff{Typed models improve boundaries but validation can add runtime cost.}
\pitfall{Treating parsed external input as semantically authorized is unsafe.}
\subsection{Type hints and protocols}
Type hints support static analysis; Protocol defines structural interfaces without inheritance.
\tradeoff{Refactoring and contracts improve, while runtime behavior is unchanged unless validated.}
\pitfall{Assuming type annotations enforce production inputs creates vulnerabilities.}
\subsection{Exceptions}
Exceptions separate error propagation from normal return values and can preserve causal chains.
\tradeoff{Central handling simplifies APIs, while overly broad catches erase actionable context.}
\pitfall{Catching Exception and returning None turns failures into corrupt state.}
\subsection{Memory management}
CPython combines reference counting with cyclic garbage collection; object graphs and native buffers affect real memory.
\tradeoff{Automatic cleanup is convenient, while cycles and caches can retain large objects.}
\pitfall{Watching only Python heap misses tensors allocated in native or GPU memory.}
\subsection{Testing and mocking}
Unit, integration, property, contract, and load tests cover different risks; mocks isolate only well-understood boundaries.
\tradeoff{Fast tests aid iteration, while excessive mocking tests implementation rather than behavior.}
\pitfall{Mocking the model, database, and network in the same test proves little about integration.}
\subsection{Packaging and environments}
pyproject.toml, lockfiles, wheels, virtual environments, and reproducible builds define deliverable Python software.
\tradeoff{Isolation prevents dependency drift but needs platform and supply-chain controls.}
\pitfall{Loose version ranges in production make rebuilds non-reproducible.}
\section{Security, Safety, and Governance}
\subsection{Prompt injection}
Untrusted content attempts to override instructions, exfiltrate data, or induce tool actions.
\tradeoff{Content isolation and least privilege reduce impact but cannot make untrusted text authoritative.}
\pitfall{Telling the model to ignore attacks is not a security boundary.}
\subsection{Tool authorization}
The host checks identity, scope, resource, action, policy, and user intent before every tool execution.
\tradeoff{Fine-grained authorization limits blast radius but adds policy complexity.}
\pitfall{Authorizing the agent once at session start enables confused-deputy attacks.}
\subsection{Secrets management}
Credentials are issued narrowly, stored outside prompts and code, rotated, audited, and redacted from telemetry.
\tradeoff{Short-lived credentials reduce exposure but require reliable identity infrastructure.}
\pitfall{Environment variables copied into sandboxes or crash logs leak broad authority.}
\subsection{Sandboxing}
Defense in depth combines process, filesystem, network, syscall, credential, resource, and time boundaries.
\tradeoff{Strong isolation raises operational cost and may restrict legitimate tools.}
\pitfall{Running as non-root inside an otherwise privileged container is insufficient.}
\subsection{Data privacy}
Data minimization, purpose limitation, access control, retention, deletion, and regional handling apply to prompts and traces.
\tradeoff{Privacy protection can reduce debugging data and personalization.}
\pitfall{Redacting only user messages while retaining retrieved documents leaves exposure.}
\subsection{Output validation}
Structured schema, semantic checks, policy filters, allowlists, and deterministic verification gate model output.
\tradeoff{Fail-closed validation improves safety but may lower completion rate.}
\pitfall{A valid JSON schema does not prove safe SQL, URLs, or business intent.}
\subsection{Model supply chain}
Weights, datasets, containers, packages, adapters, and prompts require provenance, scanning, signing, and review.
\tradeoff{Assurance improves while promotion pipelines become more deliberate.}
\pitfall{Loading arbitrary pickle-based artifacts can execute code.}
\subsection{AI risk management}
Risk processes map context, measure impacts, manage controls, document evidence, and govern residual risk.
\tradeoff{Governance improves accountability but can become checkbox compliance without technical signals.}
\pitfall{A model card without deployment-specific threat modeling is incomplete.}
\subsection{Red teaming}
Adversarial testing explores misuse, prompt injection, privacy, bias, jailbreaks, tool abuse, and failure recovery.
\tradeoff{It reveals unknown weaknesses but provides samples, not a guarantee of safety.}
\pitfall{Running one generic jailbreak list misses application-specific assets and actions.}
\subsection{Auditability}
Immutable event records connect identity, input, policy, model, evidence, tool calls, decisions, and outcomes.
\tradeoff{Investigations and compliance improve, while logs themselves become sensitive assets.}
\pitfall{Logging reasoning text instead of decisive actions and evidence can reduce clarity.}
\section{Distributed Systems and Reliability}
\subsection{CAP and consistency}
Under a network partition, a distributed system must trade availability against consistent reads and writes.
\tradeoff{Stronger guarantees simplify application invariants but increase latency and unavailability.}
\pitfall{Quoting CAP without identifying partitions and operations does not guide design.}
\subsection{Consensus}
Protocols such as Raft replicate an ordered log so nodes agree despite failures.
\tradeoff{Coordination provides a consistent control plane but costs extra messages and quorum latency.}
\pitfall{Deploying a quorum across unreliable links can reduce availability.}
\subsection{Queues and backpressure}
Bounded queues absorb bursts while admission control and backpressure prevent producers from overwhelming consumers.
\tradeoff{Smoothing improves utilization, while waiting increases latency and stale work.}
\pitfall{An unbounded queue converts overload into memory exhaustion.}
\subsection{Retries and timeouts}
Timeouts bound uncertainty; retries use budgets, backoff, jitter, and idempotency for transient failures.
\tradeoff{Success rate improves, while retries amplify overload and tail latency.}
\pitfall{Layered automatic retries create multiplicative request storms.}
\subsection{Caching}
Caches trade staleness and invalidation complexity for lower latency and load.
\tradeoff{Hit rates improve cost, but correctness depends on keys, TTLs, and invalidation.}
\pitfall{Omitting tenant, model, prompt version, or authorization from a cache key leaks data.}
\subsection{Load balancing}
Balancers route by health, load, locality, affinity, or capacity across replicas.
\tradeoff{Distribution improves availability but generic least-connections may ignore token workload size.}
\pitfall{Sending a long-context request to the least-busy replica can still cause head-of-line blocking.}
\subsection{Fault isolation}
Bulkheads, cells, quotas, and tenant boundaries prevent one failure or workload from consuming shared resources.
\tradeoff{Blast radius shrinks at some efficiency cost.}
\pitfall{A single global vector index or queue can defeat otherwise isolated services.}
\subsection{Exactly-once effects}
Systems typically combine at-least-once delivery with idempotent consumers, dedupe keys, and transactional outboxes.
\tradeoff{Business effects can be effectively once, but state and retention are needed.}
\pitfall{Equating broker delivery semantics with end-to-end exactly-once effects is incorrect.}
\subsection{Multi-region design}
Regions trade user latency, residency, failover speed, replication lag, and operational complexity.
\tradeoff{Active-active improves availability and latency but makes writes and conflict resolution harder.}
\pitfall{A standby region that never receives realistic traffic may fail during disaster recovery.}
\subsection{Capacity planning}
Workload models connect arrival rate, prompt and output tokens, service time, concurrency, memory, and SLO targets.
\tradeoff{Headroom handles bursts but idle accelerators are expensive.}
\pitfall{Planning with average sequence length underestimates tail memory and latency.}
\clearpage\section{Core 50 interview questions}
\subsection*{CORE-01 - Why should BM25 remain in a modern RAG baseline?}
\badge{MEDIUM} \quad Retrieval and Search Theory
\par\smallskip\textbf{Answer rubric.} BM25 is inexpensive, interpretable, and strong on rare terms, identifiers, and exact wording. Dense retrieval covers paraphrases but can miss those cases. Compare lexical, dense, and hybrid retrieval on the same judged set, then use slice metrics rather than assuming one method dominates.
\subsection*{CORE-02 - Derive the role of term-frequency saturation and document-length normalization in BM25.}
\badge{HARD} \quad Retrieval and Search Theory
\par\smallskip\textbf{Answer rubric.} BM25 increases score with term frequency but uses a saturating fraction so repeated terms have diminishing returns. Its length normalization compares document length with the collection average, preventing long documents from winning just because they contain more words. Explain how k1 controls saturation and b controls normalization, then discuss tuning by corpus.
\subsection*{CORE-03 - When would you choose pgvector over a dedicated vector database?}
\badge{MEDIUM} \quad Vector Databases and Search Engines
\par\smallskip\textbf{Answer rubric.} Choose pgvector when relational joins, transactions, operational simplicity, and moderate scale dominate. A dedicated service becomes attractive for independently scaled vector workloads, specialized distributed indexing, or managed operations. Validate both with filtered recall, latency, update rate, backup, and total cost.
\subsection*{CORE-04 - Compare HNSW with IVF-PQ for a billion-vector workload.}
\badge{HARD} \quad Vector Databases and Search Engines
\par\smallskip\textbf{Answer rubric.} HNSW usually provides strong latency-recall but consumes substantial RAM and has costly construction. IVF-PQ compresses vectors and searches selected clusters, reducing memory at the cost of training and quantization error. Benchmark representative filters and updates; do not make the choice from unfiltered synthetic vectors.
\subsection*{CORE-05 - How do you select a chunk size?}
\badge{MEDIUM} \quad RAG Architecture and Evaluation
\par\smallskip\textbf{Answer rubric.} Use document structure and question evidence span, not one universal token count. Measure retrieval recall, context precision, answer quality, and cost across headings, tables, code, and prose. Parent-child retrieval can decouple precise retrieval chunks from larger generation context.
\subsection*{CORE-06 - Design an evaluation suite that localizes RAG failures.}
\badge{HARD} \quad RAG Architecture and Evaluation
\par\smallskip\textbf{Answer rubric.} Separate ingestion correctness, retriever recall and ranking, context packing, faithfulness, answer correctness, citation entailment, latency, and cost. Build golden and synthetic cases with important slices, calibrate automated judges against humans, and retain per-stage traces so a final-answer regression can be attributed.
\subsection*{CORE-07 - What is the difference between faithfulness and factual correctness?}
\badge{HARD} \quad RAG Architecture and Evaluation
\par\smallskip\textbf{Answer rubric.} Faithfulness asks whether claims follow from supplied context; factual correctness asks whether they are true in the world or against a reference. An answer can be faithful to incorrect context, or factually correct but unsupported by retrieved evidence. Production evaluation often needs both plus citation quality.
\subsection*{CORE-08 - When is GraphRAG justified over vector RAG?}
\badge{MEDIUM} \quad Knowledge Graphs and GraphRAG
\par\smallskip\textbf{Answer rubric.} GraphRAG is justified for relationship-heavy, multi-hop, entity-centric, or whole-corpus thematic questions where graph structure adds retrieval value. It costs more to construct, resolve, update, and evaluate. Start with a question taxonomy and prove lift on graph-shaped slices.
\subsection*{CORE-09 - Compare GraphRAG local search, global search, and KAG.}
\badge{HARD} \quad Knowledge Graphs and GraphRAG
\par\smallskip\textbf{Answer rubric.} Local search expands around query entities for detailed multi-hop context. Global search aggregates community summaries for corpus-level themes. KAG adds schema-aware representation, chunk-graph mutual indexing, and logical-form-guided reasoning for professional domains; it demands stronger domain modeling.
\subsection*{CORE-10 - How do you prevent hallucinated graph edges from becoming authoritative?}
\badge{HARD} \quad Knowledge Graphs and GraphRAG
\par\smallskip\textbf{Answer rubric.} Attach every entity and edge to source spans, extraction model versions, confidence, and temporal validity. Apply schema constraints, deterministic validators, duplicate resolution, sampled human review, and downstream contradiction checks. Queries should be able to require verified provenance.
\subsection*{CORE-11 - What makes an agent different from a deterministic workflow?}
\badge{MEDIUM} \quad Agents, MCP, and Control
\par\smallskip\textbf{Answer rubric.} A workflow has predefined transitions; an agent chooses actions dynamically from observations and tools. Many production systems combine both: deterministic boundaries and approvals around agentic decision points. Use the least autonomy needed for the task.
\subsection*{CORE-12 - Design safe retries for an agent that can send payments.}
\badge{HARD} \quad Agents, MCP, and Control
\par\smallskip\textbf{Answer rubric.} Generate a durable business idempotency key before the side effect, validate authorization and amount at execution, persist the result, and make retries return that result. Separate planning from execution, require approval for policy-defined thresholds, and reconcile ambiguous timeouts with the payment provider before retrying.
\subsection*{CORE-13 - Explain MCP's trust boundary.}
\badge{MEDIUM} \quad Agents, MCP, and Control
\par\smallskip\textbf{Answer rubric.} MCP standardizes capability discovery and invocation; it does not make a server, tool result, or prompt trusted. The host must authenticate servers, scope credentials, authorize every action, surface consent, validate inputs and outputs, and isolate untrusted content from instructions.
\subsection*{CORE-14 - How would you bound an autonomous coding agent?}
\badge{HARD} \quad Agents, MCP, and Control
\par\smallskip\textbf{Answer rubric.} Use an ephemeral sandbox, scoped repository checkout, no ambient credentials, default-deny network, resource and time limits, command allow/deny policy, and reviewed patches. Separate read, write, test, and publish permissions; require explicit approval for destructive or external actions and retain an audit trail.
\subsection*{CORE-15 - What memory should an agent store, and what should it forget?}
\badge{HARD} \quad Agents, MCP, and Control
\par\smallskip\textbf{Answer rubric.} Store only durable, user-approved facts or proven procedural lessons with provenance, timestamps, scope, and deletion controls. Keep ephemeral task state separate. Avoid raw transcripts, secrets, uncertain inferences, and sensitive facts without purpose and retention policy.
\subsection*{CORE-16 - When should you use LangGraph rather than a simple function pipeline?}
\badge{MEDIUM} \quad Orchestration Frameworks
\par\smallskip\textbf{Answer rubric.} Use it when execution is long-running, stateful, branching, resumable, streamed, or approval-gated. A short deterministic request often needs ordinary code instead. Complexity should be justified by checkpointing, interrupts, or graph-level observability.
\subsection*{CORE-17 - Why must durable workflow nodes be idempotent?}
\badge{HARD} \quad Orchestration Frameworks
\par\smallskip\textbf{Answer rubric.} Checkpoint replay or activity retry may execute a node more than once after ambiguous failure. If side effects are not idempotent, a resume can duplicate them. Persist operation keys and outcomes, and isolate side effects in retry-aware activities.
\subsection*{CORE-18 - What should an LLM trace contain?}
\badge{MEDIUM} \quad Observability and Monitoring
\par\smallskip\textbf{Answer rubric.} Include request and session correlation, model and prompt versions, sanitized inputs and outputs, retrieval and rerank steps, tool calls, token usage, latency, cost, errors, evaluations, and policy decisions. Capture content only under explicit privacy and retention controls.
\subsection*{CORE-19 - How do you define SLOs for a streaming LLM application?}
\badge{HARD} \quad Observability and Monitoring
\par\smallskip\textbf{Answer rubric.} Separate availability, time to first token, inter-token latency, end-to-end completion, and semantic quality. Define percentiles by workload slice and model route, then pair error budgets with capacity and fallback policy. Average latency is not sufficient.
\subsection*{CORE-20 - How do you detect quality drift without immediate labels?}
\badge{HARD} \quad Observability and Monitoring
\par\smallskip\textbf{Answer rubric.} Track input and retrieval distributions, refusal and tool patterns, deterministic validators, sampled judge scores, citation coverage, user corrections, and business proxies. Calibrate alarms against delayed human labels and segment by cohort; drift signals are hypotheses, not proof of degradation.
\subsection*{CORE-21 - Explain continuous batching.}
\badge{MEDIUM} \quad Serving, Deployment, and LLMOps
\par\smallskip\textbf{Answer rubric.} A serving engine admits new requests and removes finished ones between decode iterations, keeping the GPU busy despite variable sequence lengths. It improves throughput but requires fairness, queue control, and latency isolation. Benchmark concurrency and sequence distributions, not only batch size.
\subsection*{CORE-22 - What metrics should drive LLM autoscaling?}
\badge{HARD} \quad Serving, Deployment, and LLMOps
\par\smallskip\textbf{Answer rubric.} Use queue depth or queued tokens, time to first token, decode throughput, active sequences, KV-cache utilization, memory headroom, and startup time. CPU often has weak correlation with GPU saturation. Predictive warm pools may be needed because model loading is slow.
\subsection*{CORE-23 - Design a safe multi-provider LLM gateway.}
\badge{HARD} \quad Serving, Deployment, and LLMOps
\par\smallskip\textbf{Answer rubric.} Normalize contracts but preserve provider-specific capability metadata. Add authentication, tenant budgets, token-aware limits, policy routing, retries with budgets, circuit breakers, audit logs, and semantic fallback rules. Pin model versions and expose route changes so fallbacks do not silently alter behavior.
\subsection*{CORE-24 - What does data versioning add beyond Git?}
\badge{MEDIUM} \quad Data and ML Lifecycle
\par\smallskip\textbf{Answer rubric.} Large artifacts live in object storage while Git records content-addressed pointers, pipeline definitions, and lineage. Reproducibility requires code, parameters, environment, and data versions together. A mutable remote path is not a version.
\subsection*{CORE-25 - Explain point-in-time correct feature joins.}
\badge{HARD} \quad Data and ML Lifecycle
\par\smallskip\textbf{Answer rubric.} For each training label timestamp, select only feature values available at or before that time. This prevents future information from leaking into historical examples. The design needs entity keys, event and availability time, late-data policy, and parity tests with online serving.
\subsection*{CORE-26 - Why is attention scaled by the square root of the head dimension?}
\badge{MEDIUM} \quad Transformer Theory and Training
\par\smallskip\textbf{Answer rubric.} Without scaling, query-key dot-product variance grows with dimension and pushes softmax into saturated regions with tiny gradients. Dividing by sqrt(d\_k) stabilizes logits under common initialization assumptions. It is a variance-control argument, not a learned temperature.
\subsection*{CORE-27 - Compare pre-norm and post-norm transformers.}
\badge{HARD} \quad Transformer Theory and Training
\par\smallskip\textbf{Answer rubric.} Pre-norm applies normalization before each sublayer and leaves a cleaner residual path, usually improving gradient flow in deep models. Post-norm normalizes after the residual addition and can need careful warmup or scaling. Architecture changes require retuning rather than a mechanical swap.
\subsection*{CORE-28 - Explain why MoE increases parameters without proportional FLOPs.}
\badge{HARD} \quad Transformer Theory and Training
\par\smallskip\textbf{Answer rubric.} A router activates only a small number of experts per token, so total expert parameters can grow while per-token compute follows the selected experts. The hidden costs are expert memory, routing balance, all-to-all communication, and capacity overflow.
\subsection*{CORE-29 - Compare SFT, PPO-based RLHF, and DPO.}
\badge{MEDIUM} \quad Transformer Theory and Training
\par\smallskip\textbf{Answer rubric.} SFT imitates demonstrations. PPO-based RLHF optimizes a learned reward with a policy constraint but has a complex training loop. DPO learns directly from preference pairs relative to a reference policy; it is simpler but still inherits preference-data bias and coverage limits.
\subsection*{CORE-30 - How do LoRA rank and target modules affect adaptation?}
\badge{HARD} \quad Transformer Theory and Training
\par\smallskip\textbf{Answer rubric.} Rank bounds the update subspace; higher rank increases capacity and memory. Targeting attention projections, MLPs, or both changes what behavior can move. Select with validation by task slice and inspect adapter interference rather than using a universal rank.
\subsection*{CORE-31 - What exactly does a KV cache save?}
\badge{MEDIUM} \quad Inference Optimization
\par\smallskip\textbf{Answer rubric.} It stores prior-layer key and value tensors so decoding a new token computes only its new projections and attends over cached history. It does not eliminate attention over prior positions. Memory grows with layers, sequence length, batch, KV heads, head dimension, and precision.
\subsection*{CORE-32 - Why can lower-bit quantization be slower?}
\badge{HARD} \quad Inference Optimization
\par\smallskip\textbf{Answer rubric.} A quantized model wins only if the hardware and runtime provide efficient dequantization and matrix kernels for its shapes. Unsupported formats add conversion overhead or fall back to slow paths. Measure end-to-end latency, memory, energy, and quality on target hardware.
\subsection*{CORE-33 - Explain speculative decoding's acceptance condition and benefit.}
\badge{HARD} \quad Inference Optimization
\par\smallskip\textbf{Answer rubric.} A draft model proposes several tokens; the target evaluates them in parallel and accepts a prefix using a correction rule that preserves the target distribution. Benefit depends on draft cost, acceptance rate, verification efficiency, and output length. Low acceptance can make it slower.
\subsection*{CORE-34 - FastAPI or Flask for an inference API?}
\badge{MEDIUM} \quad Backend, APIs, and Microservices
\par\smallskip\textbf{Answer rubric.} FastAPI's ASGI model, typing, validation, OpenAPI, and streaming fit concurrent APIs. Flask remains excellent for simple synchronous services and its ecosystem. The decisive issues are workload, deployment server, blocking model calls, team experience, and operational needs.
\subsection*{CORE-35 - Design backpressure for token streaming.}
\badge{HARD} \quad Backend, APIs, and Microservices
\par\smallskip\textbf{Answer rubric.} Bound admission and per-connection buffers, propagate slow-consumer signals, cancel model work on disconnect, and separate interactive from batch queues. Use token-weighted concurrency and deadlines. Never let unbounded output buffers consume server memory.
\subsection*{CORE-36 - Where should microservice boundaries sit in an AI platform?}
\badge{HARD} \quad Backend, APIs, and Microservices
\par\smallskip\textbf{Answer rubric.} Separate capabilities with distinct ownership, data, scaling, failure, and release characteristics, such as ingestion, retrieval, inference, and evaluation. Keep strongly coupled low-latency steps together until evidence supports a split. Avoid a distributed monolith with shared databases and synchronized releases.
\subsection*{CORE-37 - When do Python threads help despite the GIL?}
\badge{MEDIUM} \quad Python Engineering
\par\smallskip\textbf{Answer rubric.} They help when work spends time in blocking I/O or native code that releases the GIL. Pure Python CPU-bound loops generally need processes, native vectorized code, or another runtime. Measure contention and library behavior instead of applying a slogan.
\subsection*{CORE-38 - How do you prevent blocking an asyncio event loop?}
\badge{HARD} \quad Python Engineering
\par\smallskip\textbf{Answer rubric.} Use genuinely asynchronous clients for network I/O, move unavoidable blocking functions to a bounded thread pool, use processes for CPU-bound Python, and instrument event-loop lag. Bound concurrency so offloading does not merely move overload to an executor queue.
\subsection*{CORE-39 - Why use functools.wraps in a decorator?}
\badge{MEDIUM} \quad Python Engineering
\par\smallskip\textbf{Answer rubric.} wraps copies key metadata and establishes a \_\_wrapped\_\_ chain, preserving names, documentation, signatures through introspection, and framework behavior. Without it, debugging, dependency injection, route registration, and tests may see the wrapper instead of the original callable.
\subsection*{CORE-40 - Compare threads, processes, and coroutines for an embedding service.}
\badge{HARD} \quad Python Engineering
\par\smallskip\textbf{Answer rubric.} Coroutines efficiently multiplex many network calls; threads integrate blocking SDKs; processes parallelize CPU-heavy Python and isolate failures. Native GPU or BLAS work may already release the GIL and batch internally. Choose after profiling serialization, memory duplication, and batching.
\subsection*{CORE-41 - Why is prompt injection not solved by a stronger system prompt?}
\badge{MEDIUM} \quad Security, Safety, and Governance
\par\smallskip\textbf{Answer rubric.} The model still processes attacker-controlled instructions in the same inference context and cannot enforce authority like an operating system. Security must come from least-privilege tools, content provenance, authorization, validation, isolation, and confirmation for consequential actions.
\subsection*{CORE-42 - Threat-model a RAG system.}
\badge{HARD} \quad Security, Safety, and Governance
\par\smallskip\textbf{Answer rubric.} Assets include private documents, credentials, indexes, prompts, outputs, and tools. Threats include poisoned ingestion, cross-tenant retrieval, indirect prompt injection, membership inference, data exfiltration, insecure citations, and logging leaks. Map controls to each boundary and test bypasses.
\subsection*{CORE-43 - What does least privilege mean for agents?}
\badge{HARD} \quad Security, Safety, and Governance
\par\smallskip\textbf{Answer rubric.} Grant capabilities per action, resource, tenant, and time, not a broad session token. Separate planning credentials from execution credentials, require step-up approval for high-risk actions, and issue ephemeral scoped tokens after policy checks. Audit both decisions and side effects.
\subsection*{CORE-44 - Why can retries make an outage worse?}
\badge{MEDIUM} \quad Distributed Systems and Reliability
\par\smallskip\textbf{Answer rubric.} If a dependency is overloaded, retries multiply traffic and consume more queues, threads, and connections. Use deadlines, small retry budgets, exponential backoff with jitter, circuit breakers, and idempotency. Retry only errors likely to be transient.
\subsection*{CORE-45 - How would you design effective exactly-once document ingestion?}
\badge{HARD} \quad Distributed Systems and Reliability
\par\smallskip\textbf{Answer rubric.} Use at-least-once delivery with a content or source-version idempotency key, transactional state transitions, deterministic chunk IDs, and an outbox for downstream indexing events. Consumers upsert by version and reconciliation detects partial completion. Exactly-once is a business invariant, not a broker slogan.
\subsection*{CORE-46 - Model capacity for an LLM service.}
\badge{HARD} \quad Distributed Systems and Reliability
\par\smallskip\textbf{Answer rubric.} Estimate arrival rate by slice, prompt and output token distributions, time to first token, decode rate, active sequence memory, batching efficiency, and SLO headroom. Use Little's Law for concurrency intuition, then validate with load tests including bursts, long contexts, and failures.
\subsection*{CORE-47 - How do OpenTelemetry traces, metrics, and logs complement each other?}
\badge{MEDIUM} \quad Observability and Monitoring
\par\smallskip\textbf{Answer rubric.} Traces explain individual request paths, metrics quantify fleet trends and SLOs, and logs capture discrete detailed events. Correlation fields connect them. Do not put high-cardinality or sensitive prompt content into metric labels.
\subsection*{CORE-48 - How do you evaluate an LLM-as-a-judge?}
\badge{HARD} \quad RAG Architecture and Evaluation
\par\smallskip\textbf{Answer rubric.} Create human-labeled calibration sets with disagreements and important slices, measure agreement and ranking consistency, test position and verbosity bias, blind model identities, and version judge prompts and models. Use deterministic checks where possible and do not let one judge define truth.
\subsection*{CORE-49 - How do you evaluate an autonomous agent beyond final success?}
\badge{HARD} \quad Agents, MCP, and Control
\par\smallskip\textbf{Answer rubric.} Score task success, partial progress, trajectory efficiency, tool selection, argument correctness, recovery, policy compliance, side effects, latency, and cost. Replay deterministic environments, inject failures, inspect traces, and use human review for ambiguous outcomes.
\subsection*{CORE-50 - Why do metadata filters complicate ANN recall?}
\badge{HARD} \quad Vector Databases and Search Engines
\par\smallskip\textbf{Answer rubric.} A graph or clustered index was built for vector neighborhoods, not necessarily the filtered subset. Post-filtering can underfill results, while pre-filtering can fragment traversal. Measure filtered recall by selectivity and use index or partition strategies that support common filters.
\clearpage\section{Progressive system-design challenges}
\subsection{Multi-tenant enterprise RAG}
Design a citation-first assistant over 50 million private documents for 500 organizations.
\begin{enumerate}
\item \textbf{Initial design} Separate an access-aware ingestion plane from the query plane. Preserve tenant, ACL, source version, and stable chunk IDs; retrieve lexical and dense candidates under authorization, rerank, pack diverse evidence, generate constrained citations, and log a redacted trace.
\item \textbf{Challenge: 10x documents and strict deletion} Partition by tenant and locality, version embedding indexes, process change-data events idempotently, and maintain a deletion ledger that removes source, chunks, caches, and trace content. Use blue-green reindexing and reconciliation scans.
\item \textbf{Challenge: answer quality falls after an embedding upgrade} Shadow the new index on frozen and live evaluation sets, compare recall and downstream quality by query slice, and inspect nearest-neighbor shifts. Keep dual-read capability and roll back the alias if guardrails fail.
\item \textbf{Tradeoffs and choices} pgvector reduces services when relational ACL joins dominate; a dedicated store may scale vector traffic independently. Hybrid retrieval protects exact terms. Strong ACL filtering can reduce ANN recall, so tenant partitioning and filtered-recall tests are essential.
\end{enumerate}
\subsection{Approval-gated agent platform}
Design a platform where agents use email, calendar, CRM, and payment tools with different risk levels.
\begin{enumerate}
\item \textbf{Initial design} Maintain an explicit state machine. The planner proposes typed actions; a deterministic policy engine checks identity, scope, risk, and consent; high-risk actions pause durably; the executor receives a short-lived scoped token and writes an immutable result.
\item \textbf{Challenge: a tool times out after charging a card} Never blindly retry. Query by idempotency key, reconcile provider state, and return the stored outcome. Keep planning and execution separate so model retries cannot invent new operation keys.
\item \textbf{Challenge: malicious email contains tool instructions} Mark email as untrusted data, do not concatenate it into authority-bearing instructions, restrict available capabilities, validate recipients and amounts from trusted state, and require approval for consequential actions.
\item \textbf{Tradeoffs and choices} Dynamic agents handle novel requests but deterministic workflows are safer for financial steps. Per-action least privilege and durable approvals add latency but reduce blast radius. Store concise state, not unrestricted reasoning transcripts.
\end{enumerate}
\subsection{Global multi-provider LLM gateway}
Design a gateway serving 20 products across three model providers with budgets and regional policy.
\begin{enumerate}
\item \textbf{Initial design} Authenticate services and users, normalize a minimal request contract, retain capability metadata, enforce regional and data policies, estimate tokens before admission, route by model requirements, and stream through bounded buffers with full usage accounting.
\item \textbf{Challenge: provider A has elevated p99 latency} Use deadline-aware circuit breaking and route eligible traffic to pinned alternatives. Preserve non-equivalent capability constraints, surface the actual model route, and protect the fallback from a retry storm.
\item \textbf{Challenge: one tenant consumes half the GPU fleet} Apply tenant token buckets, weighted fair queues, concurrent-token limits, and reserved capacity. Attribute cache, input, output, and tool cost to the tenant and alert on anomalies.
\item \textbf{Tradeoffs and choices} A central gateway simplifies governance but is a critical dependency. Global caching lowers cost but creates privacy and invalidation risk. Provider abstraction should not erase semantic differences such as tool schemas or context limits.
\end{enumerate}
\subsection{Autonomous code agent sandbox}
Design a code agent that can modify repositories and run tests without endangering infrastructure.
\begin{enumerate}
\item \textbf{Initial design} Clone a pinned snapshot into an ephemeral unprivileged sandbox. Mount only the workspace writable, deny network by default, issue no ambient credentials, cap CPU, memory, processes, disk, and time, then return a patch, test evidence, and audit record.
\item \textbf{Challenge: tests need package downloads} Use an allowlisted dependency proxy with pinned hashes and egress logging, or prebuild trusted images. Do not grant general internet access or repository write credentials to the sandbox.
\item \textbf{Challenge: repository contains malicious test code} Treat repository code as untrusted. Harden the runtime boundary, block host sockets and metadata endpoints, use disposable workers, scan artifacts, and separate the merge identity from the execution identity.
\item \textbf{Tradeoffs and choices} MicroVMs offer stronger tenant isolation than ordinary containers but start more slowly. Warm pools reduce latency yet require secure reset. Automated merging improves speed but should be limited to low-risk paths with strong tests.
\end{enumerate}
\subsection{Continuous AI evaluation platform}
Design a platform that evaluates prompts, models, RAG pipelines, and agents before and after release.
\begin{enumerate}
\item \textbf{Initial design} Version datasets, prompts, models, tool environments, scorers, and code. Run deterministic tests first, then model judges and sampled human review. Store per-case traces and aggregate by task, risk, language, and tenant slice.
\item \textbf{Challenge: judge scores improve but users complain} Audit judge calibration, position and verbosity bias, and dataset representativeness. Add complaint-derived cases, blind candidates, and connect online business and correction signals without treating them as perfect labels.
\item \textbf{Challenge: agent evaluations are flaky} Use deterministic simulators, frozen tool snapshots, seeded stochastic settings where supported, retry classification, and trajectory-level assertions. Separate environment failures from policy and reasoning failures.
\item \textbf{Tradeoffs and choices} Large golden sets improve coverage but slow releases; tier tests by speed and risk. LLM judges scale but need calibration. A single aggregate score is convenient and dangerous, so gate critical slices independently.
\end{enumerate}
\subsection{Resilient document ingestion}
Design an ingestion pipeline for PDFs, HTML, email, tables, and code with near-real-time updates.
\begin{enumerate}
\item \textbf{Initial design} Use source-version events and an idempotent state machine. Store immutable raw content, parse to a canonical document model, validate and classify, create deterministic chunks, enrich, and atomically promote lexical and vector index versions.
\item \textbf{Challenge: a parser upgrade changes every chunk boundary} Version parser and chunker output, build a parallel index, compare retrieval and citation stability, and switch an alias only after acceptance. Keep old versions for rollback and delete them by retention policy.
\item \textbf{Challenge: events arrive out of order} Use source sequence or modification versions, reject stale transitions, and reconcile current source inventory against indexed state. Tombstones must dominate older create events.
\item \textbf{Tradeoffs and choices} Exactly-once broker delivery is unnecessary if effects are idempotent. Rich parsing improves quality but increases cost. Near-real-time indexing may use smaller incremental shards later compacted into efficient layouts.
\end{enumerate}
\subsection{Compliance GraphRAG}
Design a system that answers policy questions across regulations, contracts, and internal controls over time.
\begin{enumerate}
\item \textbf{Initial design} Model regulations, obligations, controls, owners, jurisdictions, and validity intervals. Attach every node and edge to source spans and extraction versions. Route precise entity questions to local traversal and thematic questions to community summaries.
\item \textbf{Challenge: two regulations conflict} Preserve both facts with jurisdiction, authority, and valid time rather than overwriting. Apply explicit precedence rules where approved and surface the conflict with citations for human interpretation.
\item \textbf{Challenge: auditors need as-of answers} Use bitemporal or validity-aware facts, immutable source versions, and query-time cutoffs. Store the exact graph, prompt, model, and evidence versions used for each answer.
\item \textbf{Tradeoffs and choices} Formal schema and temporal semantics add modeling cost but are justified by auditability. Vector-only RAG is simpler for prose lookup; graphs earn their cost on relationships, conflicts, and multi-hop obligations.
\end{enumerate}
\subsection{Real-time support copilot}
Design an assistant that listens to a live support conversation and suggests grounded responses under 700 ms.
\begin{enumerate}
\item \textbf{Initial design} Stream ASR partials, maintain a compact conversation state, detect stable intent boundaries, retrieve from a tenant-aware cache and hybrid index, and use a low-latency model with citations. Never auto-send; suggestions remain human controlled.
\item \textbf{Challenge: latency spikes during peak hours} Use token-aware admission, prewarm capacity, degrade to cached macros or smaller models, cap retrieval and output budgets, and monitor time to first suggestion separately from final completion.
\item \textbf{Challenge: ASR mishears an account number} Do not infer high-impact identifiers from uncertain audio. Display confidence, request confirmation, and resolve accounts through authenticated UI context rather than model text.
\item \textbf{Tradeoffs and choices} Frequent suggestions feel responsive but can distract and amplify partial-transcript errors. Smaller models reduce latency but may need stronger templates. Human control lowers automation but is appropriate for uncertain real-time signals.
\end{enumerate}
\subsection{Multimodal product search}
Design text, image, and attribute search for 200 million commerce products.
\begin{enumerate}
\item \textbf{Initial design} Create versioned text and image embeddings, retain structured attributes, retrieve modality-specific candidates, fuse ranks, filter availability and policy, then rerank with behavioral features. Evaluate by query type and inventory segment.
\item \textbf{Challenge: new catalog items have no interactions} Rely on content embeddings and attributes, use exploration, and separate cold-start metrics. Avoid training labels that make popular products the only relevant answers.
\item \textbf{Challenge: exact SKU searches regress} Route identifier-shaped queries to lexical or key-value lookup and combine with semantic search only as fallback. Keep BM25 or exact matching in the hybrid design.
\item \textbf{Tradeoffs and choices} A shared multimodal space simplifies cross-modal retrieval but specialized encoders can be stronger. Rich reranking improves relevance at latency cost. Behavior data improves conversion but introduces position and popularity bias.
\end{enumerate}
\subsection{Multi-region LLM inference}
Design self-hosted inference for users in North America, Europe, and Asia with residency constraints.
\begin{enumerate}
\item \textbf{Initial design} Build independent regional cells with pinned model artifacts, local telemetry, token-aware queues, and residency-safe storage. A control plane distributes signed releases and policy, while the data plane serves without cross-region prompt transfer.
\item \textbf{Challenge: Europe loses half its GPU capacity} Apply regional admission and degradation first. Fail over only requests whose policy permits leaving the region; otherwise use a smaller local model or queue within deadlines. Communicate degraded capability explicitly.
\item \textbf{Challenge: model release behaves differently on one GPU type} Treat hardware and runtime as part of the release identity. Run per-hardware conformance, quality, and performance tests, and canary independently by cell before global promotion.
\item \textbf{Tradeoffs and choices} Active-active cells improve latency and availability but duplicate expensive weights and caches. Cross-region fallback raises residency and consistency concerns. Centralized control should not be required for serving an already-approved model.
\end{enumerate}
\subsection{LLM-assisted recommendation system}
Design recommendations that combine collaborative signals, content embeddings, and natural-language explanations.
\begin{enumerate}
\item \textbf{Initial design} Generate candidates from collaborative, popularity, and vector channels; rank with point-in-time correct features; apply inventory, policy, and diversity constraints; generate explanations only from verified product attributes and recommendation factors.
\item \textbf{Challenge: explanations claim unsupported benefits} Use structured evidence templates or constrained generation, validate every claim against catalog fields, and omit explanations when evidence is insufficient. Evaluate explanation faithfulness separately from ranking quality.
\item \textbf{Challenge: filter bubbles increase} Add calibrated diversity and exploration, monitor exposure by cohort and supplier, and distinguish user satisfaction from short-term click maximization.
\item \textbf{Tradeoffs and choices} LLMs add interface quality but should not replace a measured ranker by default. Exploration improves long-term learning with short-term risk. Feature freshness improves relevance but increases online-store complexity.
\end{enumerate}
\subsection{Fraud investigation agent}
Design an agent that assembles evidence and recommends actions but cannot silently freeze accounts.
\begin{enumerate}
\item \textbf{Initial design} The agent retrieves immutable transaction evidence, constructs a case timeline and graph, cites every claim, and proposes actions. Deterministic policy calculates risk and an authenticated investigator approves account-impacting actions.
\item \textbf{Challenge: adversary places prompt injection in transaction notes} Treat notes as untrusted evidence, escape and label them, never expose general tools to the evidence-processing model, and derive actions only through typed policy-checked proposals.
\item \textbf{Challenge: false positives disproportionately affect one cohort} Monitor error and review outcomes by legally and ethically appropriate slices, audit features and thresholds, add escalation and appeal paths, and pause automated recommendations where evidence is weak.
\item \textbf{Tradeoffs and choices} Graph views help investigators follow relationships but extraction error must be visible. Human review costs time but protects high-impact decisions. Online learning from investigator outcomes needs controls for reviewer bias.
\end{enumerate}
\clearpage\section{Complete practice bank}
Questions are grouped by topic. Four-option questions include the correct answer and a compact explanation; flashcards and long answers provide model rubrics.
\subsection{Retrieval and Search Theory}
\paragraph{MCQ-0001 [medium]} Which statement best describes BM25?
\begin{enumerate}[label=\Alph*.]
\item Expansion adds synonyms, generated subqueries, or pseudo-relevance terms before retrieval.
\item A sparse weighting scheme that combines within-document frequency with inverse collection frequency.
\item A probabilistic lexical ranker that rewards term frequency while saturating repeated terms and normalizing document length.
\item A bi-encoder maps queries and passages into a shared embedding space for nearest-neighbor search.
\end{enumerate}
\noindent\textbf{Answer.} A probabilistic lexical ranker that rewards term frequency while saturating repeated terms and normalizing document length.
\par\textit{A probabilistic lexical ranker that rewards term frequency while saturating repeated terms and normalizing document length. Tradeoff: It is fast, interpretable, and exact for rare tokens but cannot directly match paraphrases. Watch for: Treating it as obsolete removes a strong baseline and hurts identifier-heavy queries.}
\paragraph{MCQ-0002 [hard]} What is the central engineering tradeoff of BM25?
\begin{enumerate}[label=\Alph*.]
\item It is transparent and cheap but lacks BM25's term-frequency saturation and tuned length normalization.
\item It is fast, interpretable, and exact for rare tokens but cannot directly match paraphrases.
\item Semantic recall improves, while exact names, numbers, and out-of-domain language may regress.
\item Recall can improve at the expense of latency and topic drift.
\end{enumerate}
\noindent\textbf{Answer.} It is fast, interpretable, and exact for rare tokens but cannot directly match paraphrases.
\par\textit{A probabilistic lexical ranker that rewards term frequency while saturating repeated terms and normalizing document length. Tradeoff: It is fast, interpretable, and exact for rare tokens but cannot directly match paraphrases. Watch for: Treating it as obsolete removes a strong baseline and hurts identifier-heavy queries.}
\paragraph{MCQ-0003 [hard]} Which failure mode should an interviewer expect you to identify for BM25?
\begin{enumerate}[label=\Alph*.]
\item Comparing raw cosine scores across changing corpora without rebuilding IDF causes drift.
\item Letting an LLM expand unconstrained queries can inject unsupported intent.
\item Evaluating only on in-domain examples hides lexical misses and embedding drift.
\item Treating it as obsolete removes a strong baseline and hurts identifier-heavy queries.
\end{enumerate}
\noindent\textbf{Answer.} Treating it as obsolete removes a strong baseline and hurts identifier-heavy queries.
\par\textit{A probabilistic lexical ranker that rewards term frequency while saturating repeated terms and normalizing document length. Tradeoff: It is fast, interpretable, and exact for rare tokens but cannot directly match paraphrases. Watch for: Treating it as obsolete removes a strong baseline and hurts identifier-heavy queries.}
\paragraph{MCQ-0004 [medium]} A design review is specifically concerned with this risk: Treating it as obsolete removes a strong baseline and hurts identifier-heavy queries. Which topic should be reviewed first?
\begin{enumerate}[label=\Alph*.]
\item Query expansion
\item Dense retrieval
\item TF-IDF
\item BM25
\end{enumerate}
\noindent\textbf{Answer.} BM25
\par\textit{A probabilistic lexical ranker that rewards term frequency while saturating repeated terms and normalizing document length. Tradeoff: It is fast, interpretable, and exact for rare tokens but cannot directly match paraphrases. Watch for: Treating it as obsolete removes a strong baseline and hurts identifier-heavy queries.}
\paragraph{MCQ-0005 [hard]} Which design note most accurately balances the benefits and costs of BM25?
\begin{enumerate}[label=\Alph*.]
\item Recall can improve at the expense of latency and topic drift.
\item Semantic recall improves, while exact names, numbers, and out-of-domain language may regress.
\item It is fast, interpretable, and exact for rare tokens but cannot directly match paraphrases.
\item It is transparent and cheap but lacks BM25's term-frequency saturation and tuned length normalization.
\end{enumerate}
\noindent\textbf{Answer.} It is fast, interpretable, and exact for rare tokens but cannot directly match paraphrases.
\par\textit{A probabilistic lexical ranker that rewards term frequency while saturating repeated terms and normalizing document length. Tradeoff: It is fast, interpretable, and exact for rare tokens but cannot directly match paraphrases. Watch for: Treating it as obsolete removes a strong baseline and hurts identifier-heavy queries.}
\paragraph{MCQ-0006 [medium]} Which explanation would be strongest in a system-design interview about BM25?
\begin{enumerate}[label=\Alph*.]
\item A probabilistic lexical ranker that rewards term frequency while saturating repeated terms and normalizing document length.
\item Expansion adds synonyms, generated subqueries, or pseudo-relevance terms before retrieval.
\item A bi-encoder maps queries and passages into a shared embedding space for nearest-neighbor search.
\item A sparse weighting scheme that combines within-document frequency with inverse collection frequency.
\end{enumerate}
\noindent\textbf{Answer.} A probabilistic lexical ranker that rewards term frequency while saturating repeated terms and normalizing document length.
\par\textit{A probabilistic lexical ranker that rewards term frequency while saturating repeated terms and normalizing document length. Tradeoff: It is fast, interpretable, and exact for rare tokens but cannot directly match paraphrases. Watch for: Treating it as obsolete removes a strong baseline and hurts identifier-heavy queries.}
\paragraph{MCQ-0007 [hard]} Which production warning is most directly associated with BM25?
\begin{enumerate}[label=\Alph*.]
\item Treating it as obsolete removes a strong baseline and hurts identifier-heavy queries.
\item Evaluating only on in-domain examples hides lexical misses and embedding drift.
\item Letting an LLM expand unconstrained queries can inject unsupported intent.
\item Comparing raw cosine scores across changing corpora without rebuilding IDF causes drift.
\end{enumerate}
\noindent\textbf{Answer.} Treating it as obsolete removes a strong baseline and hurts identifier-heavy queries.
\par\textit{A probabilistic lexical ranker that rewards term frequency while saturating repeated terms and normalizing document length. Tradeoff: It is fast, interpretable, and exact for rare tokens but cannot directly match paraphrases. Watch for: Treating it as obsolete removes a strong baseline and hurts identifier-heavy queries.}
\paragraph{FC-0008 [medium]} Define BM25 in interview-ready language.
\noindent\textbf{Answer.} A probabilistic lexical ranker that rewards term frequency while saturating repeated terms and normalizing document length.
\par\textit{A probabilistic lexical ranker that rewards term frequency while saturating repeated terms and normalizing document length.}
\paragraph{FC-0009 [hard]} State the main tradeoff and production pitfall for BM25.
\noindent\textbf{Answer.} Tradeoff: It is fast, interpretable, and exact for rare tokens but cannot directly match paraphrases. Pitfall: Treating it as obsolete removes a strong baseline and hurts identifier-heavy queries.
\par\textit{Tradeoff: It is fast, interpretable, and exact for rare tokens but cannot directly match paraphrases. Pitfall: Treating it as obsolete removes a strong baseline and hurts identifier-heavy queries.}
\paragraph{LA-0010 [hard]} Explain BM25 from first principles, then show how you would evaluate and operate it in production.
\noindent\textbf{Answer.} Start with the mechanism: A probabilistic lexical ranker that rewards term frequency while saturating repeated terms and normalizing document length. Make the tradeoff explicit: It is fast, interpretable, and exact for rare tokens but cannot directly match paraphrases. Define workload-specific offline metrics and a latency/cost budget, test important slices, instrument the critical path, and create a rollback criterion. The production review must explicitly guard against this failure: Treating it as obsolete removes a strong baseline and hurts identifier-heavy queries.
\par\textit{A strong answer connects mechanism, measurable tradeoffs, failure isolation, observability, and rollback rather than listing product features.}
\paragraph{MCQ-0011 [medium]} Which statement best describes TF-IDF?
\begin{enumerate}[label=\Alph*.]
\item RRF combines ranked lists using reciprocal rank positions rather than raw score scales.
\item A sparse weighting scheme that combines within-document frequency with inverse collection frequency.
\item A probabilistic lexical ranker that rewards term frequency while saturating repeated terms and normalizing document length.
\item Learned sparse models retain token-aligned inverted-index retrieval while expanding or reweighting terms.
\end{enumerate}
\noindent\textbf{Answer.} A sparse weighting scheme that combines within-document frequency with inverse collection frequency.
\par\textit{A sparse weighting scheme that combines within-document frequency with inverse collection frequency. Tradeoff: It is transparent and cheap but lacks BM25's term-frequency saturation and tuned length normalization. Watch for: Comparing raw cosine scores across changing corpora without rebuilding IDF causes drift.}
\paragraph{MCQ-0012 [hard]} What is the central engineering tradeoff of TF-IDF?
\begin{enumerate}[label=\Alph*.]
\item They bridge lexical and semantic matching but increase index size and model cost.
\item It is transparent and cheap but lacks BM25's term-frequency saturation and tuned length normalization.
\item It is robust without score calibration but ignores the magnitude and confidence of scores.
\item It is fast, interpretable, and exact for rare tokens but cannot directly match paraphrases.
\end{enumerate}
\noindent\textbf{Answer.} It is transparent and cheap but lacks BM25's term-frequency saturation and tuned length normalization.
\par\textit{A sparse weighting scheme that combines within-document frequency with inverse collection frequency. Tradeoff: It is transparent and cheap but lacks BM25's term-frequency saturation and tuned length normalization. Watch for: Comparing raw cosine scores across changing corpora without rebuilding IDF causes drift.}
\paragraph{MCQ-0013 [hard]} Which failure mode should an interviewer expect you to identify for TF-IDF?
\begin{enumerate}[label=\Alph*.]
\item Unbounded expansion creates latency and storage blowups.
\item Using an overly small rank constant can overreward noisy top results.
\item Treating it as obsolete removes a strong baseline and hurts identifier-heavy queries.
\item Comparing raw cosine scores across changing corpora without rebuilding IDF causes drift.
\end{enumerate}
\noindent\textbf{Answer.} Comparing raw cosine scores across changing corpora without rebuilding IDF causes drift.
\par\textit{A sparse weighting scheme that combines within-document frequency with inverse collection frequency. Tradeoff: It is transparent and cheap but lacks BM25's term-frequency saturation and tuned length normalization. Watch for: Comparing raw cosine scores across changing corpora without rebuilding IDF causes drift.}
\paragraph{MCQ-0014 [medium]} A design review is specifically concerned with this risk: Comparing raw cosine scores across changing corpora without rebuilding IDF causes drift. Which topic should be reviewed first?
\begin{enumerate}[label=\Alph*.]
\item TF-IDF
\item Reciprocal rank fusion
\item BM25
\item Sparse neural retrieval
\end{enumerate}
\noindent\textbf{Answer.} TF-IDF
\par\textit{A sparse weighting scheme that combines within-document frequency with inverse collection frequency. Tradeoff: It is transparent and cheap but lacks BM25's term-frequency saturation and tuned length normalization. Watch for: Comparing raw cosine scores across changing corpora without rebuilding IDF causes drift.}
\paragraph{MCQ-0015 [hard]} Which design note most accurately balances the benefits and costs of TF-IDF?
\begin{enumerate}[label=\Alph*.]
\item It is fast, interpretable, and exact for rare tokens but cannot directly match paraphrases.
\item It is robust without score calibration but ignores the magnitude and confidence of scores.
\item They bridge lexical and semantic matching but increase index size and model cost.
\item It is transparent and cheap but lacks BM25's term-frequency saturation and tuned length normalization.
\end{enumerate}
\noindent\textbf{Answer.} It is transparent and cheap but lacks BM25's term-frequency saturation and tuned length normalization.
\par\textit{A sparse weighting scheme that combines within-document frequency with inverse collection frequency. Tradeoff: It is transparent and cheap but lacks BM25's term-frequency saturation and tuned length normalization. Watch for: Comparing raw cosine scores across changing corpora without rebuilding IDF causes drift.}
\paragraph{MCQ-0016 [medium]} Which explanation would be strongest in a system-design interview about TF-IDF?
\begin{enumerate}[label=\Alph*.]
\item Learned sparse models retain token-aligned inverted-index retrieval while expanding or reweighting terms.
\item A probabilistic lexical ranker that rewards term frequency while saturating repeated terms and normalizing document length.
\item RRF combines ranked lists using reciprocal rank positions rather than raw score scales.
\item A sparse weighting scheme that combines within-document frequency with inverse collection frequency.
\end{enumerate}
\noindent\textbf{Answer.} A sparse weighting scheme that combines within-document frequency with inverse collection frequency.
\par\textit{A sparse weighting scheme that combines within-document frequency with inverse collection frequency. Tradeoff: It is transparent and cheap but lacks BM25's term-frequency saturation and tuned length normalization. Watch for: Comparing raw cosine scores across changing corpora without rebuilding IDF causes drift.}
\paragraph{MCQ-0017 [hard]} Which production warning is most directly associated with TF-IDF?
\begin{enumerate}[label=\Alph*.]
\item Comparing raw cosine scores across changing corpora without rebuilding IDF causes drift.
\item Using an overly small rank constant can overreward noisy top results.
\item Treating it as obsolete removes a strong baseline and hurts identifier-heavy queries.
\item Unbounded expansion creates latency and storage blowups.
\end{enumerate}
\noindent\textbf{Answer.} Comparing raw cosine scores across changing corpora without rebuilding IDF causes drift.
\par\textit{A sparse weighting scheme that combines within-document frequency with inverse collection frequency. Tradeoff: It is transparent and cheap but lacks BM25's term-frequency saturation and tuned length normalization. Watch for: Comparing raw cosine scores across changing corpora without rebuilding IDF causes drift.}
\paragraph{FC-0018 [medium]} Define TF-IDF in interview-ready language.
\noindent\textbf{Answer.} A sparse weighting scheme that combines within-document frequency with inverse collection frequency.
\par\textit{A sparse weighting scheme that combines within-document frequency with inverse collection frequency.}
\paragraph{FC-0019 [hard]} State the main tradeoff and production pitfall for TF-IDF.
\noindent\textbf{Answer.} Tradeoff: It is transparent and cheap but lacks BM25's term-frequency saturation and tuned length normalization. Pitfall: Comparing raw cosine scores across changing corpora without rebuilding IDF causes drift.
\par\textit{Tradeoff: It is transparent and cheap but lacks BM25's term-frequency saturation and tuned length normalization. Pitfall: Comparing raw cosine scores across changing corpora without rebuilding IDF causes drift.}
\paragraph{LA-0020 [hard]} Explain TF-IDF from first principles, then show how you would evaluate and operate it in production.
\noindent\textbf{Answer.} Start with the mechanism: A sparse weighting scheme that combines within-document frequency with inverse collection frequency. Make the tradeoff explicit: It is transparent and cheap but lacks BM25's term-frequency saturation and tuned length normalization. Define workload-specific offline metrics and a latency/cost budget, test important slices, instrument the critical path, and create a rollback criterion. The production review must explicitly guard against this failure: Comparing raw cosine scores across changing corpora without rebuilding IDF causes drift.
\par\textit{A strong answer connects mechanism, measurable tradeoffs, failure isolation, observability, and rollback rather than listing product features.}
\paragraph{MCQ-0021 [medium]} Which statement best describes Dense retrieval?
\begin{enumerate}[label=\Alph*.]
\item Similarity functions, normalization, anisotropy, dimensionality, and training objectives shape neighborhood quality.
\item A sparse weighting scheme that combines within-document frequency with inverse collection frequency.
\item ColBERT-style systems preserve token-level embeddings and aggregate fine-grained similarities at query time.
\item A bi-encoder maps queries and passages into a shared embedding space for nearest-neighbor search.
\end{enumerate}
\noindent\textbf{Answer.} A bi-encoder maps queries and passages into a shared embedding space for nearest-neighbor search.
\par\textit{A bi-encoder maps queries and passages into a shared embedding space for nearest-neighbor search. Tradeoff: Semantic recall improves, while exact names, numbers, and out-of-domain language may regress. Watch for: Evaluating only on in-domain examples hides lexical misses and embedding drift.}
\paragraph{MCQ-0022 [hard]} What is the central engineering tradeoff of Dense retrieval?
\begin{enumerate}[label=\Alph*.]
\item Cosine, dot product, and Euclidean distance can be equivalent only under specific normalization assumptions.
\item It is transparent and cheap but lacks BM25's term-frequency saturation and tuned length normalization.
\item Semantic recall improves, while exact names, numbers, and out-of-domain language may regress.
\item Accuracy often beats single-vector retrieval, with a larger index and heavier scoring path.
\end{enumerate}
\noindent\textbf{Answer.} Semantic recall improves, while exact names, numbers, and out-of-domain language may regress.
\par\textit{A bi-encoder maps queries and passages into a shared embedding space for nearest-neighbor search. Tradeoff: Semantic recall improves, while exact names, numbers, and out-of-domain language may regress. Watch for: Evaluating only on in-domain examples hides lexical misses and embedding drift.}
\paragraph{MCQ-0023 [hard]} Which failure mode should an interviewer expect you to identify for Dense retrieval?
\begin{enumerate}[label=\Alph*.]
\item Comparing raw cosine scores across changing corpora without rebuilding IDF causes drift.
\item Treating a late-interaction index like one-vector-per-document underestimates capacity needs.
\item Mixing an index metric with a differently trained embedding objective silently harms ranking.
\item Evaluating only on in-domain examples hides lexical misses and embedding drift.
\end{enumerate}
\noindent\textbf{Answer.} Evaluating only on in-domain examples hides lexical misses and embedding drift.
\par\textit{A bi-encoder maps queries and passages into a shared embedding space for nearest-neighbor search. Tradeoff: Semantic recall improves, while exact names, numbers, and out-of-domain language may regress. Watch for: Evaluating only on in-domain examples hides lexical misses and embedding drift.}
\paragraph{MCQ-0024 [medium]} A design review is specifically concerned with this risk: Evaluating only on in-domain examples hides lexical misses and embedding drift. Which topic should be reviewed first?
\begin{enumerate}[label=\Alph*.]
\item Dense retrieval
\item Embedding geometry
\item Late interaction
\item TF-IDF
\end{enumerate}
\noindent\textbf{Answer.} Dense retrieval
\par\textit{A bi-encoder maps queries and passages into a shared embedding space for nearest-neighbor search. Tradeoff: Semantic recall improves, while exact names, numbers, and out-of-domain language may regress. Watch for: Evaluating only on in-domain examples hides lexical misses and embedding drift.}
\paragraph{MCQ-0025 [hard]} Which design note most accurately balances the benefits and costs of Dense retrieval?
\begin{enumerate}[label=\Alph*.]
\item Semantic recall improves, while exact names, numbers, and out-of-domain language may regress.
\item Accuracy often beats single-vector retrieval, with a larger index and heavier scoring path.
\item Cosine, dot product, and Euclidean distance can be equivalent only under specific normalization assumptions.
\item It is transparent and cheap but lacks BM25's term-frequency saturation and tuned length normalization.
\end{enumerate}
\noindent\textbf{Answer.} Semantic recall improves, while exact names, numbers, and out-of-domain language may regress.
\par\textit{A bi-encoder maps queries and passages into a shared embedding space for nearest-neighbor search. Tradeoff: Semantic recall improves, while exact names, numbers, and out-of-domain language may regress. Watch for: Evaluating only on in-domain examples hides lexical misses and embedding drift.}
\paragraph{MCQ-0026 [medium]} Which explanation would be strongest in a system-design interview about Dense retrieval?
\begin{enumerate}[label=\Alph*.]
\item A bi-encoder maps queries and passages into a shared embedding space for nearest-neighbor search.
\item Similarity functions, normalization, anisotropy, dimensionality, and training objectives shape neighborhood quality.
\item ColBERT-style systems preserve token-level embeddings and aggregate fine-grained similarities at query time.
\item A sparse weighting scheme that combines within-document frequency with inverse collection frequency.
\end{enumerate}
\noindent\textbf{Answer.} A bi-encoder maps queries and passages into a shared embedding space for nearest-neighbor search.
\par\textit{A bi-encoder maps queries and passages into a shared embedding space for nearest-neighbor search. Tradeoff: Semantic recall improves, while exact names, numbers, and out-of-domain language may regress. Watch for: Evaluating only on in-domain examples hides lexical misses and embedding drift.}
\paragraph{MCQ-0027 [hard]} Which production warning is most directly associated with Dense retrieval?
\begin{enumerate}[label=\Alph*.]
\item Comparing raw cosine scores across changing corpora without rebuilding IDF causes drift.
\item Mixing an index metric with a differently trained embedding objective silently harms ranking.
\item Treating a late-interaction index like one-vector-per-document underestimates capacity needs.
\item Evaluating only on in-domain examples hides lexical misses and embedding drift.
\end{enumerate}
\noindent\textbf{Answer.} Evaluating only on in-domain examples hides lexical misses and embedding drift.
\par\textit{A bi-encoder maps queries and passages into a shared embedding space for nearest-neighbor search. Tradeoff: Semantic recall improves, while exact names, numbers, and out-of-domain language may regress. Watch for: Evaluating only on in-domain examples hides lexical misses and embedding drift.}
\paragraph{FC-0028 [medium]} Define Dense retrieval in interview-ready language.
\noindent\textbf{Answer.} A bi-encoder maps queries and passages into a shared embedding space for nearest-neighbor search.
\par\textit{A bi-encoder maps queries and passages into a shared embedding space for nearest-neighbor search.}
\paragraph{FC-0029 [hard]} State the main tradeoff and production pitfall for Dense retrieval.
\noindent\textbf{Answer.} Tradeoff: Semantic recall improves, while exact names, numbers, and out-of-domain language may regress. Pitfall: Evaluating only on in-domain examples hides lexical misses and embedding drift.
\par\textit{Tradeoff: Semantic recall improves, while exact names, numbers, and out-of-domain language may regress. Pitfall: Evaluating only on in-domain examples hides lexical misses and embedding drift.}
\paragraph{LA-0030 [hard]} Explain Dense retrieval from first principles, then show how you would evaluate and operate it in production.
\noindent\textbf{Answer.} Start with the mechanism: A bi-encoder maps queries and passages into a shared embedding space for nearest-neighbor search. Make the tradeoff explicit: Semantic recall improves, while exact names, numbers, and out-of-domain language may regress. Define workload-specific offline metrics and a latency/cost budget, test important slices, instrument the critical path, and create a rollback criterion. The production review must explicitly guard against this failure: Evaluating only on in-domain examples hides lexical misses and embedding drift.
\par\textit{A strong answer connects mechanism, measurable tradeoffs, failure isolation, observability, and rollback rather than listing product features.}
\paragraph{MCQ-0031 [medium]} Which statement best describes Sparse neural retrieval?
\begin{enumerate}[label=\Alph*.]
\item Learned sparse models retain token-aligned inverted-index retrieval while expanding or reweighting terms.
\item ColBERT-style systems preserve token-level embeddings and aggregate fine-grained similarities at query time.
\item Recall@k, precision@k, MRR, MAP, and nDCG measure different ranking behaviors.
\item A bi-encoder maps queries and passages into a shared embedding space for nearest-neighbor search.
\end{enumerate}
\noindent\textbf{Answer.} Learned sparse models retain token-aligned inverted-index retrieval while expanding or reweighting terms.
\par\textit{Learned sparse models retain token-aligned inverted-index retrieval while expanding or reweighting terms. Tradeoff: They bridge lexical and semantic matching but increase index size and model cost. Watch for: Unbounded expansion creates latency and storage blowups.}
\paragraph{MCQ-0032 [hard]} What is the central engineering tradeoff of Sparse neural retrieval?
\begin{enumerate}[label=\Alph*.]
\item They bridge lexical and semantic matching but increase index size and model cost.
\item Semantic recall improves, while exact names, numbers, and out-of-domain language may regress.
\item Accuracy often beats single-vector retrieval, with a larger index and heavier scoring path.
\item A metric must match the product objective and judgment granularity.
\end{enumerate}
\noindent\textbf{Answer.} They bridge lexical and semantic matching but increase index size and model cost.
\par\textit{Learned sparse models retain token-aligned inverted-index retrieval while expanding or reweighting terms. Tradeoff: They bridge lexical and semantic matching but increase index size and model cost. Watch for: Unbounded expansion creates latency and storage blowups.}
\paragraph{MCQ-0033 [hard]} Which failure mode should an interviewer expect you to identify for Sparse neural retrieval?
\begin{enumerate}[label=\Alph*.]
\item Optimizing MRR for a task that consumes many passages can worsen context coverage.
\item Treating a late-interaction index like one-vector-per-document underestimates capacity needs.
\item Evaluating only on in-domain examples hides lexical misses and embedding drift.
\item Unbounded expansion creates latency and storage blowups.
\end{enumerate}
\noindent\textbf{Answer.} Unbounded expansion creates latency and storage blowups.
\par\textit{Learned sparse models retain token-aligned inverted-index retrieval while expanding or reweighting terms. Tradeoff: They bridge lexical and semantic matching but increase index size and model cost. Watch for: Unbounded expansion creates latency and storage blowups.}
\paragraph{MCQ-0034 [medium]} A design review is specifically concerned with this risk: Unbounded expansion creates latency and storage blowups. Which topic should be reviewed first?
\begin{enumerate}[label=\Alph*.]
\item Sparse neural retrieval
\item Dense retrieval
\item Retrieval metrics
\item Late interaction
\end{enumerate}
\noindent\textbf{Answer.} Sparse neural retrieval
\par\textit{Learned sparse models retain token-aligned inverted-index retrieval while expanding or reweighting terms. Tradeoff: They bridge lexical and semantic matching but increase index size and model cost. Watch for: Unbounded expansion creates latency and storage blowups.}
\paragraph{MCQ-0035 [hard]} Which design note most accurately balances the benefits and costs of Sparse neural retrieval?
\begin{enumerate}[label=\Alph*.]
\item Accuracy often beats single-vector retrieval, with a larger index and heavier scoring path.
\item Semantic recall improves, while exact names, numbers, and out-of-domain language may regress.
\item A metric must match the product objective and judgment granularity.
\item They bridge lexical and semantic matching but increase index size and model cost.
\end{enumerate}
\noindent\textbf{Answer.} They bridge lexical and semantic matching but increase index size and model cost.
\par\textit{Learned sparse models retain token-aligned inverted-index retrieval while expanding or reweighting terms. Tradeoff: They bridge lexical and semantic matching but increase index size and model cost. Watch for: Unbounded expansion creates latency and storage blowups.}
\paragraph{MCQ-0036 [medium]} Which explanation would be strongest in a system-design interview about Sparse neural retrieval?
\begin{enumerate}[label=\Alph*.]
\item Learned sparse models retain token-aligned inverted-index retrieval while expanding or reweighting terms.
\item Recall@k, precision@k, MRR, MAP, and nDCG measure different ranking behaviors.
\item ColBERT-style systems preserve token-level embeddings and aggregate fine-grained similarities at query time.
\item A bi-encoder maps queries and passages into a shared embedding space for nearest-neighbor search.
\end{enumerate}
\noindent\textbf{Answer.} Learned sparse models retain token-aligned inverted-index retrieval while expanding or reweighting terms.
\par\textit{Learned sparse models retain token-aligned inverted-index retrieval while expanding or reweighting terms. Tradeoff: They bridge lexical and semantic matching but increase index size and model cost. Watch for: Unbounded expansion creates latency and storage blowups.}
\paragraph{MCQ-0037 [hard]} Which production warning is most directly associated with Sparse neural retrieval?
\begin{enumerate}[label=\Alph*.]
\item Unbounded expansion creates latency and storage blowups.
\item Evaluating only on in-domain examples hides lexical misses and embedding drift.
\item Treating a late-interaction index like one-vector-per-document underestimates capacity needs.
\item Optimizing MRR for a task that consumes many passages can worsen context coverage.
\end{enumerate}
\noindent\textbf{Answer.} Unbounded expansion creates latency and storage blowups.
\par\textit{Learned sparse models retain token-aligned inverted-index retrieval while expanding or reweighting terms. Tradeoff: They bridge lexical and semantic matching but increase index size and model cost. Watch for: Unbounded expansion creates latency and storage blowups.}
\paragraph{FC-0038 [medium]} Define Sparse neural retrieval in interview-ready language.
\noindent\textbf{Answer.} Learned sparse models retain token-aligned inverted-index retrieval while expanding or reweighting terms.
\par\textit{Learned sparse models retain token-aligned inverted-index retrieval while expanding or reweighting terms.}
\paragraph{FC-0039 [hard]} State the main tradeoff and production pitfall for Sparse neural retrieval.
\noindent\textbf{Answer.} Tradeoff: They bridge lexical and semantic matching but increase index size and model cost. Pitfall: Unbounded expansion creates latency and storage blowups.
\par\textit{Tradeoff: They bridge lexical and semantic matching but increase index size and model cost. Pitfall: Unbounded expansion creates latency and storage blowups.}
\paragraph{LA-0040 [hard]} Explain Sparse neural retrieval from first principles, then show how you would evaluate and operate it in production.
\noindent\textbf{Answer.} Start with the mechanism: Learned sparse models retain token-aligned inverted-index retrieval while expanding or reweighting terms. Make the tradeoff explicit: They bridge lexical and semantic matching but increase index size and model cost. Define workload-specific offline metrics and a latency/cost budget, test important slices, instrument the critical path, and create a rollback criterion. The production review must explicitly guard against this failure: Unbounded expansion creates latency and storage blowups.
\par\textit{A strong answer connects mechanism, measurable tradeoffs, failure isolation, observability, and rollback rather than listing product features.}
\paragraph{MCQ-0041 [medium]} Which statement best describes Late interaction?
\begin{enumerate}[label=\Alph*.]
\item ColBERT-style systems preserve token-level embeddings and aggregate fine-grained similarities at query time.
\item Expansion adds synonyms, generated subqueries, or pseudo-relevance terms before retrieval.
\item A model jointly encodes query and candidate text to assign a relevance score.
\item A bi-encoder maps queries and passages into a shared embedding space for nearest-neighbor search.
\end{enumerate}
\noindent\textbf{Answer.} ColBERT-style systems preserve token-level embeddings and aggregate fine-grained similarities at query time.
\par\textit{ColBERT-style systems preserve token-level embeddings and aggregate fine-grained similarities at query time. Tradeoff: Accuracy often beats single-vector retrieval, with a larger index and heavier scoring path. Watch for: Treating a late-interaction index like one-vector-per-document underestimates capacity needs.}
\paragraph{MCQ-0042 [hard]} What is the central engineering tradeoff of Late interaction?
\begin{enumerate}[label=\Alph*.]
\item Accuracy often beats single-vector retrieval, with a larger index and heavier scoring path.
\item Semantic recall improves, while exact names, numbers, and out-of-domain language may regress.
\item It gives strong precision on a small candidate set but is too expensive for full-corpus retrieval.
\item Recall can improve at the expense of latency and topic drift.
\end{enumerate}
\noindent\textbf{Answer.} Accuracy often beats single-vector retrieval, with a larger index and heavier scoring path.
\par\textit{ColBERT-style systems preserve token-level embeddings and aggregate fine-grained similarities at query time. Tradeoff: Accuracy often beats single-vector retrieval, with a larger index and heavier scoring path. Watch for: Treating a late-interaction index like one-vector-per-document underestimates capacity needs.}
\paragraph{MCQ-0043 [hard]} Which failure mode should an interviewer expect you to identify for Late interaction?
\begin{enumerate}[label=\Alph*.]
\item Evaluating only on in-domain examples hides lexical misses and embedding drift.
\item Letting an LLM expand unconstrained queries can inject unsupported intent.
\item Treating a late-interaction index like one-vector-per-document underestimates capacity needs.
\item Reranking too many long passages dominates p95 latency.
\end{enumerate}
\noindent\textbf{Answer.} Treating a late-interaction index like one-vector-per-document underestimates capacity needs.
\par\textit{ColBERT-style systems preserve token-level embeddings and aggregate fine-grained similarities at query time. Tradeoff: Accuracy often beats single-vector retrieval, with a larger index and heavier scoring path. Watch for: Treating a late-interaction index like one-vector-per-document underestimates capacity needs.}
\paragraph{MCQ-0044 [medium]} A design review is specifically concerned with this risk: Treating a late-interaction index like one-vector-per-document underestimates capacity needs. Which topic should be reviewed first?
\begin{enumerate}[label=\Alph*.]
\item Dense retrieval
\item Cross-encoder reranking
\item Late interaction
\item Query expansion
\end{enumerate}
\noindent\textbf{Answer.} Late interaction
\par\textit{ColBERT-style systems preserve token-level embeddings and aggregate fine-grained similarities at query time. Tradeoff: Accuracy often beats single-vector retrieval, with a larger index and heavier scoring path. Watch for: Treating a late-interaction index like one-vector-per-document underestimates capacity needs.}
\paragraph{MCQ-0045 [hard]} Which design note most accurately balances the benefits and costs of Late interaction?
\begin{enumerate}[label=\Alph*.]
\item Recall can improve at the expense of latency and topic drift.
\item It gives strong precision on a small candidate set but is too expensive for full-corpus retrieval.
\item Semantic recall improves, while exact names, numbers, and out-of-domain language may regress.
\item Accuracy often beats single-vector retrieval, with a larger index and heavier scoring path.
\end{enumerate}
\noindent\textbf{Answer.} Accuracy often beats single-vector retrieval, with a larger index and heavier scoring path.
\par\textit{ColBERT-style systems preserve token-level embeddings and aggregate fine-grained similarities at query time. Tradeoff: Accuracy often beats single-vector retrieval, with a larger index and heavier scoring path. Watch for: Treating a late-interaction index like one-vector-per-document underestimates capacity needs.}
\paragraph{MCQ-0046 [medium]} Which explanation would be strongest in a system-design interview about Late interaction?
\begin{enumerate}[label=\Alph*.]
\item Expansion adds synonyms, generated subqueries, or pseudo-relevance terms before retrieval.
\item ColBERT-style systems preserve token-level embeddings and aggregate fine-grained similarities at query time.
\item A model jointly encodes query and candidate text to assign a relevance score.
\item A bi-encoder maps queries and passages into a shared embedding space for nearest-neighbor search.
\end{enumerate}
\noindent\textbf{Answer.} ColBERT-style systems preserve token-level embeddings and aggregate fine-grained similarities at query time.
\par\textit{ColBERT-style systems preserve token-level embeddings and aggregate fine-grained similarities at query time. Tradeoff: Accuracy often beats single-vector retrieval, with a larger index and heavier scoring path. Watch for: Treating a late-interaction index like one-vector-per-document underestimates capacity needs.}
\paragraph{MCQ-0047 [hard]} Which production warning is most directly associated with Late interaction?
\begin{enumerate}[label=\Alph*.]
\item Treating a late-interaction index like one-vector-per-document underestimates capacity needs.
\item Reranking too many long passages dominates p95 latency.
\item Letting an LLM expand unconstrained queries can inject unsupported intent.
\item Evaluating only on in-domain examples hides lexical misses and embedding drift.
\end{enumerate}
\noindent\textbf{Answer.} Treating a late-interaction index like one-vector-per-document underestimates capacity needs.
\par\textit{ColBERT-style systems preserve token-level embeddings and aggregate fine-grained similarities at query time. Tradeoff: Accuracy often beats single-vector retrieval, with a larger index and heavier scoring path. Watch for: Treating a late-interaction index like one-vector-per-document underestimates capacity needs.}
\paragraph{FC-0048 [medium]} Define Late interaction in interview-ready language.
\noindent\textbf{Answer.} ColBERT-style systems preserve token-level embeddings and aggregate fine-grained similarities at query time.
\par\textit{ColBERT-style systems preserve token-level embeddings and aggregate fine-grained similarities at query time.}
\paragraph{FC-0049 [hard]} State the main tradeoff and production pitfall for Late interaction.
\noindent\textbf{Answer.} Tradeoff: Accuracy often beats single-vector retrieval, with a larger index and heavier scoring path. Pitfall: Treating a late-interaction index like one-vector-per-document underestimates capacity needs.
\par\textit{Tradeoff: Accuracy often beats single-vector retrieval, with a larger index and heavier scoring path. Pitfall: Treating a late-interaction index like one-vector-per-document underestimates capacity needs.}
\paragraph{LA-0050 [hard]} Explain Late interaction from first principles, then show how you would evaluate and operate it in production.
\noindent\textbf{Answer.} Start with the mechanism: ColBERT-style systems preserve token-level embeddings and aggregate fine-grained similarities at query time. Make the tradeoff explicit: Accuracy often beats single-vector retrieval, with a larger index and heavier scoring path. Define workload-specific offline metrics and a latency/cost budget, test important slices, instrument the critical path, and create a rollback criterion. The production review must explicitly guard against this failure: Treating a late-interaction index like one-vector-per-document underestimates capacity needs.
\par\textit{A strong answer connects mechanism, measurable tradeoffs, failure isolation, observability, and rollback rather than listing product features.}
\paragraph{MCQ-0051 [medium]} Which statement best describes Cross-encoder reranking?
\begin{enumerate}[label=\Alph*.]
\item Recall@k, precision@k, MRR, MAP, and nDCG measure different ranking behaviors.
\item Expansion adds synonyms, generated subqueries, or pseudo-relevance terms before retrieval.
\item ColBERT-style systems preserve token-level embeddings and aggregate fine-grained similarities at query time.
\item A model jointly encodes query and candidate text to assign a relevance score.
\end{enumerate}
\noindent\textbf{Answer.} A model jointly encodes query and candidate text to assign a relevance score.
\par\textit{A model jointly encodes query and candidate text to assign a relevance score. Tradeoff: It gives strong precision on a small candidate set but is too expensive for full-corpus retrieval. Watch for: Reranking too many long passages dominates p95 latency.}
\paragraph{MCQ-0052 [hard]} What is the central engineering tradeoff of Cross-encoder reranking?
\begin{enumerate}[label=\Alph*.]
\item It gives strong precision on a small candidate set but is too expensive for full-corpus retrieval.
\item A metric must match the product objective and judgment granularity.
\item Recall can improve at the expense of latency and topic drift.
\item Accuracy often beats single-vector retrieval, with a larger index and heavier scoring path.
\end{enumerate}
\noindent\textbf{Answer.} It gives strong precision on a small candidate set but is too expensive for full-corpus retrieval.
\par\textit{A model jointly encodes query and candidate text to assign a relevance score. Tradeoff: It gives strong precision on a small candidate set but is too expensive for full-corpus retrieval. Watch for: Reranking too many long passages dominates p95 latency.}
\paragraph{MCQ-0053 [hard]} Which failure mode should an interviewer expect you to identify for Cross-encoder reranking?
\begin{enumerate}[label=\Alph*.]
\item Letting an LLM expand unconstrained queries can inject unsupported intent.
\item Treating a late-interaction index like one-vector-per-document underestimates capacity needs.
\item Optimizing MRR for a task that consumes many passages can worsen context coverage.
\item Reranking too many long passages dominates p95 latency.
\end{enumerate}
\noindent\textbf{Answer.} Reranking too many long passages dominates p95 latency.
\par\textit{A model jointly encodes query and candidate text to assign a relevance score. Tradeoff: It gives strong precision on a small candidate set but is too expensive for full-corpus retrieval. Watch for: Reranking too many long passages dominates p95 latency.}
\paragraph{MCQ-0054 [medium]} A design review is specifically concerned with this risk: Reranking too many long passages dominates p95 latency. Which topic should be reviewed first?
\begin{enumerate}[label=\Alph*.]
\item Late interaction
\item Cross-encoder reranking
\item Query expansion
\item Retrieval metrics
\end{enumerate}
\noindent\textbf{Answer.} Cross-encoder reranking
\par\textit{A model jointly encodes query and candidate text to assign a relevance score. Tradeoff: It gives strong precision on a small candidate set but is too expensive for full-corpus retrieval. Watch for: Reranking too many long passages dominates p95 latency.}
\paragraph{MCQ-0055 [hard]} Which design note most accurately balances the benefits and costs of Cross-encoder reranking?
\begin{enumerate}[label=\Alph*.]
\item Recall can improve at the expense of latency and topic drift.
\item It gives strong precision on a small candidate set but is too expensive for full-corpus retrieval.
\item A metric must match the product objective and judgment granularity.
\item Accuracy often beats single-vector retrieval, with a larger index and heavier scoring path.
\end{enumerate}
\noindent\textbf{Answer.} It gives strong precision on a small candidate set but is too expensive for full-corpus retrieval.
\par\textit{A model jointly encodes query and candidate text to assign a relevance score. Tradeoff: It gives strong precision on a small candidate set but is too expensive for full-corpus retrieval. Watch for: Reranking too many long passages dominates p95 latency.}
\paragraph{MCQ-0056 [medium]} Which explanation would be strongest in a system-design interview about Cross-encoder reranking?
\begin{enumerate}[label=\Alph*.]
\item Recall@k, precision@k, MRR, MAP, and nDCG measure different ranking behaviors.
\item Expansion adds synonyms, generated subqueries, or pseudo-relevance terms before retrieval.
\item ColBERT-style systems preserve token-level embeddings and aggregate fine-grained similarities at query time.
\item A model jointly encodes query and candidate text to assign a relevance score.
\end{enumerate}
\noindent\textbf{Answer.} A model jointly encodes query and candidate text to assign a relevance score.
\par\textit{A model jointly encodes query and candidate text to assign a relevance score. Tradeoff: It gives strong precision on a small candidate set but is too expensive for full-corpus retrieval. Watch for: Reranking too many long passages dominates p95 latency.}
\paragraph{MCQ-0057 [hard]} Which production warning is most directly associated with Cross-encoder reranking?
\begin{enumerate}[label=\Alph*.]
\item Treating a late-interaction index like one-vector-per-document underestimates capacity needs.
\item Optimizing MRR for a task that consumes many passages can worsen context coverage.
\item Reranking too many long passages dominates p95 latency.
\item Letting an LLM expand unconstrained queries can inject unsupported intent.
\end{enumerate}
\noindent\textbf{Answer.} Reranking too many long passages dominates p95 latency.
\par\textit{A model jointly encodes query and candidate text to assign a relevance score. Tradeoff: It gives strong precision on a small candidate set but is too expensive for full-corpus retrieval. Watch for: Reranking too many long passages dominates p95 latency.}
\paragraph{FC-0058 [medium]} Define Cross-encoder reranking in interview-ready language.
\noindent\textbf{Answer.} A model jointly encodes query and candidate text to assign a relevance score.
\par\textit{A model jointly encodes query and candidate text to assign a relevance score.}
\paragraph{FC-0059 [hard]} State the main tradeoff and production pitfall for Cross-encoder reranking.
\noindent\textbf{Answer.} Tradeoff: It gives strong precision on a small candidate set but is too expensive for full-corpus retrieval. Pitfall: Reranking too many long passages dominates p95 latency.
\par\textit{Tradeoff: It gives strong precision on a small candidate set but is too expensive for full-corpus retrieval. Pitfall: Reranking too many long passages dominates p95 latency.}
\paragraph{LA-0060 [hard]} Explain Cross-encoder reranking from first principles, then show how you would evaluate and operate it in production.
\noindent\textbf{Answer.} Start with the mechanism: A model jointly encodes query and candidate text to assign a relevance score. Make the tradeoff explicit: It gives strong precision on a small candidate set but is too expensive for full-corpus retrieval. Define workload-specific offline metrics and a latency/cost budget, test important slices, instrument the critical path, and create a rollback criterion. The production review must explicitly guard against this failure: Reranking too many long passages dominates p95 latency.
\par\textit{A strong answer connects mechanism, measurable tradeoffs, failure isolation, observability, and rollback rather than listing product features.}
\paragraph{MCQ-0061 [medium]} Which statement best describes Hybrid retrieval?
\begin{enumerate}[label=\Alph*.]
\item Similarity functions, normalization, anisotropy, dimensionality, and training objectives shape neighborhood quality.
\item A bi-encoder maps queries and passages into a shared embedding space for nearest-neighbor search.
\item Recall@k, precision@k, MRR, MAP, and nDCG measure different ranking behaviors.
\item Dense and lexical signals are retrieved together and fused, often with weighted scores or reciprocal-rank fusion.
\end{enumerate}
\noindent\textbf{Answer.} Dense and lexical signals are retrieved together and fused, often with weighted scores or reciprocal-rank fusion.
\par\textit{Dense and lexical signals are retrieved together and fused, often with weighted scores or reciprocal-rank fusion. Tradeoff: It improves robustness across semantic and exact-match queries but requires score calibration and duplicate handling. Watch for: Adding incomparable raw scores without normalization makes one retriever dominate.}
\paragraph{MCQ-0062 [hard]} What is the central engineering tradeoff of Hybrid retrieval?
\begin{enumerate}[label=\Alph*.]
\item Semantic recall improves, while exact names, numbers, and out-of-domain language may regress.
\item A metric must match the product objective and judgment granularity.
\item It improves robustness across semantic and exact-match queries but requires score calibration and duplicate handling.
\item Cosine, dot product, and Euclidean distance can be equivalent only under specific normalization assumptions.
\end{enumerate}
\noindent\textbf{Answer.} It improves robustness across semantic and exact-match queries but requires score calibration and duplicate handling.
\par\textit{Dense and lexical signals are retrieved together and fused, often with weighted scores or reciprocal-rank fusion. Tradeoff: It improves robustness across semantic and exact-match queries but requires score calibration and duplicate handling. Watch for: Adding incomparable raw scores without normalization makes one retriever dominate.}
\paragraph{MCQ-0063 [hard]} Which failure mode should an interviewer expect you to identify for Hybrid retrieval?
\begin{enumerate}[label=\Alph*.]
\item Mixing an index metric with a differently trained embedding objective silently harms ranking.
\item Adding incomparable raw scores without normalization makes one retriever dominate.
\item Evaluating only on in-domain examples hides lexical misses and embedding drift.
\item Optimizing MRR for a task that consumes many passages can worsen context coverage.
\end{enumerate}
\noindent\textbf{Answer.} Adding incomparable raw scores without normalization makes one retriever dominate.
\par\textit{Dense and lexical signals are retrieved together and fused, often with weighted scores or reciprocal-rank fusion. Tradeoff: It improves robustness across semantic and exact-match queries but requires score calibration and duplicate handling. Watch for: Adding incomparable raw scores without normalization makes one retriever dominate.}
\paragraph{MCQ-0064 [medium]} A design review is specifically concerned with this risk: Adding incomparable raw scores without normalization makes one retriever dominate. Which topic should be reviewed first?
\begin{enumerate}[label=\Alph*.]
\item Dense retrieval
\item Retrieval metrics
\item Embedding geometry
\item Hybrid retrieval
\end{enumerate}
\noindent\textbf{Answer.} Hybrid retrieval
\par\textit{Dense and lexical signals are retrieved together and fused, often with weighted scores or reciprocal-rank fusion. Tradeoff: It improves robustness across semantic and exact-match queries but requires score calibration and duplicate handling. Watch for: Adding incomparable raw scores without normalization makes one retriever dominate.}
\paragraph{MCQ-0065 [hard]} Which design note most accurately balances the benefits and costs of Hybrid retrieval?
\begin{enumerate}[label=\Alph*.]
\item Semantic recall improves, while exact names, numbers, and out-of-domain language may regress.
\item Cosine, dot product, and Euclidean distance can be equivalent only under specific normalization assumptions.
\item A metric must match the product objective and judgment granularity.
\item It improves robustness across semantic and exact-match queries but requires score calibration and duplicate handling.
\end{enumerate}
\noindent\textbf{Answer.} It improves robustness across semantic and exact-match queries but requires score calibration and duplicate handling.
\par\textit{Dense and lexical signals are retrieved together and fused, often with weighted scores or reciprocal-rank fusion. Tradeoff: It improves robustness across semantic and exact-match queries but requires score calibration and duplicate handling. Watch for: Adding incomparable raw scores without normalization makes one retriever dominate.}
\paragraph{MCQ-0066 [medium]} Which explanation would be strongest in a system-design interview about Hybrid retrieval?
\begin{enumerate}[label=\Alph*.]
\item Similarity functions, normalization, anisotropy, dimensionality, and training objectives shape neighborhood quality.
\item A bi-encoder maps queries and passages into a shared embedding space for nearest-neighbor search.
\item Recall@k, precision@k, MRR, MAP, and nDCG measure different ranking behaviors.
\item Dense and lexical signals are retrieved together and fused, often with weighted scores or reciprocal-rank fusion.
\end{enumerate}
\noindent\textbf{Answer.} Dense and lexical signals are retrieved together and fused, often with weighted scores or reciprocal-rank fusion.
\par\textit{Dense and lexical signals are retrieved together and fused, often with weighted scores or reciprocal-rank fusion. Tradeoff: It improves robustness across semantic and exact-match queries but requires score calibration and duplicate handling. Watch for: Adding incomparable raw scores without normalization makes one retriever dominate.}
\paragraph{MCQ-0067 [hard]} Which production warning is most directly associated with Hybrid retrieval?
\begin{enumerate}[label=\Alph*.]
\item Adding incomparable raw scores without normalization makes one retriever dominate.
\item Evaluating only on in-domain examples hides lexical misses and embedding drift.
\item Optimizing MRR for a task that consumes many passages can worsen context coverage.
\item Mixing an index metric with a differently trained embedding objective silently harms ranking.
\end{enumerate}
\noindent\textbf{Answer.} Adding incomparable raw scores without normalization makes one retriever dominate.
\par\textit{Dense and lexical signals are retrieved together and fused, often with weighted scores or reciprocal-rank fusion. Tradeoff: It improves robustness across semantic and exact-match queries but requires score calibration and duplicate handling. Watch for: Adding incomparable raw scores without normalization makes one retriever dominate.}
\paragraph{FC-0068 [medium]} Define Hybrid retrieval in interview-ready language.
\noindent\textbf{Answer.} Dense and lexical signals are retrieved together and fused, often with weighted scores or reciprocal-rank fusion.
\par\textit{Dense and lexical signals are retrieved together and fused, often with weighted scores or reciprocal-rank fusion.}
\paragraph{FC-0069 [hard]} State the main tradeoff and production pitfall for Hybrid retrieval.
\noindent\textbf{Answer.} Tradeoff: It improves robustness across semantic and exact-match queries but requires score calibration and duplicate handling. Pitfall: Adding incomparable raw scores without normalization makes one retriever dominate.
\par\textit{Tradeoff: It improves robustness across semantic and exact-match queries but requires score calibration and duplicate handling. Pitfall: Adding incomparable raw scores without normalization makes one retriever dominate.}
\paragraph{LA-0070 [hard]} Explain Hybrid retrieval from first principles, then show how you would evaluate and operate it in production.
\noindent\textbf{Answer.} Start with the mechanism: Dense and lexical signals are retrieved together and fused, often with weighted scores or reciprocal-rank fusion. Make the tradeoff explicit: It improves robustness across semantic and exact-match queries but requires score calibration and duplicate handling. Define workload-specific offline metrics and a latency/cost budget, test important slices, instrument the critical path, and create a rollback criterion. The production review must explicitly guard against this failure: Adding incomparable raw scores without normalization makes one retriever dominate.
\par\textit{A strong answer connects mechanism, measurable tradeoffs, failure isolation, observability, and rollback rather than listing product features.}
\paragraph{MCQ-0071 [medium]} Which statement best describes Reciprocal rank fusion?
\begin{enumerate}[label=\Alph*.]
\item Recall@k, precision@k, MRR, MAP, and nDCG measure different ranking behaviors.
\item A bi-encoder maps queries and passages into a shared embedding space for nearest-neighbor search.
\item RRF combines ranked lists using reciprocal rank positions rather than raw score scales.
\item ColBERT-style systems preserve token-level embeddings and aggregate fine-grained similarities at query time.
\end{enumerate}
\noindent\textbf{Answer.} RRF combines ranked lists using reciprocal rank positions rather than raw score scales.
\par\textit{RRF combines ranked lists using reciprocal rank positions rather than raw score scales. Tradeoff: It is robust without score calibration but ignores the magnitude and confidence of scores. Watch for: Using an overly small rank constant can overreward noisy top results.}
\paragraph{MCQ-0072 [hard]} What is the central engineering tradeoff of Reciprocal rank fusion?
\begin{enumerate}[label=\Alph*.]
\item It is robust without score calibration but ignores the magnitude and confidence of scores.
\item Accuracy often beats single-vector retrieval, with a larger index and heavier scoring path.
\item A metric must match the product objective and judgment granularity.
\item Semantic recall improves, while exact names, numbers, and out-of-domain language may regress.
\end{enumerate}
\noindent\textbf{Answer.} It is robust without score calibration but ignores the magnitude and confidence of scores.
\par\textit{RRF combines ranked lists using reciprocal rank positions rather than raw score scales. Tradeoff: It is robust without score calibration but ignores the magnitude and confidence of scores. Watch for: Using an overly small rank constant can overreward noisy top results.}
\paragraph{MCQ-0073 [hard]} Which failure mode should an interviewer expect you to identify for Reciprocal rank fusion?
\begin{enumerate}[label=\Alph*.]
\item Using an overly small rank constant can overreward noisy top results.
\item Optimizing MRR for a task that consumes many passages can worsen context coverage.
\item Evaluating only on in-domain examples hides lexical misses and embedding drift.
\item Treating a late-interaction index like one-vector-per-document underestimates capacity needs.
\end{enumerate}
\noindent\textbf{Answer.} Using an overly small rank constant can overreward noisy top results.
\par\textit{RRF combines ranked lists using reciprocal rank positions rather than raw score scales. Tradeoff: It is robust without score calibration but ignores the magnitude and confidence of scores. Watch for: Using an overly small rank constant can overreward noisy top results.}
\paragraph{MCQ-0074 [medium]} A design review is specifically concerned with this risk: Using an overly small rank constant can overreward noisy top results. Which topic should be reviewed first?
\begin{enumerate}[label=\Alph*.]
\item Retrieval metrics
\item Dense retrieval
\item Late interaction
\item Reciprocal rank fusion
\end{enumerate}
\noindent\textbf{Answer.} Reciprocal rank fusion
\par\textit{RRF combines ranked lists using reciprocal rank positions rather than raw score scales. Tradeoff: It is robust without score calibration but ignores the magnitude and confidence of scores. Watch for: Using an overly small rank constant can overreward noisy top results.}
\paragraph{MCQ-0075 [hard]} Which design note most accurately balances the benefits and costs of Reciprocal rank fusion?
\begin{enumerate}[label=\Alph*.]
\item Semantic recall improves, while exact names, numbers, and out-of-domain language may regress.
\item A metric must match the product objective and judgment granularity.
\item Accuracy often beats single-vector retrieval, with a larger index and heavier scoring path.
\item It is robust without score calibration but ignores the magnitude and confidence of scores.
\end{enumerate}
\noindent\textbf{Answer.} It is robust without score calibration but ignores the magnitude and confidence of scores.
\par\textit{RRF combines ranked lists using reciprocal rank positions rather than raw score scales. Tradeoff: It is robust without score calibration but ignores the magnitude and confidence of scores. Watch for: Using an overly small rank constant can overreward noisy top results.}
\paragraph{MCQ-0076 [medium]} Which explanation would be strongest in a system-design interview about Reciprocal rank fusion?
\begin{enumerate}[label=\Alph*.]
\item Recall@k, precision@k, MRR, MAP, and nDCG measure different ranking behaviors.
\item RRF combines ranked lists using reciprocal rank positions rather than raw score scales.
\item A bi-encoder maps queries and passages into a shared embedding space for nearest-neighbor search.
\item ColBERT-style systems preserve token-level embeddings and aggregate fine-grained similarities at query time.
\end{enumerate}
\noindent\textbf{Answer.} RRF combines ranked lists using reciprocal rank positions rather than raw score scales.
\par\textit{RRF combines ranked lists using reciprocal rank positions rather than raw score scales. Tradeoff: It is robust without score calibration but ignores the magnitude and confidence of scores. Watch for: Using an overly small rank constant can overreward noisy top results.}
\paragraph{MCQ-0077 [hard]} Which production warning is most directly associated with Reciprocal rank fusion?
\begin{enumerate}[label=\Alph*.]
\item Treating a late-interaction index like one-vector-per-document underestimates capacity needs.
\item Evaluating only on in-domain examples hides lexical misses and embedding drift.
\item Optimizing MRR for a task that consumes many passages can worsen context coverage.
\item Using an overly small rank constant can overreward noisy top results.
\end{enumerate}
\noindent\textbf{Answer.} Using an overly small rank constant can overreward noisy top results.
\par\textit{RRF combines ranked lists using reciprocal rank positions rather than raw score scales. Tradeoff: It is robust without score calibration but ignores the magnitude and confidence of scores. Watch for: Using an overly small rank constant can overreward noisy top results.}
\paragraph{FC-0078 [medium]} Define Reciprocal rank fusion in interview-ready language.
\noindent\textbf{Answer.} RRF combines ranked lists using reciprocal rank positions rather than raw score scales.
\par\textit{RRF combines ranked lists using reciprocal rank positions rather than raw score scales.}
\paragraph{FC-0079 [hard]} State the main tradeoff and production pitfall for Reciprocal rank fusion.
\noindent\textbf{Answer.} Tradeoff: It is robust without score calibration but ignores the magnitude and confidence of scores. Pitfall: Using an overly small rank constant can overreward noisy top results.
\par\textit{Tradeoff: It is robust without score calibration but ignores the magnitude and confidence of scores. Pitfall: Using an overly small rank constant can overreward noisy top results.}
\paragraph{LA-0080 [hard]} Explain Reciprocal rank fusion from first principles, then show how you would evaluate and operate it in production.
\noindent\textbf{Answer.} Start with the mechanism: RRF combines ranked lists using reciprocal rank positions rather than raw score scales. Make the tradeoff explicit: It is robust without score calibration but ignores the magnitude and confidence of scores. Define workload-specific offline metrics and a latency/cost budget, test important slices, instrument the critical path, and create a rollback criterion. The production review must explicitly guard against this failure: Using an overly small rank constant can overreward noisy top results.
\par\textit{A strong answer connects mechanism, measurable tradeoffs, failure isolation, observability, and rollback rather than listing product features.}
\paragraph{MCQ-0081 [medium]} Which statement best describes ANN recall?
\begin{enumerate}[label=\Alph*.]
\item Learned sparse models retain token-aligned inverted-index retrieval while expanding or reweighting terms.
\item A probabilistic lexical ranker that rewards term frequency while saturating repeated terms and normalizing document length.
\item Approximate nearest-neighbor recall measures how many true nearest items are recovered by an index configuration.
\item A sparse weighting scheme that combines within-document frequency with inverse collection frequency.
\end{enumerate}
\noindent\textbf{Answer.} Approximate nearest-neighbor recall measures how many true nearest items are recovered by an index configuration.
\par\textit{Approximate nearest-neighbor recall measures how many true nearest items are recovered by an index configuration. Tradeoff: Higher recall usually costs memory, build time, or query latency. Watch for: Reporting application answer quality without measuring index recall hides retrieval infrastructure defects.}
\paragraph{MCQ-0082 [hard]} What is the central engineering tradeoff of ANN recall?
\begin{enumerate}[label=\Alph*.]
\item Higher recall usually costs memory, build time, or query latency.
\item They bridge lexical and semantic matching but increase index size and model cost.
\item It is fast, interpretable, and exact for rare tokens but cannot directly match paraphrases.
\item It is transparent and cheap but lacks BM25's term-frequency saturation and tuned length normalization.
\end{enumerate}
\noindent\textbf{Answer.} Higher recall usually costs memory, build time, or query latency.
\par\textit{Approximate nearest-neighbor recall measures how many true nearest items are recovered by an index configuration. Tradeoff: Higher recall usually costs memory, build time, or query latency. Watch for: Reporting application answer quality without measuring index recall hides retrieval infrastructure defects.}
\paragraph{MCQ-0083 [hard]} Which failure mode should an interviewer expect you to identify for ANN recall?
\begin{enumerate}[label=\Alph*.]
\item Treating it as obsolete removes a strong baseline and hurts identifier-heavy queries.
\item Reporting application answer quality without measuring index recall hides retrieval infrastructure defects.
\item Comparing raw cosine scores across changing corpora without rebuilding IDF causes drift.
\item Unbounded expansion creates latency and storage blowups.
\end{enumerate}
\noindent\textbf{Answer.} Reporting application answer quality without measuring index recall hides retrieval infrastructure defects.
\par\textit{Approximate nearest-neighbor recall measures how many true nearest items are recovered by an index configuration. Tradeoff: Higher recall usually costs memory, build time, or query latency. Watch for: Reporting application answer quality without measuring index recall hides retrieval infrastructure defects.}
\paragraph{MCQ-0084 [medium]} A design review is specifically concerned with this risk: Reporting application answer quality without measuring index recall hides retrieval infrastructure defects. Which topic should be reviewed first?
\begin{enumerate}[label=\Alph*.]
\item TF-IDF
\item Sparse neural retrieval
\item ANN recall
\item BM25
\end{enumerate}
\noindent\textbf{Answer.} ANN recall
\par\textit{Approximate nearest-neighbor recall measures how many true nearest items are recovered by an index configuration. Tradeoff: Higher recall usually costs memory, build time, or query latency. Watch for: Reporting application answer quality without measuring index recall hides retrieval infrastructure defects.}
\paragraph{MCQ-0085 [hard]} Which design note most accurately balances the benefits and costs of ANN recall?
\begin{enumerate}[label=\Alph*.]
\item It is fast, interpretable, and exact for rare tokens but cannot directly match paraphrases.
\item It is transparent and cheap but lacks BM25's term-frequency saturation and tuned length normalization.
\item Higher recall usually costs memory, build time, or query latency.
\item They bridge lexical and semantic matching but increase index size and model cost.
\end{enumerate}
\noindent\textbf{Answer.} Higher recall usually costs memory, build time, or query latency.
\par\textit{Approximate nearest-neighbor recall measures how many true nearest items are recovered by an index configuration. Tradeoff: Higher recall usually costs memory, build time, or query latency. Watch for: Reporting application answer quality without measuring index recall hides retrieval infrastructure defects.}
\paragraph{MCQ-0086 [medium]} Which explanation would be strongest in a system-design interview about ANN recall?
\begin{enumerate}[label=\Alph*.]
\item Learned sparse models retain token-aligned inverted-index retrieval while expanding or reweighting terms.
\item A probabilistic lexical ranker that rewards term frequency while saturating repeated terms and normalizing document length.
\item A sparse weighting scheme that combines within-document frequency with inverse collection frequency.
\item Approximate nearest-neighbor recall measures how many true nearest items are recovered by an index configuration.
\end{enumerate}
\noindent\textbf{Answer.} Approximate nearest-neighbor recall measures how many true nearest items are recovered by an index configuration.
\par\textit{Approximate nearest-neighbor recall measures how many true nearest items are recovered by an index configuration. Tradeoff: Higher recall usually costs memory, build time, or query latency. Watch for: Reporting application answer quality without measuring index recall hides retrieval infrastructure defects.}
\paragraph{MCQ-0087 [hard]} Which production warning is most directly associated with ANN recall?
\begin{enumerate}[label=\Alph*.]
\item Reporting application answer quality without measuring index recall hides retrieval infrastructure defects.
\item Unbounded expansion creates latency and storage blowups.
\item Treating it as obsolete removes a strong baseline and hurts identifier-heavy queries.
\item Comparing raw cosine scores across changing corpora without rebuilding IDF causes drift.
\end{enumerate}
\noindent\textbf{Answer.} Reporting application answer quality without measuring index recall hides retrieval infrastructure defects.
\par\textit{Approximate nearest-neighbor recall measures how many true nearest items are recovered by an index configuration. Tradeoff: Higher recall usually costs memory, build time, or query latency. Watch for: Reporting application answer quality without measuring index recall hides retrieval infrastructure defects.}
\paragraph{FC-0088 [medium]} Define ANN recall in interview-ready language.
\noindent\textbf{Answer.} Approximate nearest-neighbor recall measures how many true nearest items are recovered by an index configuration.
\par\textit{Approximate nearest-neighbor recall measures how many true nearest items are recovered by an index configuration.}
\paragraph{FC-0089 [hard]} State the main tradeoff and production pitfall for ANN recall.
\noindent\textbf{Answer.} Tradeoff: Higher recall usually costs memory, build time, or query latency. Pitfall: Reporting application answer quality without measuring index recall hides retrieval infrastructure defects.
\par\textit{Tradeoff: Higher recall usually costs memory, build time, or query latency. Pitfall: Reporting application answer quality without measuring index recall hides retrieval infrastructure defects.}
\paragraph{LA-0090 [hard]} Explain ANN recall from first principles, then show how you would evaluate and operate it in production.
\noindent\textbf{Answer.} Start with the mechanism: Approximate nearest-neighbor recall measures how many true nearest items are recovered by an index configuration. Make the tradeoff explicit: Higher recall usually costs memory, build time, or query latency. Define workload-specific offline metrics and a latency/cost budget, test important slices, instrument the critical path, and create a rollback criterion. The production review must explicitly guard against this failure: Reporting application answer quality without measuring index recall hides retrieval infrastructure defects.
\par\textit{A strong answer connects mechanism, measurable tradeoffs, failure isolation, observability, and rollback rather than listing product features.}
\paragraph{MCQ-0091 [medium]} Which statement best describes Retrieval metrics?
\begin{enumerate}[label=\Alph*.]
\item A sparse weighting scheme that combines within-document frequency with inverse collection frequency.
\item A probabilistic lexical ranker that rewards term frequency while saturating repeated terms and normalizing document length.
\item RRF combines ranked lists using reciprocal rank positions rather than raw score scales.
\item Recall@k, precision@k, MRR, MAP, and nDCG measure different ranking behaviors.
\end{enumerate}
\noindent\textbf{Answer.} Recall@k, precision@k, MRR, MAP, and nDCG measure different ranking behaviors.
\par\textit{Recall@k, precision@k, MRR, MAP, and nDCG measure different ranking behaviors. Tradeoff: A metric must match the product objective and judgment granularity. Watch for: Optimizing MRR for a task that consumes many passages can worsen context coverage.}
\paragraph{MCQ-0092 [hard]} What is the central engineering tradeoff of Retrieval metrics?
\begin{enumerate}[label=\Alph*.]
\item It is robust without score calibration but ignores the magnitude and confidence of scores.
\item It is fast, interpretable, and exact for rare tokens but cannot directly match paraphrases.
\item It is transparent and cheap but lacks BM25's term-frequency saturation and tuned length normalization.
\item A metric must match the product objective and judgment granularity.
\end{enumerate}
\noindent\textbf{Answer.} A metric must match the product objective and judgment granularity.
\par\textit{Recall@k, precision@k, MRR, MAP, and nDCG measure different ranking behaviors. Tradeoff: A metric must match the product objective and judgment granularity. Watch for: Optimizing MRR for a task that consumes many passages can worsen context coverage.}
\paragraph{MCQ-0093 [hard]} Which failure mode should an interviewer expect you to identify for Retrieval metrics?
\begin{enumerate}[label=\Alph*.]
\item Optimizing MRR for a task that consumes many passages can worsen context coverage.
\item Comparing raw cosine scores across changing corpora without rebuilding IDF causes drift.
\item Treating it as obsolete removes a strong baseline and hurts identifier-heavy queries.
\item Using an overly small rank constant can overreward noisy top results.
\end{enumerate}
\noindent\textbf{Answer.} Optimizing MRR for a task that consumes many passages can worsen context coverage.
\par\textit{Recall@k, precision@k, MRR, MAP, and nDCG measure different ranking behaviors. Tradeoff: A metric must match the product objective and judgment granularity. Watch for: Optimizing MRR for a task that consumes many passages can worsen context coverage.}
\paragraph{MCQ-0094 [medium]} A design review is specifically concerned with this risk: Optimizing MRR for a task that consumes many passages can worsen context coverage. Which topic should be reviewed first?
\begin{enumerate}[label=\Alph*.]
\item BM25
\item Retrieval metrics
\item TF-IDF
\item Reciprocal rank fusion
\end{enumerate}
\noindent\textbf{Answer.} Retrieval metrics
\par\textit{Recall@k, precision@k, MRR, MAP, and nDCG measure different ranking behaviors. Tradeoff: A metric must match the product objective and judgment granularity. Watch for: Optimizing MRR for a task that consumes many passages can worsen context coverage.}
\paragraph{MCQ-0095 [hard]} Which design note most accurately balances the benefits and costs of Retrieval metrics?
\begin{enumerate}[label=\Alph*.]
\item It is robust without score calibration but ignores the magnitude and confidence of scores.
\item A metric must match the product objective and judgment granularity.
\item It is transparent and cheap but lacks BM25's term-frequency saturation and tuned length normalization.
\item It is fast, interpretable, and exact for rare tokens but cannot directly match paraphrases.
\end{enumerate}
\noindent\textbf{Answer.} A metric must match the product objective and judgment granularity.
\par\textit{Recall@k, precision@k, MRR, MAP, and nDCG measure different ranking behaviors. Tradeoff: A metric must match the product objective and judgment granularity. Watch for: Optimizing MRR for a task that consumes many passages can worsen context coverage.}
\paragraph{MCQ-0096 [medium]} Which explanation would be strongest in a system-design interview about Retrieval metrics?
\begin{enumerate}[label=\Alph*.]
\item A probabilistic lexical ranker that rewards term frequency while saturating repeated terms and normalizing document length.
\item RRF combines ranked lists using reciprocal rank positions rather than raw score scales.
\item A sparse weighting scheme that combines within-document frequency with inverse collection frequency.
\item Recall@k, precision@k, MRR, MAP, and nDCG measure different ranking behaviors.
\end{enumerate}
\noindent\textbf{Answer.} Recall@k, precision@k, MRR, MAP, and nDCG measure different ranking behaviors.
\par\textit{Recall@k, precision@k, MRR, MAP, and nDCG measure different ranking behaviors. Tradeoff: A metric must match the product objective and judgment granularity. Watch for: Optimizing MRR for a task that consumes many passages can worsen context coverage.}
\paragraph{MCQ-0097 [hard]} Which production warning is most directly associated with Retrieval metrics?
\begin{enumerate}[label=\Alph*.]
\item Comparing raw cosine scores across changing corpora without rebuilding IDF causes drift.
\item Using an overly small rank constant can overreward noisy top results.
\item Treating it as obsolete removes a strong baseline and hurts identifier-heavy queries.
\item Optimizing MRR for a task that consumes many passages can worsen context coverage.
\end{enumerate}
\noindent\textbf{Answer.} Optimizing MRR for a task that consumes many passages can worsen context coverage.
\par\textit{Recall@k, precision@k, MRR, MAP, and nDCG measure different ranking behaviors. Tradeoff: A metric must match the product objective and judgment granularity. Watch for: Optimizing MRR for a task that consumes many passages can worsen context coverage.}
\paragraph{FC-0098 [medium]} Define Retrieval metrics in interview-ready language.
\noindent\textbf{Answer.} Recall@k, precision@k, MRR, MAP, and nDCG measure different ranking behaviors.
\par\textit{Recall@k, precision@k, MRR, MAP, and nDCG measure different ranking behaviors.}
\paragraph{FC-0099 [hard]} State the main tradeoff and production pitfall for Retrieval metrics.
\noindent\textbf{Answer.} Tradeoff: A metric must match the product objective and judgment granularity. Pitfall: Optimizing MRR for a task that consumes many passages can worsen context coverage.
\par\textit{Tradeoff: A metric must match the product objective and judgment granularity. Pitfall: Optimizing MRR for a task that consumes many passages can worsen context coverage.}
\paragraph{LA-0100 [hard]} Explain Retrieval metrics from first principles, then show how you would evaluate and operate it in production.
\noindent\textbf{Answer.} Start with the mechanism: Recall@k, precision@k, MRR, MAP, and nDCG measure different ranking behaviors. Make the tradeoff explicit: A metric must match the product objective and judgment granularity. Define workload-specific offline metrics and a latency/cost budget, test important slices, instrument the critical path, and create a rollback criterion. The production review must explicitly guard against this failure: Optimizing MRR for a task that consumes many passages can worsen context coverage.
\par\textit{A strong answer connects mechanism, measurable tradeoffs, failure isolation, observability, and rollback rather than listing product features.}
\paragraph{MCQ-0101 [medium]} Which statement best describes Embedding geometry?
\begin{enumerate}[label=\Alph*.]
\item Similarity functions, normalization, anisotropy, dimensionality, and training objectives shape neighborhood quality.
\item Learned sparse models retain token-aligned inverted-index retrieval while expanding or reweighting terms.
\item A sparse weighting scheme that combines within-document frequency with inverse collection frequency.
\item A bi-encoder maps queries and passages into a shared embedding space for nearest-neighbor search.
\end{enumerate}
\noindent\textbf{Answer.} Similarity functions, normalization, anisotropy, dimensionality, and training objectives shape neighborhood quality.
\par\textit{Similarity functions, normalization, anisotropy, dimensionality, and training objectives shape neighborhood quality. Tradeoff: Cosine, dot product, and Euclidean distance can be equivalent only under specific normalization assumptions. Watch for: Mixing an index metric with a differently trained embedding objective silently harms ranking.}
\paragraph{MCQ-0102 [hard]} What is the central engineering tradeoff of Embedding geometry?
\begin{enumerate}[label=\Alph*.]
\item They bridge lexical and semantic matching but increase index size and model cost.
\item Cosine, dot product, and Euclidean distance can be equivalent only under specific normalization assumptions.
\item Semantic recall improves, while exact names, numbers, and out-of-domain language may regress.
\item It is transparent and cheap but lacks BM25's term-frequency saturation and tuned length normalization.
\end{enumerate}
\noindent\textbf{Answer.} Cosine, dot product, and Euclidean distance can be equivalent only under specific normalization assumptions.
\par\textit{Similarity functions, normalization, anisotropy, dimensionality, and training objectives shape neighborhood quality. Tradeoff: Cosine, dot product, and Euclidean distance can be equivalent only under specific normalization assumptions. Watch for: Mixing an index metric with a differently trained embedding objective silently harms ranking.}
\paragraph{MCQ-0103 [hard]} Which failure mode should an interviewer expect you to identify for Embedding geometry?
\begin{enumerate}[label=\Alph*.]
\item Unbounded expansion creates latency and storage blowups.
\item Mixing an index metric with a differently trained embedding objective silently harms ranking.
\item Evaluating only on in-domain examples hides lexical misses and embedding drift.
\item Comparing raw cosine scores across changing corpora without rebuilding IDF causes drift.
\end{enumerate}
\noindent\textbf{Answer.} Mixing an index metric with a differently trained embedding objective silently harms ranking.
\par\textit{Similarity functions, normalization, anisotropy, dimensionality, and training objectives shape neighborhood quality. Tradeoff: Cosine, dot product, and Euclidean distance can be equivalent only under specific normalization assumptions. Watch for: Mixing an index metric with a differently trained embedding objective silently harms ranking.}
\paragraph{MCQ-0104 [medium]} A design review is specifically concerned with this risk: Mixing an index metric with a differently trained embedding objective silently harms ranking. Which topic should be reviewed first?
\begin{enumerate}[label=\Alph*.]
\item Embedding geometry
\item Sparse neural retrieval
\item TF-IDF
\item Dense retrieval
\end{enumerate}
\noindent\textbf{Answer.} Embedding geometry
\par\textit{Similarity functions, normalization, anisotropy, dimensionality, and training objectives shape neighborhood quality. Tradeoff: Cosine, dot product, and Euclidean distance can be equivalent only under specific normalization assumptions. Watch for: Mixing an index metric with a differently trained embedding objective silently harms ranking.}
\paragraph{MCQ-0105 [hard]} Which design note most accurately balances the benefits and costs of Embedding geometry?
\begin{enumerate}[label=\Alph*.]
\item Cosine, dot product, and Euclidean distance can be equivalent only under specific normalization assumptions.
\item It is transparent and cheap but lacks BM25's term-frequency saturation and tuned length normalization.
\item Semantic recall improves, while exact names, numbers, and out-of-domain language may regress.
\item They bridge lexical and semantic matching but increase index size and model cost.
\end{enumerate}
\noindent\textbf{Answer.} Cosine, dot product, and Euclidean distance can be equivalent only under specific normalization assumptions.
\par\textit{Similarity functions, normalization, anisotropy, dimensionality, and training objectives shape neighborhood quality. Tradeoff: Cosine, dot product, and Euclidean distance can be equivalent only under specific normalization assumptions. Watch for: Mixing an index metric with a differently trained embedding objective silently harms ranking.}
\paragraph{MCQ-0106 [medium]} Which explanation would be strongest in a system-design interview about Embedding geometry?
\begin{enumerate}[label=\Alph*.]
\item Learned sparse models retain token-aligned inverted-index retrieval while expanding or reweighting terms.
\item A sparse weighting scheme that combines within-document frequency with inverse collection frequency.
\item Similarity functions, normalization, anisotropy, dimensionality, and training objectives shape neighborhood quality.
\item A bi-encoder maps queries and passages into a shared embedding space for nearest-neighbor search.
\end{enumerate}
\noindent\textbf{Answer.} Similarity functions, normalization, anisotropy, dimensionality, and training objectives shape neighborhood quality.
\par\textit{Similarity functions, normalization, anisotropy, dimensionality, and training objectives shape neighborhood quality. Tradeoff: Cosine, dot product, and Euclidean distance can be equivalent only under specific normalization assumptions. Watch for: Mixing an index metric with a differently trained embedding objective silently harms ranking.}
\paragraph{MCQ-0107 [hard]} Which production warning is most directly associated with Embedding geometry?
\begin{enumerate}[label=\Alph*.]
\item Comparing raw cosine scores across changing corpora without rebuilding IDF causes drift.
\item Unbounded expansion creates latency and storage blowups.
\item Mixing an index metric with a differently trained embedding objective silently harms ranking.
\item Evaluating only on in-domain examples hides lexical misses and embedding drift.
\end{enumerate}
\noindent\textbf{Answer.} Mixing an index metric with a differently trained embedding objective silently harms ranking.
\par\textit{Similarity functions, normalization, anisotropy, dimensionality, and training objectives shape neighborhood quality. Tradeoff: Cosine, dot product, and Euclidean distance can be equivalent only under specific normalization assumptions. Watch for: Mixing an index metric with a differently trained embedding objective silently harms ranking.}
\paragraph{FC-0108 [medium]} Define Embedding geometry in interview-ready language.
\noindent\textbf{Answer.} Similarity functions, normalization, anisotropy, dimensionality, and training objectives shape neighborhood quality.
\par\textit{Similarity functions, normalization, anisotropy, dimensionality, and training objectives shape neighborhood quality.}
\paragraph{FC-0109 [hard]} State the main tradeoff and production pitfall for Embedding geometry.
\noindent\textbf{Answer.} Tradeoff: Cosine, dot product, and Euclidean distance can be equivalent only under specific normalization assumptions. Pitfall: Mixing an index metric with a differently trained embedding objective silently harms ranking.
\par\textit{Tradeoff: Cosine, dot product, and Euclidean distance can be equivalent only under specific normalization assumptions. Pitfall: Mixing an index metric with a differently trained embedding objective silently harms ranking.}
\paragraph{LA-0110 [hard]} Explain Embedding geometry from first principles, then show how you would evaluate and operate it in production.
\noindent\textbf{Answer.} Start with the mechanism: Similarity functions, normalization, anisotropy, dimensionality, and training objectives shape neighborhood quality. Make the tradeoff explicit: Cosine, dot product, and Euclidean distance can be equivalent only under specific normalization assumptions. Define workload-specific offline metrics and a latency/cost budget, test important slices, instrument the critical path, and create a rollback criterion. The production review must explicitly guard against this failure: Mixing an index metric with a differently trained embedding objective silently harms ranking.
\par\textit{A strong answer connects mechanism, measurable tradeoffs, failure isolation, observability, and rollback rather than listing product features.}
\paragraph{MCQ-0111 [medium]} Which statement best describes Query expansion?
\begin{enumerate}[label=\Alph*.]
\item A model jointly encodes query and candidate text to assign a relevance score.
\item Expansion adds synonyms, generated subqueries, or pseudo-relevance terms before retrieval.
\item A probabilistic lexical ranker that rewards term frequency while saturating repeated terms and normalizing document length.
\item Recall@k, precision@k, MRR, MAP, and nDCG measure different ranking behaviors.
\end{enumerate}
\noindent\textbf{Answer.} Expansion adds synonyms, generated subqueries, or pseudo-relevance terms before retrieval.
\par\textit{Expansion adds synonyms, generated subqueries, or pseudo-relevance terms before retrieval. Tradeoff: Recall can improve at the expense of latency and topic drift. Watch for: Letting an LLM expand unconstrained queries can inject unsupported intent.}
\paragraph{MCQ-0112 [hard]} What is the central engineering tradeoff of Query expansion?
\begin{enumerate}[label=\Alph*.]
\item It gives strong precision on a small candidate set but is too expensive for full-corpus retrieval.
\item A metric must match the product objective and judgment granularity.
\item It is fast, interpretable, and exact for rare tokens but cannot directly match paraphrases.
\item Recall can improve at the expense of latency and topic drift.
\end{enumerate}
\noindent\textbf{Answer.} Recall can improve at the expense of latency and topic drift.
\par\textit{Expansion adds synonyms, generated subqueries, or pseudo-relevance terms before retrieval. Tradeoff: Recall can improve at the expense of latency and topic drift. Watch for: Letting an LLM expand unconstrained queries can inject unsupported intent.}
\paragraph{MCQ-0113 [hard]} Which failure mode should an interviewer expect you to identify for Query expansion?
\begin{enumerate}[label=\Alph*.]
\item Optimizing MRR for a task that consumes many passages can worsen context coverage.
\item Treating it as obsolete removes a strong baseline and hurts identifier-heavy queries.
\item Letting an LLM expand unconstrained queries can inject unsupported intent.
\item Reranking too many long passages dominates p95 latency.
\end{enumerate}
\noindent\textbf{Answer.} Letting an LLM expand unconstrained queries can inject unsupported intent.
\par\textit{Expansion adds synonyms, generated subqueries, or pseudo-relevance terms before retrieval. Tradeoff: Recall can improve at the expense of latency and topic drift. Watch for: Letting an LLM expand unconstrained queries can inject unsupported intent.}
\paragraph{MCQ-0114 [medium]} A design review is specifically concerned with this risk: Letting an LLM expand unconstrained queries can inject unsupported intent. Which topic should be reviewed first?
\begin{enumerate}[label=\Alph*.]
\item BM25
\item Query expansion
\item Retrieval metrics
\item Cross-encoder reranking
\end{enumerate}
\noindent\textbf{Answer.} Query expansion
\par\textit{Expansion adds synonyms, generated subqueries, or pseudo-relevance terms before retrieval. Tradeoff: Recall can improve at the expense of latency and topic drift. Watch for: Letting an LLM expand unconstrained queries can inject unsupported intent.}
\paragraph{MCQ-0115 [hard]} Which design note most accurately balances the benefits and costs of Query expansion?
\begin{enumerate}[label=\Alph*.]
\item It gives strong precision on a small candidate set but is too expensive for full-corpus retrieval.
\item Recall can improve at the expense of latency and topic drift.
\item It is fast, interpretable, and exact for rare tokens but cannot directly match paraphrases.
\item A metric must match the product objective and judgment granularity.
\end{enumerate}
\noindent\textbf{Answer.} Recall can improve at the expense of latency and topic drift.
\par\textit{Expansion adds synonyms, generated subqueries, or pseudo-relevance terms before retrieval. Tradeoff: Recall can improve at the expense of latency and topic drift. Watch for: Letting an LLM expand unconstrained queries can inject unsupported intent.}
\paragraph{MCQ-0116 [medium]} Which explanation would be strongest in a system-design interview about Query expansion?
\begin{enumerate}[label=\Alph*.]
\item A model jointly encodes query and candidate text to assign a relevance score.
\item A probabilistic lexical ranker that rewards term frequency while saturating repeated terms and normalizing document length.
\item Recall@k, precision@k, MRR, MAP, and nDCG measure different ranking behaviors.
\item Expansion adds synonyms, generated subqueries, or pseudo-relevance terms before retrieval.
\end{enumerate}
\noindent\textbf{Answer.} Expansion adds synonyms, generated subqueries, or pseudo-relevance terms before retrieval.
\par\textit{Expansion adds synonyms, generated subqueries, or pseudo-relevance terms before retrieval. Tradeoff: Recall can improve at the expense of latency and topic drift. Watch for: Letting an LLM expand unconstrained queries can inject unsupported intent.}
\paragraph{MCQ-0117 [hard]} Which production warning is most directly associated with Query expansion?
\begin{enumerate}[label=\Alph*.]
\item Reranking too many long passages dominates p95 latency.
\item Optimizing MRR for a task that consumes many passages can worsen context coverage.
\item Letting an LLM expand unconstrained queries can inject unsupported intent.
\item Treating it as obsolete removes a strong baseline and hurts identifier-heavy queries.
\end{enumerate}
\noindent\textbf{Answer.} Letting an LLM expand unconstrained queries can inject unsupported intent.
\par\textit{Expansion adds synonyms, generated subqueries, or pseudo-relevance terms before retrieval. Tradeoff: Recall can improve at the expense of latency and topic drift. Watch for: Letting an LLM expand unconstrained queries can inject unsupported intent.}
\paragraph{FC-0118 [medium]} Define Query expansion in interview-ready language.
\noindent\textbf{Answer.} Expansion adds synonyms, generated subqueries, or pseudo-relevance terms before retrieval.
\par\textit{Expansion adds synonyms, generated subqueries, or pseudo-relevance terms before retrieval.}
\paragraph{FC-0119 [hard]} State the main tradeoff and production pitfall for Query expansion.
\noindent\textbf{Answer.} Tradeoff: Recall can improve at the expense of latency and topic drift. Pitfall: Letting an LLM expand unconstrained queries can inject unsupported intent.
\par\textit{Tradeoff: Recall can improve at the expense of latency and topic drift. Pitfall: Letting an LLM expand unconstrained queries can inject unsupported intent.}
\paragraph{LA-0120 [hard]} Explain Query expansion from first principles, then show how you would evaluate and operate it in production.
\noindent\textbf{Answer.} Start with the mechanism: Expansion adds synonyms, generated subqueries, or pseudo-relevance terms before retrieval. Make the tradeoff explicit: Recall can improve at the expense of latency and topic drift. Define workload-specific offline metrics and a latency/cost budget, test important slices, instrument the critical path, and create a rollback criterion. The production review must explicitly guard against this failure: Letting an LLM expand unconstrained queries can inject unsupported intent.
\par\textit{A strong answer connects mechanism, measurable tradeoffs, failure isolation, observability, and rollback rather than listing product features.}
\subsection{Vector Databases and Search Engines}
\paragraph{MCQ-0121 [medium]} Which statement best describes pgvector?
\begin{enumerate}[label=\Alph*.]
\item A PostgreSQL extension that stores vectors beside relational data and supports exact, HNSW, and IVFFlat search.
\item Redis Query Engine adds vector indexes and filters to an in-memory data platform.
\item Elasticsearch combines inverted indexes, filters, aggregations, and approximate kNN in one search platform.
\item Milvus is a distributed vector database; Zilliz provides its managed service with multiple index choices.
\end{enumerate}
\noindent\textbf{Answer.} A PostgreSQL extension that stores vectors beside relational data and supports exact, HNSW, and IVFFlat search.
\par\textit{A PostgreSQL extension that stores vectors beside relational data and supports exact, HNSW, and IVFFlat search. Tradeoff: It simplifies transactions and joins, while specialized stores may scale independent vector workloads more easily. Watch for: Adding an ANN index without tuning filtering and planner behavior can reduce recall or miss the index.}
\paragraph{MCQ-0122 [hard]} What is the central engineering tradeoff of pgvector?
\begin{enumerate}[label=\Alph*.]
\item It is compelling for existing Elastic estates but vector memory and shard design require care.
\item It simplifies transactions and joins, while specialized stores may scale independent vector workloads more easily.
\item It targets large-scale separation of compute and storage, with more operational concepts than embedded options.
\item Very low latency is possible, but RAM cost and persistence design matter at scale.
\end{enumerate}
\noindent\textbf{Answer.} It simplifies transactions and joins, while specialized stores may scale independent vector workloads more easily.
\par\textit{A PostgreSQL extension that stores vectors beside relational data and supports exact, HNSW, and IVFFlat search. Tradeoff: It simplifies transactions and joins, while specialized stores may scale independent vector workloads more easily. Watch for: Adding an ANN index without tuning filtering and planner behavior can reduce recall or miss the index.}
\paragraph{MCQ-0123 [hard]} Which failure mode should an interviewer expect you to identify for pgvector?
\begin{enumerate}[label=\Alph*.]
\item Oversharding fragments candidate pools and increases coordination overhead.
\item Using Redis only as a vector database without capacity planning can evict critical data.
\item Selecting an index by popularity instead of dataset size and latency-recall tests wastes resources.
\item Adding an ANN index without tuning filtering and planner behavior can reduce recall or miss the index.
\end{enumerate}
\noindent\textbf{Answer.} Adding an ANN index without tuning filtering and planner behavior can reduce recall or miss the index.
\par\textit{A PostgreSQL extension that stores vectors beside relational data and supports exact, HNSW, and IVFFlat search. Tradeoff: It simplifies transactions and joins, while specialized stores may scale independent vector workloads more easily. Watch for: Adding an ANN index without tuning filtering and planner behavior can reduce recall or miss the index.}
\paragraph{MCQ-0124 [medium]} A design review is specifically concerned with this risk: Adding an ANN index without tuning filtering and planner behavior can reduce recall or miss the index. Which topic should be reviewed first?
\begin{enumerate}[label=\Alph*.]
\item Milvus and Zilliz
\item Redis vector search
\item Elasticsearch vector search
\item pgvector
\end{enumerate}
\noindent\textbf{Answer.} pgvector
\par\textit{A PostgreSQL extension that stores vectors beside relational data and supports exact, HNSW, and IVFFlat search. Tradeoff: It simplifies transactions and joins, while specialized stores may scale independent vector workloads more easily. Watch for: Adding an ANN index without tuning filtering and planner behavior can reduce recall or miss the index.}
\paragraph{MCQ-0125 [hard]} Which design note most accurately balances the benefits and costs of pgvector?
\begin{enumerate}[label=\Alph*.]
\item It targets large-scale separation of compute and storage, with more operational concepts than embedded options.
\item Very low latency is possible, but RAM cost and persistence design matter at scale.
\item It simplifies transactions and joins, while specialized stores may scale independent vector workloads more easily.
\item It is compelling for existing Elastic estates but vector memory and shard design require care.
\end{enumerate}
\noindent\textbf{Answer.} It simplifies transactions and joins, while specialized stores may scale independent vector workloads more easily.
\par\textit{A PostgreSQL extension that stores vectors beside relational data and supports exact, HNSW, and IVFFlat search. Tradeoff: It simplifies transactions and joins, while specialized stores may scale independent vector workloads more easily. Watch for: Adding an ANN index without tuning filtering and planner behavior can reduce recall or miss the index.}
\paragraph{MCQ-0126 [medium]} Which explanation would be strongest in a system-design interview about pgvector?
\begin{enumerate}[label=\Alph*.]
\item Elasticsearch combines inverted indexes, filters, aggregations, and approximate kNN in one search platform.
\item Redis Query Engine adds vector indexes and filters to an in-memory data platform.
\item Milvus is a distributed vector database; Zilliz provides its managed service with multiple index choices.
\item A PostgreSQL extension that stores vectors beside relational data and supports exact, HNSW, and IVFFlat search.
\end{enumerate}
\noindent\textbf{Answer.} A PostgreSQL extension that stores vectors beside relational data and supports exact, HNSW, and IVFFlat search.
\par\textit{A PostgreSQL extension that stores vectors beside relational data and supports exact, HNSW, and IVFFlat search. Tradeoff: It simplifies transactions and joins, while specialized stores may scale independent vector workloads more easily. Watch for: Adding an ANN index without tuning filtering and planner behavior can reduce recall or miss the index.}
\paragraph{MCQ-0127 [hard]} Which production warning is most directly associated with pgvector?
\begin{enumerate}[label=\Alph*.]
\item Selecting an index by popularity instead of dataset size and latency-recall tests wastes resources.
\item Using Redis only as a vector database without capacity planning can evict critical data.
\item Adding an ANN index without tuning filtering and planner behavior can reduce recall or miss the index.
\item Oversharding fragments candidate pools and increases coordination overhead.
\end{enumerate}
\noindent\textbf{Answer.} Adding an ANN index without tuning filtering and planner behavior can reduce recall or miss the index.
\par\textit{A PostgreSQL extension that stores vectors beside relational data and supports exact, HNSW, and IVFFlat search. Tradeoff: It simplifies transactions and joins, while specialized stores may scale independent vector workloads more easily. Watch for: Adding an ANN index without tuning filtering and planner behavior can reduce recall or miss the index.}
\paragraph{FC-0128 [medium]} Define pgvector in interview-ready language.
\noindent\textbf{Answer.} A PostgreSQL extension that stores vectors beside relational data and supports exact, HNSW, and IVFFlat search.
\par\textit{A PostgreSQL extension that stores vectors beside relational data and supports exact, HNSW, and IVFFlat search.}
\paragraph{FC-0129 [hard]} State the main tradeoff and production pitfall for pgvector.
\noindent\textbf{Answer.} Tradeoff: It simplifies transactions and joins, while specialized stores may scale independent vector workloads more easily. Pitfall: Adding an ANN index without tuning filtering and planner behavior can reduce recall or miss the index.
\par\textit{Tradeoff: It simplifies transactions and joins, while specialized stores may scale independent vector workloads more easily. Pitfall: Adding an ANN index without tuning filtering and planner behavior can reduce recall or miss the index.}
\paragraph{LA-0130 [hard]} Explain pgvector from first principles, then show how you would evaluate and operate it in production.
\noindent\textbf{Answer.} Start with the mechanism: A PostgreSQL extension that stores vectors beside relational data and supports exact, HNSW, and IVFFlat search. Make the tradeoff explicit: It simplifies transactions and joins, while specialized stores may scale independent vector workloads more easily. Define workload-specific offline metrics and a latency/cost budget, test important slices, instrument the critical path, and create a rollback criterion. The production review must explicitly guard against this failure: Adding an ANN index without tuning filtering and planner behavior can reduce recall or miss the index.
\par\textit{A strong answer connects mechanism, measurable tradeoffs, failure isolation, observability, and rollback rather than listing product features.}
\paragraph{MCQ-0131 [medium]} Which statement best describes Pinecone?
\begin{enumerate}[label=\Alph*.]
\item Elasticsearch combines inverted indexes, filters, aggregations, and approximate kNN in one search platform.
\item LanceDB uses a columnar, versioned data format and supports vector, full-text, and multimodal data locally or in cloud storage.
\item Inverted files narrow the search to clusters; product quantization compresses vectors for cheaper approximate scoring.
\item A managed vector service with metadata filtering, namespaces, and dense-sparse hybrid retrieval.
\end{enumerate}
\noindent\textbf{Answer.} A managed vector service with metadata filtering, namespaces, and dense-sparse hybrid retrieval.
\par\textit{A managed vector service with metadata filtering, namespaces, and dense-sparse hybrid retrieval. Tradeoff: Operations are outsourced, trading infrastructure control and portability for managed scale. Watch for: Poor namespace and metadata design creates noisy tenants and expensive fan-out queries.}
\paragraph{MCQ-0132 [hard]} What is the central engineering tradeoff of Pinecone?
\begin{enumerate}[label=\Alph*.]
\item Memory and latency fall, while training quality and quantization error reduce recall.
\item It fits data-lake and embedded workflows but differs operationally from a server-first database.
\item It is compelling for existing Elastic estates but vector memory and shard design require care.
\item Operations are outsourced, trading infrastructure control and portability for managed scale.
\end{enumerate}
\noindent\textbf{Answer.} Operations are outsourced, trading infrastructure control and portability for managed scale.
\par\textit{A managed vector service with metadata filtering, namespaces, and dense-sparse hybrid retrieval. Tradeoff: Operations are outsourced, trading infrastructure control and portability for managed scale. Watch for: Poor namespace and metadata design creates noisy tenants and expensive fan-out queries.}
\paragraph{MCQ-0133 [hard]} Which failure mode should an interviewer expect you to identify for Pinecone?
\begin{enumerate}[label=\Alph*.]
\item Ignoring compaction and version retention can degrade read performance and inflate storage.
\item Training centroids or codebooks on an unrepresentative sample bakes bias into the index.
\item Oversharding fragments candidate pools and increases coordination overhead.
\item Poor namespace and metadata design creates noisy tenants and expensive fan-out queries.
\end{enumerate}
\noindent\textbf{Answer.} Poor namespace and metadata design creates noisy tenants and expensive fan-out queries.
\par\textit{A managed vector service with metadata filtering, namespaces, and dense-sparse hybrid retrieval. Tradeoff: Operations are outsourced, trading infrastructure control and portability for managed scale. Watch for: Poor namespace and metadata design creates noisy tenants and expensive fan-out queries.}
\paragraph{MCQ-0134 [medium]} A design review is specifically concerned with this risk: Poor namespace and metadata design creates noisy tenants and expensive fan-out queries. Which topic should be reviewed first?
\begin{enumerate}[label=\Alph*.]
\item Elasticsearch vector search
\item Pinecone
\item LanceDB
\item IVF and PQ
\end{enumerate}
\noindent\textbf{Answer.} Pinecone
\par\textit{A managed vector service with metadata filtering, namespaces, and dense-sparse hybrid retrieval. Tradeoff: Operations are outsourced, trading infrastructure control and portability for managed scale. Watch for: Poor namespace and metadata design creates noisy tenants and expensive fan-out queries.}
\paragraph{MCQ-0135 [hard]} Which design note most accurately balances the benefits and costs of Pinecone?
\begin{enumerate}[label=\Alph*.]
\item It fits data-lake and embedded workflows but differs operationally from a server-first database.
\item It is compelling for existing Elastic estates but vector memory and shard design require care.
\item Operations are outsourced, trading infrastructure control and portability for managed scale.
\item Memory and latency fall, while training quality and quantization error reduce recall.
\end{enumerate}
\noindent\textbf{Answer.} Operations are outsourced, trading infrastructure control and portability for managed scale.
\par\textit{A managed vector service with metadata filtering, namespaces, and dense-sparse hybrid retrieval. Tradeoff: Operations are outsourced, trading infrastructure control and portability for managed scale. Watch for: Poor namespace and metadata design creates noisy tenants and expensive fan-out queries.}
\paragraph{MCQ-0136 [medium]} Which explanation would be strongest in a system-design interview about Pinecone?
\begin{enumerate}[label=\Alph*.]
\item Inverted files narrow the search to clusters; product quantization compresses vectors for cheaper approximate scoring.
\item A managed vector service with metadata filtering, namespaces, and dense-sparse hybrid retrieval.
\item Elasticsearch combines inverted indexes, filters, aggregations, and approximate kNN in one search platform.
\item LanceDB uses a columnar, versioned data format and supports vector, full-text, and multimodal data locally or in cloud storage.
\end{enumerate}
\noindent\textbf{Answer.} A managed vector service with metadata filtering, namespaces, and dense-sparse hybrid retrieval.
\par\textit{A managed vector service with metadata filtering, namespaces, and dense-sparse hybrid retrieval. Tradeoff: Operations are outsourced, trading infrastructure control and portability for managed scale. Watch for: Poor namespace and metadata design creates noisy tenants and expensive fan-out queries.}
\paragraph{MCQ-0137 [hard]} Which production warning is most directly associated with Pinecone?
\begin{enumerate}[label=\Alph*.]
\item Ignoring compaction and version retention can degrade read performance and inflate storage.
\item Oversharding fragments candidate pools and increases coordination overhead.
\item Training centroids or codebooks on an unrepresentative sample bakes bias into the index.
\item Poor namespace and metadata design creates noisy tenants and expensive fan-out queries.
\end{enumerate}
\noindent\textbf{Answer.} Poor namespace and metadata design creates noisy tenants and expensive fan-out queries.
\par\textit{A managed vector service with metadata filtering, namespaces, and dense-sparse hybrid retrieval. Tradeoff: Operations are outsourced, trading infrastructure control and portability for managed scale. Watch for: Poor namespace and metadata design creates noisy tenants and expensive fan-out queries.}
\paragraph{FC-0138 [medium]} Define Pinecone in interview-ready language.
\noindent\textbf{Answer.} A managed vector service with metadata filtering, namespaces, and dense-sparse hybrid retrieval.
\par\textit{A managed vector service with metadata filtering, namespaces, and dense-sparse hybrid retrieval.}
\paragraph{FC-0139 [hard]} State the main tradeoff and production pitfall for Pinecone.
\noindent\textbf{Answer.} Tradeoff: Operations are outsourced, trading infrastructure control and portability for managed scale. Pitfall: Poor namespace and metadata design creates noisy tenants and expensive fan-out queries.
\par\textit{Tradeoff: Operations are outsourced, trading infrastructure control and portability for managed scale. Pitfall: Poor namespace and metadata design creates noisy tenants and expensive fan-out queries.}
\paragraph{LA-0140 [hard]} Explain Pinecone from first principles, then show how you would evaluate and operate it in production.
\noindent\textbf{Answer.} Start with the mechanism: A managed vector service with metadata filtering, namespaces, and dense-sparse hybrid retrieval. Make the tradeoff explicit: Operations are outsourced, trading infrastructure control and portability for managed scale. Define workload-specific offline metrics and a latency/cost budget, test important slices, instrument the critical path, and create a rollback criterion. The production review must explicitly guard against this failure: Poor namespace and metadata design creates noisy tenants and expensive fan-out queries.
\par\textit{A strong answer connects mechanism, measurable tradeoffs, failure isolation, observability, and rollback rather than listing product features.}
\paragraph{MCQ-0141 [medium]} Which statement best describes Qdrant?
\begin{enumerate}[label=\Alph*.]
\item Inverted files narrow the search to clusters; product quantization compresses vectors for cheaper approximate scoring.
\item Elasticsearch combines inverted indexes, filters, aggregations, and approximate kNN in one search platform.
\item An open-source vector engine centered on HNSW, payload filtering, quantization, and distributed collections.
\item A PostgreSQL extension that stores vectors beside relational data and supports exact, HNSW, and IVFFlat search.
\end{enumerate}
\noindent\textbf{Answer.} An open-source vector engine centered on HNSW, payload filtering, quantization, and distributed collections.
\par\textit{An open-source vector engine centered on HNSW, payload filtering, quantization, and distributed collections. Tradeoff: It offers operational control and rich filtering but adds a stateful distributed service to run. Watch for: Applying post-filtering after ANN can return too few eligible results.}
\paragraph{MCQ-0142 [hard]} What is the central engineering tradeoff of Qdrant?
\begin{enumerate}[label=\Alph*.]
\item It offers operational control and rich filtering but adds a stateful distributed service to run.
\item Memory and latency fall, while training quality and quantization error reduce recall.
\item It simplifies transactions and joins, while specialized stores may scale independent vector workloads more easily.
\item It is compelling for existing Elastic estates but vector memory and shard design require care.
\end{enumerate}
\noindent\textbf{Answer.} It offers operational control and rich filtering but adds a stateful distributed service to run.
\par\textit{An open-source vector engine centered on HNSW, payload filtering, quantization, and distributed collections. Tradeoff: It offers operational control and rich filtering but adds a stateful distributed service to run. Watch for: Applying post-filtering after ANN can return too few eligible results.}
\paragraph{MCQ-0143 [hard]} Which failure mode should an interviewer expect you to identify for Qdrant?
\begin{enumerate}[label=\Alph*.]
\item Oversharding fragments candidate pools and increases coordination overhead.
\item Applying post-filtering after ANN can return too few eligible results.
\item Adding an ANN index without tuning filtering and planner behavior can reduce recall or miss the index.
\item Training centroids or codebooks on an unrepresentative sample bakes bias into the index.
\end{enumerate}
\noindent\textbf{Answer.} Applying post-filtering after ANN can return too few eligible results.
\par\textit{An open-source vector engine centered on HNSW, payload filtering, quantization, and distributed collections. Tradeoff: It offers operational control and rich filtering but adds a stateful distributed service to run. Watch for: Applying post-filtering after ANN can return too few eligible results.}
\paragraph{MCQ-0144 [medium]} A design review is specifically concerned with this risk: Applying post-filtering after ANN can return too few eligible results. Which topic should be reviewed first?
\begin{enumerate}[label=\Alph*.]
\item Elasticsearch vector search
\item pgvector
\item IVF and PQ
\item Qdrant
\end{enumerate}
\noindent\textbf{Answer.} Qdrant
\par\textit{An open-source vector engine centered on HNSW, payload filtering, quantization, and distributed collections. Tradeoff: It offers operational control and rich filtering but adds a stateful distributed service to run. Watch for: Applying post-filtering after ANN can return too few eligible results.}
\paragraph{MCQ-0145 [hard]} Which design note most accurately balances the benefits and costs of Qdrant?
\begin{enumerate}[label=\Alph*.]
\item It simplifies transactions and joins, while specialized stores may scale independent vector workloads more easily.
\item It is compelling for existing Elastic estates but vector memory and shard design require care.
\item Memory and latency fall, while training quality and quantization error reduce recall.
\item It offers operational control and rich filtering but adds a stateful distributed service to run.
\end{enumerate}
\noindent\textbf{Answer.} It offers operational control and rich filtering but adds a stateful distributed service to run.
\par\textit{An open-source vector engine centered on HNSW, payload filtering, quantization, and distributed collections. Tradeoff: It offers operational control and rich filtering but adds a stateful distributed service to run. Watch for: Applying post-filtering after ANN can return too few eligible results.}
\paragraph{MCQ-0146 [medium]} Which explanation would be strongest in a system-design interview about Qdrant?
\begin{enumerate}[label=\Alph*.]
\item Inverted files narrow the search to clusters; product quantization compresses vectors for cheaper approximate scoring.
\item A PostgreSQL extension that stores vectors beside relational data and supports exact, HNSW, and IVFFlat search.
\item An open-source vector engine centered on HNSW, payload filtering, quantization, and distributed collections.
\item Elasticsearch combines inverted indexes, filters, aggregations, and approximate kNN in one search platform.
\end{enumerate}
\noindent\textbf{Answer.} An open-source vector engine centered on HNSW, payload filtering, quantization, and distributed collections.
\par\textit{An open-source vector engine centered on HNSW, payload filtering, quantization, and distributed collections. Tradeoff: It offers operational control and rich filtering but adds a stateful distributed service to run. Watch for: Applying post-filtering after ANN can return too few eligible results.}
\paragraph{MCQ-0147 [hard]} Which production warning is most directly associated with Qdrant?
\begin{enumerate}[label=\Alph*.]
\item Oversharding fragments candidate pools and increases coordination overhead.
\item Applying post-filtering after ANN can return too few eligible results.
\item Training centroids or codebooks on an unrepresentative sample bakes bias into the index.
\item Adding an ANN index without tuning filtering and planner behavior can reduce recall or miss the index.
\end{enumerate}
\noindent\textbf{Answer.} Applying post-filtering after ANN can return too few eligible results.
\par\textit{An open-source vector engine centered on HNSW, payload filtering, quantization, and distributed collections. Tradeoff: It offers operational control and rich filtering but adds a stateful distributed service to run. Watch for: Applying post-filtering after ANN can return too few eligible results.}
\paragraph{FC-0148 [medium]} Define Qdrant in interview-ready language.
\noindent\textbf{Answer.} An open-source vector engine centered on HNSW, payload filtering, quantization, and distributed collections.
\par\textit{An open-source vector engine centered on HNSW, payload filtering, quantization, and distributed collections.}
\paragraph{FC-0149 [hard]} State the main tradeoff and production pitfall for Qdrant.
\noindent\textbf{Answer.} Tradeoff: It offers operational control and rich filtering but adds a stateful distributed service to run. Pitfall: Applying post-filtering after ANN can return too few eligible results.
\par\textit{Tradeoff: It offers operational control and rich filtering but adds a stateful distributed service to run. Pitfall: Applying post-filtering after ANN can return too few eligible results.}
\paragraph{LA-0150 [hard]} Explain Qdrant from first principles, then show how you would evaluate and operate it in production.
\noindent\textbf{Answer.} Start with the mechanism: An open-source vector engine centered on HNSW, payload filtering, quantization, and distributed collections. Make the tradeoff explicit: It offers operational control and rich filtering but adds a stateful distributed service to run. Define workload-specific offline metrics and a latency/cost budget, test important slices, instrument the critical path, and create a rollback criterion. The production review must explicitly guard against this failure: Applying post-filtering after ANN can return too few eligible results.
\par\textit{A strong answer connects mechanism, measurable tradeoffs, failure isolation, observability, and rollback rather than listing product features.}
\paragraph{MCQ-0151 [medium]} Which statement best describes Weaviate?
\begin{enumerate}[label=\Alph*.]
\item A navigable multilayer proximity graph controlled by construction degree and search breadth.
\item A vector database with collections, hybrid search, filters, modules, and multi-tenancy features.
\item Milvus is a distributed vector database; Zilliz provides its managed service with multiple index choices.
\item OpenSearch supports lexical and vector workloads with pluggable kNN engines and search pipelines.
\end{enumerate}
\noindent\textbf{Answer.} A vector database with collections, hybrid search, filters, modules, and multi-tenancy features.
\par\textit{A vector database with collections, hybrid search, filters, modules, and multi-tenancy features. Tradeoff: An integrated schema accelerates product work but may couple ingestion and inference choices. Watch for: Automatic vectorization without explicit model version metadata makes reindexing unsafe.}
\paragraph{MCQ-0152 [hard]} What is the central engineering tradeoff of Weaviate?
\begin{enumerate}[label=\Alph*.]
\item An integrated schema accelerates product work but may couple ingestion and inference choices.
\item Open governance and integrated search are attractive, while engine-specific behavior complicates tuning.
\item It offers strong latency-recall behavior with significant memory and build cost.
\item It targets large-scale separation of compute and storage, with more operational concepts than embedded options.
\end{enumerate}
\noindent\textbf{Answer.} An integrated schema accelerates product work but may couple ingestion and inference choices.
\par\textit{A vector database with collections, hybrid search, filters, modules, and multi-tenancy features. Tradeoff: An integrated schema accelerates product work but may couple ingestion and inference choices. Watch for: Automatic vectorization without explicit model version metadata makes reindexing unsafe.}
\paragraph{MCQ-0153 [hard]} Which failure mode should an interviewer expect you to identify for Weaviate?
\begin{enumerate}[label=\Alph*.]
\item Deleting or heavily filtering data without maintenance degrades graph quality.
\item Selecting an index by popularity instead of dataset size and latency-recall tests wastes resources.
\item Automatic vectorization without explicit model version metadata makes reindexing unsafe.
\item Assuming identical parameters across FAISS, Lucene, and NMSLIB backends produces misleading tests.
\end{enumerate}
\noindent\textbf{Answer.} Automatic vectorization without explicit model version metadata makes reindexing unsafe.
\par\textit{A vector database with collections, hybrid search, filters, modules, and multi-tenancy features. Tradeoff: An integrated schema accelerates product work but may couple ingestion and inference choices. Watch for: Automatic vectorization without explicit model version metadata makes reindexing unsafe.}
\paragraph{MCQ-0154 [medium]} A design review is specifically concerned with this risk: Automatic vectorization without explicit model version metadata makes reindexing unsafe. Which topic should be reviewed first?
\begin{enumerate}[label=\Alph*.]
\item Milvus and Zilliz
\item Weaviate
\item HNSW
\item OpenSearch vector search
\end{enumerate}
\noindent\textbf{Answer.} Weaviate
\par\textit{A vector database with collections, hybrid search, filters, modules, and multi-tenancy features. Tradeoff: An integrated schema accelerates product work but may couple ingestion and inference choices. Watch for: Automatic vectorization without explicit model version metadata makes reindexing unsafe.}
\paragraph{MCQ-0155 [hard]} Which design note most accurately balances the benefits and costs of Weaviate?
\begin{enumerate}[label=\Alph*.]
\item An integrated schema accelerates product work but may couple ingestion and inference choices.
\item Open governance and integrated search are attractive, while engine-specific behavior complicates tuning.
\item It targets large-scale separation of compute and storage, with more operational concepts than embedded options.
\item It offers strong latency-recall behavior with significant memory and build cost.
\end{enumerate}
\noindent\textbf{Answer.} An integrated schema accelerates product work but may couple ingestion and inference choices.
\par\textit{A vector database with collections, hybrid search, filters, modules, and multi-tenancy features. Tradeoff: An integrated schema accelerates product work but may couple ingestion and inference choices. Watch for: Automatic vectorization without explicit model version metadata makes reindexing unsafe.}
\paragraph{MCQ-0156 [medium]} Which explanation would be strongest in a system-design interview about Weaviate?
\begin{enumerate}[label=\Alph*.]
\item A navigable multilayer proximity graph controlled by construction degree and search breadth.
\item A vector database with collections, hybrid search, filters, modules, and multi-tenancy features.
\item Milvus is a distributed vector database; Zilliz provides its managed service with multiple index choices.
\item OpenSearch supports lexical and vector workloads with pluggable kNN engines and search pipelines.
\end{enumerate}
\noindent\textbf{Answer.} A vector database with collections, hybrid search, filters, modules, and multi-tenancy features.
\par\textit{A vector database with collections, hybrid search, filters, modules, and multi-tenancy features. Tradeoff: An integrated schema accelerates product work but may couple ingestion and inference choices. Watch for: Automatic vectorization without explicit model version metadata makes reindexing unsafe.}
\paragraph{MCQ-0157 [hard]} Which production warning is most directly associated with Weaviate?
\begin{enumerate}[label=\Alph*.]
\item Deleting or heavily filtering data without maintenance degrades graph quality.
\item Assuming identical parameters across FAISS, Lucene, and NMSLIB backends produces misleading tests.
\item Automatic vectorization without explicit model version metadata makes reindexing unsafe.
\item Selecting an index by popularity instead of dataset size and latency-recall tests wastes resources.
\end{enumerate}
\noindent\textbf{Answer.} Automatic vectorization without explicit model version metadata makes reindexing unsafe.
\par\textit{A vector database with collections, hybrid search, filters, modules, and multi-tenancy features. Tradeoff: An integrated schema accelerates product work but may couple ingestion and inference choices. Watch for: Automatic vectorization without explicit model version metadata makes reindexing unsafe.}
\paragraph{FC-0158 [medium]} Define Weaviate in interview-ready language.
\noindent\textbf{Answer.} A vector database with collections, hybrid search, filters, modules, and multi-tenancy features.
\par\textit{A vector database with collections, hybrid search, filters, modules, and multi-tenancy features.}
\paragraph{FC-0159 [hard]} State the main tradeoff and production pitfall for Weaviate.
\noindent\textbf{Answer.} Tradeoff: An integrated schema accelerates product work but may couple ingestion and inference choices. Pitfall: Automatic vectorization without explicit model version metadata makes reindexing unsafe.
\par\textit{Tradeoff: An integrated schema accelerates product work but may couple ingestion and inference choices. Pitfall: Automatic vectorization without explicit model version metadata makes reindexing unsafe.}
\paragraph{LA-0160 [hard]} Explain Weaviate from first principles, then show how you would evaluate and operate it in production.
\noindent\textbf{Answer.} Start with the mechanism: A vector database with collections, hybrid search, filters, modules, and multi-tenancy features. Make the tradeoff explicit: An integrated schema accelerates product work but may couple ingestion and inference choices. Define workload-specific offline metrics and a latency/cost budget, test important slices, instrument the critical path, and create a rollback criterion. The production review must explicitly guard against this failure: Automatic vectorization without explicit model version metadata makes reindexing unsafe.
\par\textit{A strong answer connects mechanism, measurable tradeoffs, failure isolation, observability, and rollback rather than listing product features.}
\paragraph{MCQ-0161 [medium]} Which statement best describes Milvus and Zilliz?
\begin{enumerate}[label=\Alph*.]
\item A PostgreSQL extension that stores vectors beside relational data and supports exact, HNSW, and IVFFlat search.
\item Chroma is a developer-oriented embedding store commonly used for local prototypes and small applications.
\item Milvus is a distributed vector database; Zilliz provides its managed service with multiple index choices.
\item Redis Query Engine adds vector indexes and filters to an in-memory data platform.
\end{enumerate}
\noindent\textbf{Answer.} Milvus is a distributed vector database; Zilliz provides its managed service with multiple index choices.
\par\textit{Milvus is a distributed vector database; Zilliz provides its managed service with multiple index choices. Tradeoff: It targets large-scale separation of compute and storage, with more operational concepts than embedded options. Watch for: Selecting an index by popularity instead of dataset size and latency-recall tests wastes resources.}
\paragraph{MCQ-0162 [hard]} What is the central engineering tradeoff of Milvus and Zilliz?
\begin{enumerate}[label=\Alph*.]
\item It favors simplicity and iteration speed over complex distributed operations.
\item It targets large-scale separation of compute and storage, with more operational concepts than embedded options.
\item Very low latency is possible, but RAM cost and persistence design matter at scale.
\item It simplifies transactions and joins, while specialized stores may scale independent vector workloads more easily.
\end{enumerate}
\noindent\textbf{Answer.} It targets large-scale separation of compute and storage, with more operational concepts than embedded options.
\par\textit{Milvus is a distributed vector database; Zilliz provides its managed service with multiple index choices. Tradeoff: It targets large-scale separation of compute and storage, with more operational concepts than embedded options. Watch for: Selecting an index by popularity instead of dataset size and latency-recall tests wastes resources.}
\paragraph{MCQ-0163 [hard]} Which failure mode should an interviewer expect you to identify for Milvus and Zilliz?
\begin{enumerate}[label=\Alph*.]
\item Promoting a local prototype topology directly to production skips backup, tenancy, and load tests.
\item Adding an ANN index without tuning filtering and planner behavior can reduce recall or miss the index.
\item Selecting an index by popularity instead of dataset size and latency-recall tests wastes resources.
\item Using Redis only as a vector database without capacity planning can evict critical data.
\end{enumerate}
\noindent\textbf{Answer.} Selecting an index by popularity instead of dataset size and latency-recall tests wastes resources.
\par\textit{Milvus is a distributed vector database; Zilliz provides its managed service with multiple index choices. Tradeoff: It targets large-scale separation of compute and storage, with more operational concepts than embedded options. Watch for: Selecting an index by popularity instead of dataset size and latency-recall tests wastes resources.}
\paragraph{MCQ-0164 [medium]} A design review is specifically concerned with this risk: Selecting an index by popularity instead of dataset size and latency-recall tests wastes resources. Which topic should be reviewed first?
\begin{enumerate}[label=\Alph*.]
\item Milvus and Zilliz
\item pgvector
\item Redis vector search
\item Chroma
\end{enumerate}
\noindent\textbf{Answer.} Milvus and Zilliz
\par\textit{Milvus is a distributed vector database; Zilliz provides its managed service with multiple index choices. Tradeoff: It targets large-scale separation of compute and storage, with more operational concepts than embedded options. Watch for: Selecting an index by popularity instead of dataset size and latency-recall tests wastes resources.}
\paragraph{MCQ-0165 [hard]} Which design note most accurately balances the benefits and costs of Milvus and Zilliz?
\begin{enumerate}[label=\Alph*.]
\item Very low latency is possible, but RAM cost and persistence design matter at scale.
\item It simplifies transactions and joins, while specialized stores may scale independent vector workloads more easily.
\item It favors simplicity and iteration speed over complex distributed operations.
\item It targets large-scale separation of compute and storage, with more operational concepts than embedded options.
\end{enumerate}
\noindent\textbf{Answer.} It targets large-scale separation of compute and storage, with more operational concepts than embedded options.
\par\textit{Milvus is a distributed vector database; Zilliz provides its managed service with multiple index choices. Tradeoff: It targets large-scale separation of compute and storage, with more operational concepts than embedded options. Watch for: Selecting an index by popularity instead of dataset size and latency-recall tests wastes resources.}
\paragraph{MCQ-0166 [medium]} Which explanation would be strongest in a system-design interview about Milvus and Zilliz?
\begin{enumerate}[label=\Alph*.]
\item Milvus is a distributed vector database; Zilliz provides its managed service with multiple index choices.
\item Chroma is a developer-oriented embedding store commonly used for local prototypes and small applications.
\item A PostgreSQL extension that stores vectors beside relational data and supports exact, HNSW, and IVFFlat search.
\item Redis Query Engine adds vector indexes and filters to an in-memory data platform.
\end{enumerate}
\noindent\textbf{Answer.} Milvus is a distributed vector database; Zilliz provides its managed service with multiple index choices.
\par\textit{Milvus is a distributed vector database; Zilliz provides its managed service with multiple index choices. Tradeoff: It targets large-scale separation of compute and storage, with more operational concepts than embedded options. Watch for: Selecting an index by popularity instead of dataset size and latency-recall tests wastes resources.}
\paragraph{MCQ-0167 [hard]} Which production warning is most directly associated with Milvus and Zilliz?
\begin{enumerate}[label=\Alph*.]
\item Adding an ANN index without tuning filtering and planner behavior can reduce recall or miss the index.
\item Promoting a local prototype topology directly to production skips backup, tenancy, and load tests.
\item Using Redis only as a vector database without capacity planning can evict critical data.
\item Selecting an index by popularity instead of dataset size and latency-recall tests wastes resources.
\end{enumerate}
\noindent\textbf{Answer.} Selecting an index by popularity instead of dataset size and latency-recall tests wastes resources.
\par\textit{Milvus is a distributed vector database; Zilliz provides its managed service with multiple index choices. Tradeoff: It targets large-scale separation of compute and storage, with more operational concepts than embedded options. Watch for: Selecting an index by popularity instead of dataset size and latency-recall tests wastes resources.}
\paragraph{FC-0168 [medium]} Define Milvus and Zilliz in interview-ready language.
\noindent\textbf{Answer.} Milvus is a distributed vector database; Zilliz provides its managed service with multiple index choices.
\par\textit{Milvus is a distributed vector database; Zilliz provides its managed service with multiple index choices.}
\paragraph{FC-0169 [hard]} State the main tradeoff and production pitfall for Milvus and Zilliz.
\noindent\textbf{Answer.} Tradeoff: It targets large-scale separation of compute and storage, with more operational concepts than embedded options. Pitfall: Selecting an index by popularity instead of dataset size and latency-recall tests wastes resources.
\par\textit{Tradeoff: It targets large-scale separation of compute and storage, with more operational concepts than embedded options. Pitfall: Selecting an index by popularity instead of dataset size and latency-recall tests wastes resources.}
\paragraph{LA-0170 [hard]} Explain Milvus and Zilliz from first principles, then show how you would evaluate and operate it in production.
\noindent\textbf{Answer.} Start with the mechanism: Milvus is a distributed vector database; Zilliz provides its managed service with multiple index choices. Make the tradeoff explicit: It targets large-scale separation of compute and storage, with more operational concepts than embedded options. Define workload-specific offline metrics and a latency/cost budget, test important slices, instrument the critical path, and create a rollback criterion. The production review must explicitly guard against this failure: Selecting an index by popularity instead of dataset size and latency-recall tests wastes resources.
\par\textit{A strong answer connects mechanism, measurable tradeoffs, failure isolation, observability, and rollback rather than listing product features.}
\paragraph{MCQ-0171 [medium]} Which statement best describes Elasticsearch vector search?
\begin{enumerate}[label=\Alph*.]
\item Vespa combines structured filtering, lexical search, tensors, ranking functions, and online serving.
\item An open-source vector engine centered on HNSW, payload filtering, quantization, and distributed collections.
\item Milvus is a distributed vector database; Zilliz provides its managed service with multiple index choices.
\item Elasticsearch combines inverted indexes, filters, aggregations, and approximate kNN in one search platform.
\end{enumerate}
\noindent\textbf{Answer.} Elasticsearch combines inverted indexes, filters, aggregations, and approximate kNN in one search platform.
\par\textit{Elasticsearch combines inverted indexes, filters, aggregations, and approximate kNN in one search platform. Tradeoff: It is compelling for existing Elastic estates but vector memory and shard design require care. Watch for: Oversharding fragments candidate pools and increases coordination overhead.}
\paragraph{MCQ-0172 [hard]} What is the central engineering tradeoff of Elasticsearch vector search?
\begin{enumerate}[label=\Alph*.]
\item It is compelling for existing Elastic estates but vector memory and shard design require care.
\item It offers operational control and rich filtering but adds a stateful distributed service to run.
\item It targets large-scale separation of compute and storage, with more operational concepts than embedded options.
\item It enables sophisticated multi-stage ranking but has a steeper schema and operations learning curve.
\end{enumerate}
\noindent\textbf{Answer.} It is compelling for existing Elastic estates but vector memory and shard design require care.
\par\textit{Elasticsearch combines inverted indexes, filters, aggregations, and approximate kNN in one search platform. Tradeoff: It is compelling for existing Elastic estates but vector memory and shard design require care. Watch for: Oversharding fragments candidate pools and increases coordination overhead.}
\paragraph{MCQ-0173 [hard]} Which failure mode should an interviewer expect you to identify for Elasticsearch vector search?
\begin{enumerate}[label=\Alph*.]
\item Applying post-filtering after ANN can return too few eligible results.
\item Oversharding fragments candidate pools and increases coordination overhead.
\item Putting expensive ranking expressions in the first phase can collapse throughput.
\item Selecting an index by popularity instead of dataset size and latency-recall tests wastes resources.
\end{enumerate}
\noindent\textbf{Answer.} Oversharding fragments candidate pools and increases coordination overhead.
\par\textit{Elasticsearch combines inverted indexes, filters, aggregations, and approximate kNN in one search platform. Tradeoff: It is compelling for existing Elastic estates but vector memory and shard design require care. Watch for: Oversharding fragments candidate pools and increases coordination overhead.}
\paragraph{MCQ-0174 [medium]} A design review is specifically concerned with this risk: Oversharding fragments candidate pools and increases coordination overhead. Which topic should be reviewed first?
\begin{enumerate}[label=\Alph*.]
\item Qdrant
\item Elasticsearch vector search
\item Vespa
\item Milvus and Zilliz
\end{enumerate}
\noindent\textbf{Answer.} Elasticsearch vector search
\par\textit{Elasticsearch combines inverted indexes, filters, aggregations, and approximate kNN in one search platform. Tradeoff: It is compelling for existing Elastic estates but vector memory and shard design require care. Watch for: Oversharding fragments candidate pools and increases coordination overhead.}
\paragraph{MCQ-0175 [hard]} Which design note most accurately balances the benefits and costs of Elasticsearch vector search?
\begin{enumerate}[label=\Alph*.]
\item It offers operational control and rich filtering but adds a stateful distributed service to run.
\item It targets large-scale separation of compute and storage, with more operational concepts than embedded options.
\item It is compelling for existing Elastic estates but vector memory and shard design require care.
\item It enables sophisticated multi-stage ranking but has a steeper schema and operations learning curve.
\end{enumerate}
\noindent\textbf{Answer.} It is compelling for existing Elastic estates but vector memory and shard design require care.
\par\textit{Elasticsearch combines inverted indexes, filters, aggregations, and approximate kNN in one search platform. Tradeoff: It is compelling for existing Elastic estates but vector memory and shard design require care. Watch for: Oversharding fragments candidate pools and increases coordination overhead.}
\paragraph{MCQ-0176 [medium]} Which explanation would be strongest in a system-design interview about Elasticsearch vector search?
\begin{enumerate}[label=\Alph*.]
\item Milvus is a distributed vector database; Zilliz provides its managed service with multiple index choices.
\item An open-source vector engine centered on HNSW, payload filtering, quantization, and distributed collections.
\item Vespa combines structured filtering, lexical search, tensors, ranking functions, and online serving.
\item Elasticsearch combines inverted indexes, filters, aggregations, and approximate kNN in one search platform.
\end{enumerate}
\noindent\textbf{Answer.} Elasticsearch combines inverted indexes, filters, aggregations, and approximate kNN in one search platform.
\par\textit{Elasticsearch combines inverted indexes, filters, aggregations, and approximate kNN in one search platform. Tradeoff: It is compelling for existing Elastic estates but vector memory and shard design require care. Watch for: Oversharding fragments candidate pools and increases coordination overhead.}
\paragraph{MCQ-0177 [hard]} Which production warning is most directly associated with Elasticsearch vector search?
\begin{enumerate}[label=\Alph*.]
\item Applying post-filtering after ANN can return too few eligible results.
\item Oversharding fragments candidate pools and increases coordination overhead.
\item Selecting an index by popularity instead of dataset size and latency-recall tests wastes resources.
\item Putting expensive ranking expressions in the first phase can collapse throughput.
\end{enumerate}
\noindent\textbf{Answer.} Oversharding fragments candidate pools and increases coordination overhead.
\par\textit{Elasticsearch combines inverted indexes, filters, aggregations, and approximate kNN in one search platform. Tradeoff: It is compelling for existing Elastic estates but vector memory and shard design require care. Watch for: Oversharding fragments candidate pools and increases coordination overhead.}
\paragraph{FC-0178 [medium]} Define Elasticsearch vector search in interview-ready language.
\noindent\textbf{Answer.} Elasticsearch combines inverted indexes, filters, aggregations, and approximate kNN in one search platform.
\par\textit{Elasticsearch combines inverted indexes, filters, aggregations, and approximate kNN in one search platform.}
\paragraph{FC-0179 [hard]} State the main tradeoff and production pitfall for Elasticsearch vector search.
\noindent\textbf{Answer.} Tradeoff: It is compelling for existing Elastic estates but vector memory and shard design require care. Pitfall: Oversharding fragments candidate pools and increases coordination overhead.
\par\textit{Tradeoff: It is compelling for existing Elastic estates but vector memory and shard design require care. Pitfall: Oversharding fragments candidate pools and increases coordination overhead.}
\paragraph{LA-0180 [hard]} Explain Elasticsearch vector search from first principles, then show how you would evaluate and operate it in production.
\noindent\textbf{Answer.} Start with the mechanism: Elasticsearch combines inverted indexes, filters, aggregations, and approximate kNN in one search platform. Make the tradeoff explicit: It is compelling for existing Elastic estates but vector memory and shard design require care. Define workload-specific offline metrics and a latency/cost budget, test important slices, instrument the critical path, and create a rollback criterion. The production review must explicitly guard against this failure: Oversharding fragments candidate pools and increases coordination overhead.
\par\textit{A strong answer connects mechanism, measurable tradeoffs, failure isolation, observability, and rollback rather than listing product features.}
\paragraph{MCQ-0181 [medium]} Which statement best describes OpenSearch vector search?
\begin{enumerate}[label=\Alph*.]
\item A navigable multilayer proximity graph controlled by construction degree and search breadth.
\item OpenSearch supports lexical and vector workloads with pluggable kNN engines and search pipelines.
\item Vespa combines structured filtering, lexical search, tensors, ranking functions, and online serving.
\item Elasticsearch combines inverted indexes, filters, aggregations, and approximate kNN in one search platform.
\end{enumerate}
\noindent\textbf{Answer.} OpenSearch supports lexical and vector workloads with pluggable kNN engines and search pipelines.
\par\textit{OpenSearch supports lexical and vector workloads with pluggable kNN engines and search pipelines. Tradeoff: Open governance and integrated search are attractive, while engine-specific behavior complicates tuning. Watch for: Assuming identical parameters across FAISS, Lucene, and NMSLIB backends produces misleading tests.}
\paragraph{MCQ-0182 [hard]} What is the central engineering tradeoff of OpenSearch vector search?
\begin{enumerate}[label=\Alph*.]
\item Open governance and integrated search are attractive, while engine-specific behavior complicates tuning.
\item It offers strong latency-recall behavior with significant memory and build cost.
\item It enables sophisticated multi-stage ranking but has a steeper schema and operations learning curve.
\item It is compelling for existing Elastic estates but vector memory and shard design require care.
\end{enumerate}
\noindent\textbf{Answer.} Open governance and integrated search are attractive, while engine-specific behavior complicates tuning.
\par\textit{OpenSearch supports lexical and vector workloads with pluggable kNN engines and search pipelines. Tradeoff: Open governance and integrated search are attractive, while engine-specific behavior complicates tuning. Watch for: Assuming identical parameters across FAISS, Lucene, and NMSLIB backends produces misleading tests.}
\paragraph{MCQ-0183 [hard]} Which failure mode should an interviewer expect you to identify for OpenSearch vector search?
\begin{enumerate}[label=\Alph*.]
\item Assuming identical parameters across FAISS, Lucene, and NMSLIB backends produces misleading tests.
\item Oversharding fragments candidate pools and increases coordination overhead.
\item Deleting or heavily filtering data without maintenance degrades graph quality.
\item Putting expensive ranking expressions in the first phase can collapse throughput.
\end{enumerate}
\noindent\textbf{Answer.} Assuming identical parameters across FAISS, Lucene, and NMSLIB backends produces misleading tests.
\par\textit{OpenSearch supports lexical and vector workloads with pluggable kNN engines and search pipelines. Tradeoff: Open governance and integrated search are attractive, while engine-specific behavior complicates tuning. Watch for: Assuming identical parameters across FAISS, Lucene, and NMSLIB backends produces misleading tests.}
\paragraph{MCQ-0184 [medium]} A design review is specifically concerned with this risk: Assuming identical parameters across FAISS, Lucene, and NMSLIB backends produces misleading tests. Which topic should be reviewed first?
\begin{enumerate}[label=\Alph*.]
\item HNSW
\item OpenSearch vector search
\item Vespa
\item Elasticsearch vector search
\end{enumerate}
\noindent\textbf{Answer.} OpenSearch vector search
\par\textit{OpenSearch supports lexical and vector workloads with pluggable kNN engines and search pipelines. Tradeoff: Open governance and integrated search are attractive, while engine-specific behavior complicates tuning. Watch for: Assuming identical parameters across FAISS, Lucene, and NMSLIB backends produces misleading tests.}
\paragraph{MCQ-0185 [hard]} Which design note most accurately balances the benefits and costs of OpenSearch vector search?
\begin{enumerate}[label=\Alph*.]
\item Open governance and integrated search are attractive, while engine-specific behavior complicates tuning.
\item It is compelling for existing Elastic estates but vector memory and shard design require care.
\item It offers strong latency-recall behavior with significant memory and build cost.
\item It enables sophisticated multi-stage ranking but has a steeper schema and operations learning curve.
\end{enumerate}
\noindent\textbf{Answer.} Open governance and integrated search are attractive, while engine-specific behavior complicates tuning.
\par\textit{OpenSearch supports lexical and vector workloads with pluggable kNN engines and search pipelines. Tradeoff: Open governance and integrated search are attractive, while engine-specific behavior complicates tuning. Watch for: Assuming identical parameters across FAISS, Lucene, and NMSLIB backends produces misleading tests.}
\paragraph{MCQ-0186 [medium]} Which explanation would be strongest in a system-design interview about OpenSearch vector search?
\begin{enumerate}[label=\Alph*.]
\item OpenSearch supports lexical and vector workloads with pluggable kNN engines and search pipelines.
\item A navigable multilayer proximity graph controlled by construction degree and search breadth.
\item Elasticsearch combines inverted indexes, filters, aggregations, and approximate kNN in one search platform.
\item Vespa combines structured filtering, lexical search, tensors, ranking functions, and online serving.
\end{enumerate}
\noindent\textbf{Answer.} OpenSearch supports lexical and vector workloads with pluggable kNN engines and search pipelines.
\par\textit{OpenSearch supports lexical and vector workloads with pluggable kNN engines and search pipelines. Tradeoff: Open governance and integrated search are attractive, while engine-specific behavior complicates tuning. Watch for: Assuming identical parameters across FAISS, Lucene, and NMSLIB backends produces misleading tests.}
\paragraph{MCQ-0187 [hard]} Which production warning is most directly associated with OpenSearch vector search?
\begin{enumerate}[label=\Alph*.]
\item Putting expensive ranking expressions in the first phase can collapse throughput.
\item Assuming identical parameters across FAISS, Lucene, and NMSLIB backends produces misleading tests.
\item Oversharding fragments candidate pools and increases coordination overhead.
\item Deleting or heavily filtering data without maintenance degrades graph quality.
\end{enumerate}
\noindent\textbf{Answer.} Assuming identical parameters across FAISS, Lucene, and NMSLIB backends produces misleading tests.
\par\textit{OpenSearch supports lexical and vector workloads with pluggable kNN engines and search pipelines. Tradeoff: Open governance and integrated search are attractive, while engine-specific behavior complicates tuning. Watch for: Assuming identical parameters across FAISS, Lucene, and NMSLIB backends produces misleading tests.}
\paragraph{FC-0188 [medium]} Define OpenSearch vector search in interview-ready language.
\noindent\textbf{Answer.} OpenSearch supports lexical and vector workloads with pluggable kNN engines and search pipelines.
\par\textit{OpenSearch supports lexical and vector workloads with pluggable kNN engines and search pipelines.}
\paragraph{FC-0189 [hard]} State the main tradeoff and production pitfall for OpenSearch vector search.
\noindent\textbf{Answer.} Tradeoff: Open governance and integrated search are attractive, while engine-specific behavior complicates tuning. Pitfall: Assuming identical parameters across FAISS, Lucene, and NMSLIB backends produces misleading tests.
\par\textit{Tradeoff: Open governance and integrated search are attractive, while engine-specific behavior complicates tuning. Pitfall: Assuming identical parameters across FAISS, Lucene, and NMSLIB backends produces misleading tests.}
\paragraph{LA-0190 [hard]} Explain OpenSearch vector search from first principles, then show how you would evaluate and operate it in production.
\noindent\textbf{Answer.} Start with the mechanism: OpenSearch supports lexical and vector workloads with pluggable kNN engines and search pipelines. Make the tradeoff explicit: Open governance and integrated search are attractive, while engine-specific behavior complicates tuning. Define workload-specific offline metrics and a latency/cost budget, test important slices, instrument the critical path, and create a rollback criterion. The production review must explicitly guard against this failure: Assuming identical parameters across FAISS, Lucene, and NMSLIB backends produces misleading tests.
\par\textit{A strong answer connects mechanism, measurable tradeoffs, failure isolation, observability, and rollback rather than listing product features.}
\paragraph{MCQ-0191 [medium]} Which statement best describes Redis vector search?
\begin{enumerate}[label=\Alph*.]
\item Redis Query Engine adds vector indexes and filters to an in-memory data platform.
\item Milvus is a distributed vector database; Zilliz provides its managed service with multiple index choices.
\item Attribute constraints restrict eligible vectors before, during, or after ANN search.
\item A vector database with collections, hybrid search, filters, modules, and multi-tenancy features.
\end{enumerate}
\noindent\textbf{Answer.} Redis Query Engine adds vector indexes and filters to an in-memory data platform.
\par\textit{Redis Query Engine adds vector indexes and filters to an in-memory data platform. Tradeoff: Very low latency is possible, but RAM cost and persistence design matter at scale. Watch for: Using Redis only as a vector database without capacity planning can evict critical data.}
\paragraph{MCQ-0192 [hard]} What is the central engineering tradeoff of Redis vector search?
\begin{enumerate}[label=\Alph*.]
\item Very low latency is possible, but RAM cost and persistence design matter at scale.
\item Early filtering preserves relevance but can make graph traversal sparse; post-filtering may underfill results.
\item It targets large-scale separation of compute and storage, with more operational concepts than embedded options.
\item An integrated schema accelerates product work but may couple ingestion and inference choices.
\end{enumerate}
\noindent\textbf{Answer.} Very low latency is possible, but RAM cost and persistence design matter at scale.
\par\textit{Redis Query Engine adds vector indexes and filters to an in-memory data platform. Tradeoff: Very low latency is possible, but RAM cost and persistence design matter at scale. Watch for: Using Redis only as a vector database without capacity planning can evict critical data.}
\paragraph{MCQ-0193 [hard]} Which failure mode should an interviewer expect you to identify for Redis vector search?
\begin{enumerate}[label=\Alph*.]
\item Using Redis only as a vector database without capacity planning can evict critical data.
\item Automatic vectorization without explicit model version metadata makes reindexing unsafe.
\item Failing to test selective filters separately from unfiltered queries hides severe recall loss.
\item Selecting an index by popularity instead of dataset size and latency-recall tests wastes resources.
\end{enumerate}
\noindent\textbf{Answer.} Using Redis only as a vector database without capacity planning can evict critical data.
\par\textit{Redis Query Engine adds vector indexes and filters to an in-memory data platform. Tradeoff: Very low latency is possible, but RAM cost and persistence design matter at scale. Watch for: Using Redis only as a vector database without capacity planning can evict critical data.}
\paragraph{MCQ-0194 [medium]} A design review is specifically concerned with this risk: Using Redis only as a vector database without capacity planning can evict critical data. Which topic should be reviewed first?
\begin{enumerate}[label=\Alph*.]
\item Weaviate
\item Redis vector search
\item Metadata filtering
\item Milvus and Zilliz
\end{enumerate}
\noindent\textbf{Answer.} Redis vector search
\par\textit{Redis Query Engine adds vector indexes and filters to an in-memory data platform. Tradeoff: Very low latency is possible, but RAM cost and persistence design matter at scale. Watch for: Using Redis only as a vector database without capacity planning can evict critical data.}
\paragraph{MCQ-0195 [hard]} Which design note most accurately balances the benefits and costs of Redis vector search?
\begin{enumerate}[label=\Alph*.]
\item Very low latency is possible, but RAM cost and persistence design matter at scale.
\item An integrated schema accelerates product work but may couple ingestion and inference choices.
\item It targets large-scale separation of compute and storage, with more operational concepts than embedded options.
\item Early filtering preserves relevance but can make graph traversal sparse; post-filtering may underfill results.
\end{enumerate}
\noindent\textbf{Answer.} Very low latency is possible, but RAM cost and persistence design matter at scale.
\par\textit{Redis Query Engine adds vector indexes and filters to an in-memory data platform. Tradeoff: Very low latency is possible, but RAM cost and persistence design matter at scale. Watch for: Using Redis only as a vector database without capacity planning can evict critical data.}
\paragraph{MCQ-0196 [medium]} Which explanation would be strongest in a system-design interview about Redis vector search?
\begin{enumerate}[label=\Alph*.]
\item Attribute constraints restrict eligible vectors before, during, or after ANN search.
\item A vector database with collections, hybrid search, filters, modules, and multi-tenancy features.
\item Milvus is a distributed vector database; Zilliz provides its managed service with multiple index choices.
\item Redis Query Engine adds vector indexes and filters to an in-memory data platform.
\end{enumerate}
\noindent\textbf{Answer.} Redis Query Engine adds vector indexes and filters to an in-memory data platform.
\par\textit{Redis Query Engine adds vector indexes and filters to an in-memory data platform. Tradeoff: Very low latency is possible, but RAM cost and persistence design matter at scale. Watch for: Using Redis only as a vector database without capacity planning can evict critical data.}
\paragraph{MCQ-0197 [hard]} Which production warning is most directly associated with Redis vector search?
\begin{enumerate}[label=\Alph*.]
\item Selecting an index by popularity instead of dataset size and latency-recall tests wastes resources.
\item Failing to test selective filters separately from unfiltered queries hides severe recall loss.
\item Using Redis only as a vector database without capacity planning can evict critical data.
\item Automatic vectorization without explicit model version metadata makes reindexing unsafe.
\end{enumerate}
\noindent\textbf{Answer.} Using Redis only as a vector database without capacity planning can evict critical data.
\par\textit{Redis Query Engine adds vector indexes and filters to an in-memory data platform. Tradeoff: Very low latency is possible, but RAM cost and persistence design matter at scale. Watch for: Using Redis only as a vector database without capacity planning can evict critical data.}
\paragraph{FC-0198 [medium]} Define Redis vector search in interview-ready language.
\noindent\textbf{Answer.} Redis Query Engine adds vector indexes and filters to an in-memory data platform.
\par\textit{Redis Query Engine adds vector indexes and filters to an in-memory data platform.}
\paragraph{FC-0199 [hard]} State the main tradeoff and production pitfall for Redis vector search.
\noindent\textbf{Answer.} Tradeoff: Very low latency is possible, but RAM cost and persistence design matter at scale. Pitfall: Using Redis only as a vector database without capacity planning can evict critical data.
\par\textit{Tradeoff: Very low latency is possible, but RAM cost and persistence design matter at scale. Pitfall: Using Redis only as a vector database without capacity planning can evict critical data.}
\paragraph{LA-0200 [hard]} Explain Redis vector search from first principles, then show how you would evaluate and operate it in production.
\noindent\textbf{Answer.} Start with the mechanism: Redis Query Engine adds vector indexes and filters to an in-memory data platform. Make the tradeoff explicit: Very low latency is possible, but RAM cost and persistence design matter at scale. Define workload-specific offline metrics and a latency/cost budget, test important slices, instrument the critical path, and create a rollback criterion. The production review must explicitly guard against this failure: Using Redis only as a vector database without capacity planning can evict critical data.
\par\textit{A strong answer connects mechanism, measurable tradeoffs, failure isolation, observability, and rollback rather than listing product features.}
\paragraph{MCQ-0201 [medium]} Which statement best describes Chroma?
\begin{enumerate}[label=\Alph*.]
\item A PostgreSQL extension that stores vectors beside relational data and supports exact, HNSW, and IVFFlat search.
\item Chroma is a developer-oriented embedding store commonly used for local prototypes and small applications.
\item Attribute constraints restrict eligible vectors before, during, or after ANN search.
\item A vector database with collections, hybrid search, filters, modules, and multi-tenancy features.
\end{enumerate}
\noindent\textbf{Answer.} Chroma is a developer-oriented embedding store commonly used for local prototypes and small applications.
\par\textit{Chroma is a developer-oriented embedding store commonly used for local prototypes and small applications. Tradeoff: It favors simplicity and iteration speed over complex distributed operations. Watch for: Promoting a local prototype topology directly to production skips backup, tenancy, and load tests.}
\paragraph{MCQ-0202 [hard]} What is the central engineering tradeoff of Chroma?
\begin{enumerate}[label=\Alph*.]
\item Early filtering preserves relevance but can make graph traversal sparse; post-filtering may underfill results.
\item It favors simplicity and iteration speed over complex distributed operations.
\item It simplifies transactions and joins, while specialized stores may scale independent vector workloads more easily.
\item An integrated schema accelerates product work but may couple ingestion and inference choices.
\end{enumerate}
\noindent\textbf{Answer.} It favors simplicity and iteration speed over complex distributed operations.
\par\textit{Chroma is a developer-oriented embedding store commonly used for local prototypes and small applications. Tradeoff: It favors simplicity and iteration speed over complex distributed operations. Watch for: Promoting a local prototype topology directly to production skips backup, tenancy, and load tests.}
\paragraph{MCQ-0203 [hard]} Which failure mode should an interviewer expect you to identify for Chroma?
\begin{enumerate}[label=\Alph*.]
\item Promoting a local prototype topology directly to production skips backup, tenancy, and load tests.
\item Failing to test selective filters separately from unfiltered queries hides severe recall loss.
\item Adding an ANN index without tuning filtering and planner behavior can reduce recall or miss the index.
\item Automatic vectorization without explicit model version metadata makes reindexing unsafe.
\end{enumerate}
\noindent\textbf{Answer.} Promoting a local prototype topology directly to production skips backup, tenancy, and load tests.
\par\textit{Chroma is a developer-oriented embedding store commonly used for local prototypes and small applications. Tradeoff: It favors simplicity and iteration speed over complex distributed operations. Watch for: Promoting a local prototype topology directly to production skips backup, tenancy, and load tests.}
\paragraph{MCQ-0204 [medium]} A design review is specifically concerned with this risk: Promoting a local prototype topology directly to production skips backup, tenancy, and load tests. Which topic should be reviewed first?
\begin{enumerate}[label=\Alph*.]
\item Weaviate
\item Metadata filtering
\item Chroma
\item pgvector
\end{enumerate}
\noindent\textbf{Answer.} Chroma
\par\textit{Chroma is a developer-oriented embedding store commonly used for local prototypes and small applications. Tradeoff: It favors simplicity and iteration speed over complex distributed operations. Watch for: Promoting a local prototype topology directly to production skips backup, tenancy, and load tests.}
\paragraph{MCQ-0205 [hard]} Which design note most accurately balances the benefits and costs of Chroma?
\begin{enumerate}[label=\Alph*.]
\item Early filtering preserves relevance but can make graph traversal sparse; post-filtering may underfill results.
\item It simplifies transactions and joins, while specialized stores may scale independent vector workloads more easily.
\item An integrated schema accelerates product work but may couple ingestion and inference choices.
\item It favors simplicity and iteration speed over complex distributed operations.
\end{enumerate}
\noindent\textbf{Answer.} It favors simplicity and iteration speed over complex distributed operations.
\par\textit{Chroma is a developer-oriented embedding store commonly used for local prototypes and small applications. Tradeoff: It favors simplicity and iteration speed over complex distributed operations. Watch for: Promoting a local prototype topology directly to production skips backup, tenancy, and load tests.}
\paragraph{MCQ-0206 [medium]} Which explanation would be strongest in a system-design interview about Chroma?
\begin{enumerate}[label=\Alph*.]
\item Chroma is a developer-oriented embedding store commonly used for local prototypes and small applications.
\item A PostgreSQL extension that stores vectors beside relational data and supports exact, HNSW, and IVFFlat search.
\item Attribute constraints restrict eligible vectors before, during, or after ANN search.
\item A vector database with collections, hybrid search, filters, modules, and multi-tenancy features.
\end{enumerate}
\noindent\textbf{Answer.} Chroma is a developer-oriented embedding store commonly used for local prototypes and small applications.
\par\textit{Chroma is a developer-oriented embedding store commonly used for local prototypes and small applications. Tradeoff: It favors simplicity and iteration speed over complex distributed operations. Watch for: Promoting a local prototype topology directly to production skips backup, tenancy, and load tests.}
\paragraph{MCQ-0207 [hard]} Which production warning is most directly associated with Chroma?
\begin{enumerate}[label=\Alph*.]
\item Automatic vectorization without explicit model version metadata makes reindexing unsafe.
\item Adding an ANN index without tuning filtering and planner behavior can reduce recall or miss the index.
\item Failing to test selective filters separately from unfiltered queries hides severe recall loss.
\item Promoting a local prototype topology directly to production skips backup, tenancy, and load tests.
\end{enumerate}
\noindent\textbf{Answer.} Promoting a local prototype topology directly to production skips backup, tenancy, and load tests.
\par\textit{Chroma is a developer-oriented embedding store commonly used for local prototypes and small applications. Tradeoff: It favors simplicity and iteration speed over complex distributed operations. Watch for: Promoting a local prototype topology directly to production skips backup, tenancy, and load tests.}
\paragraph{FC-0208 [medium]} Define Chroma in interview-ready language.
\noindent\textbf{Answer.} Chroma is a developer-oriented embedding store commonly used for local prototypes and small applications.
\par\textit{Chroma is a developer-oriented embedding store commonly used for local prototypes and small applications.}
\paragraph{FC-0209 [hard]} State the main tradeoff and production pitfall for Chroma.
\noindent\textbf{Answer.} Tradeoff: It favors simplicity and iteration speed over complex distributed operations. Pitfall: Promoting a local prototype topology directly to production skips backup, tenancy, and load tests.
\par\textit{Tradeoff: It favors simplicity and iteration speed over complex distributed operations. Pitfall: Promoting a local prototype topology directly to production skips backup, tenancy, and load tests.}
\paragraph{LA-0210 [hard]} Explain Chroma from first principles, then show how you would evaluate and operate it in production.
\noindent\textbf{Answer.} Start with the mechanism: Chroma is a developer-oriented embedding store commonly used for local prototypes and small applications. Make the tradeoff explicit: It favors simplicity and iteration speed over complex distributed operations. Define workload-specific offline metrics and a latency/cost budget, test important slices, instrument the critical path, and create a rollback criterion. The production review must explicitly guard against this failure: Promoting a local prototype topology directly to production skips backup, tenancy, and load tests.
\par\textit{A strong answer connects mechanism, measurable tradeoffs, failure isolation, observability, and rollback rather than listing product features.}
\paragraph{MCQ-0211 [medium]} Which statement best describes LanceDB?
\begin{enumerate}[label=\Alph*.]
\item LanceDB uses a columnar, versioned data format and supports vector, full-text, and multimodal data locally or in cloud storage.
\item A managed vector service with metadata filtering, namespaces, and dense-sparse hybrid retrieval.
\item Redis Query Engine adds vector indexes and filters to an in-memory data platform.
\item An open-source vector engine centered on HNSW, payload filtering, quantization, and distributed collections.
\end{enumerate}
\noindent\textbf{Answer.} LanceDB uses a columnar, versioned data format and supports vector, full-text, and multimodal data locally or in cloud storage.
\par\textit{LanceDB uses a columnar, versioned data format and supports vector, full-text, and multimodal data locally or in cloud storage. Tradeoff: It fits data-lake and embedded workflows but differs operationally from a server-first database. Watch for: Ignoring compaction and version retention can degrade read performance and inflate storage.}
\paragraph{MCQ-0212 [hard]} What is the central engineering tradeoff of LanceDB?
\begin{enumerate}[label=\Alph*.]
\item It fits data-lake and embedded workflows but differs operationally from a server-first database.
\item Operations are outsourced, trading infrastructure control and portability for managed scale.
\item Very low latency is possible, but RAM cost and persistence design matter at scale.
\item It offers operational control and rich filtering but adds a stateful distributed service to run.
\end{enumerate}
\noindent\textbf{Answer.} It fits data-lake and embedded workflows but differs operationally from a server-first database.
\par\textit{LanceDB uses a columnar, versioned data format and supports vector, full-text, and multimodal data locally or in cloud storage. Tradeoff: It fits data-lake and embedded workflows but differs operationally from a server-first database. Watch for: Ignoring compaction and version retention can degrade read performance and inflate storage.}
\paragraph{MCQ-0213 [hard]} Which failure mode should an interviewer expect you to identify for LanceDB?
\begin{enumerate}[label=\Alph*.]
\item Using Redis only as a vector database without capacity planning can evict critical data.
\item Poor namespace and metadata design creates noisy tenants and expensive fan-out queries.
\item Ignoring compaction and version retention can degrade read performance and inflate storage.
\item Applying post-filtering after ANN can return too few eligible results.
\end{enumerate}
\noindent\textbf{Answer.} Ignoring compaction and version retention can degrade read performance and inflate storage.
\par\textit{LanceDB uses a columnar, versioned data format and supports vector, full-text, and multimodal data locally or in cloud storage. Tradeoff: It fits data-lake and embedded workflows but differs operationally from a server-first database. Watch for: Ignoring compaction and version retention can degrade read performance and inflate storage.}
\paragraph{MCQ-0214 [medium]} A design review is specifically concerned with this risk: Ignoring compaction and version retention can degrade read performance and inflate storage. Which topic should be reviewed first?
\begin{enumerate}[label=\Alph*.]
\item Redis vector search
\item LanceDB
\item Pinecone
\item Qdrant
\end{enumerate}
\noindent\textbf{Answer.} LanceDB
\par\textit{LanceDB uses a columnar, versioned data format and supports vector, full-text, and multimodal data locally or in cloud storage. Tradeoff: It fits data-lake and embedded workflows but differs operationally from a server-first database. Watch for: Ignoring compaction and version retention can degrade read performance and inflate storage.}
\paragraph{MCQ-0215 [hard]} Which design note most accurately balances the benefits and costs of LanceDB?
\begin{enumerate}[label=\Alph*.]
\item It fits data-lake and embedded workflows but differs operationally from a server-first database.
\item It offers operational control and rich filtering but adds a stateful distributed service to run.
\item Operations are outsourced, trading infrastructure control and portability for managed scale.
\item Very low latency is possible, but RAM cost and persistence design matter at scale.
\end{enumerate}
\noindent\textbf{Answer.} It fits data-lake and embedded workflows but differs operationally from a server-first database.
\par\textit{LanceDB uses a columnar, versioned data format and supports vector, full-text, and multimodal data locally or in cloud storage. Tradeoff: It fits data-lake and embedded workflows but differs operationally from a server-first database. Watch for: Ignoring compaction and version retention can degrade read performance and inflate storage.}
\paragraph{MCQ-0216 [medium]} Which explanation would be strongest in a system-design interview about LanceDB?
\begin{enumerate}[label=\Alph*.]
\item LanceDB uses a columnar, versioned data format and supports vector, full-text, and multimodal data locally or in cloud storage.
\item Redis Query Engine adds vector indexes and filters to an in-memory data platform.
\item A managed vector service with metadata filtering, namespaces, and dense-sparse hybrid retrieval.
\item An open-source vector engine centered on HNSW, payload filtering, quantization, and distributed collections.
\end{enumerate}
\noindent\textbf{Answer.} LanceDB uses a columnar, versioned data format and supports vector, full-text, and multimodal data locally or in cloud storage.
\par\textit{LanceDB uses a columnar, versioned data format and supports vector, full-text, and multimodal data locally or in cloud storage. Tradeoff: It fits data-lake and embedded workflows but differs operationally from a server-first database. Watch for: Ignoring compaction and version retention can degrade read performance and inflate storage.}
\paragraph{MCQ-0217 [hard]} Which production warning is most directly associated with LanceDB?
\begin{enumerate}[label=\Alph*.]
\item Poor namespace and metadata design creates noisy tenants and expensive fan-out queries.
\item Applying post-filtering after ANN can return too few eligible results.
\item Using Redis only as a vector database without capacity planning can evict critical data.
\item Ignoring compaction and version retention can degrade read performance and inflate storage.
\end{enumerate}
\noindent\textbf{Answer.} Ignoring compaction and version retention can degrade read performance and inflate storage.
\par\textit{LanceDB uses a columnar, versioned data format and supports vector, full-text, and multimodal data locally or in cloud storage. Tradeoff: It fits data-lake and embedded workflows but differs operationally from a server-first database. Watch for: Ignoring compaction and version retention can degrade read performance and inflate storage.}
\paragraph{FC-0218 [medium]} Define LanceDB in interview-ready language.
\noindent\textbf{Answer.} LanceDB uses a columnar, versioned data format and supports vector, full-text, and multimodal data locally or in cloud storage.
\par\textit{LanceDB uses a columnar, versioned data format and supports vector, full-text, and multimodal data locally or in cloud storage.}
\paragraph{FC-0219 [hard]} State the main tradeoff and production pitfall for LanceDB.
\noindent\textbf{Answer.} Tradeoff: It fits data-lake and embedded workflows but differs operationally from a server-first database. Pitfall: Ignoring compaction and version retention can degrade read performance and inflate storage.
\par\textit{Tradeoff: It fits data-lake and embedded workflows but differs operationally from a server-first database. Pitfall: Ignoring compaction and version retention can degrade read performance and inflate storage.}
\paragraph{LA-0220 [hard]} Explain LanceDB from first principles, then show how you would evaluate and operate it in production.
\noindent\textbf{Answer.} Start with the mechanism: LanceDB uses a columnar, versioned data format and supports vector, full-text, and multimodal data locally or in cloud storage. Make the tradeoff explicit: It fits data-lake and embedded workflows but differs operationally from a server-first database. Define workload-specific offline metrics and a latency/cost budget, test important slices, instrument the critical path, and create a rollback criterion. The production review must explicitly guard against this failure: Ignoring compaction and version retention can degrade read performance and inflate storage.
\par\textit{A strong answer connects mechanism, measurable tradeoffs, failure isolation, observability, and rollback rather than listing product features.}
\paragraph{MCQ-0221 [medium]} Which statement best describes Vespa?
\begin{enumerate}[label=\Alph*.]
\item A vector database with collections, hybrid search, filters, modules, and multi-tenancy features.
\item Vespa combines structured filtering, lexical search, tensors, ranking functions, and online serving.
\item Elasticsearch combines inverted indexes, filters, aggregations, and approximate kNN in one search platform.
\item Milvus is a distributed vector database; Zilliz provides its managed service with multiple index choices.
\end{enumerate}
\noindent\textbf{Answer.} Vespa combines structured filtering, lexical search, tensors, ranking functions, and online serving.
\par\textit{Vespa combines structured filtering, lexical search, tensors, ranking functions, and online serving. Tradeoff: It enables sophisticated multi-stage ranking but has a steeper schema and operations learning curve. Watch for: Putting expensive ranking expressions in the first phase can collapse throughput.}
\paragraph{MCQ-0222 [hard]} What is the central engineering tradeoff of Vespa?
\begin{enumerate}[label=\Alph*.]
\item It is compelling for existing Elastic estates but vector memory and shard design require care.
\item It targets large-scale separation of compute and storage, with more operational concepts than embedded options.
\item It enables sophisticated multi-stage ranking but has a steeper schema and operations learning curve.
\item An integrated schema accelerates product work but may couple ingestion and inference choices.
\end{enumerate}
\noindent\textbf{Answer.} It enables sophisticated multi-stage ranking but has a steeper schema and operations learning curve.
\par\textit{Vespa combines structured filtering, lexical search, tensors, ranking functions, and online serving. Tradeoff: It enables sophisticated multi-stage ranking but has a steeper schema and operations learning curve. Watch for: Putting expensive ranking expressions in the first phase can collapse throughput.}
\paragraph{MCQ-0223 [hard]} Which failure mode should an interviewer expect you to identify for Vespa?
\begin{enumerate}[label=\Alph*.]
\item Automatic vectorization without explicit model version metadata makes reindexing unsafe.
\item Putting expensive ranking expressions in the first phase can collapse throughput.
\item Selecting an index by popularity instead of dataset size and latency-recall tests wastes resources.
\item Oversharding fragments candidate pools and increases coordination overhead.
\end{enumerate}
\noindent\textbf{Answer.} Putting expensive ranking expressions in the first phase can collapse throughput.
\par\textit{Vespa combines structured filtering, lexical search, tensors, ranking functions, and online serving. Tradeoff: It enables sophisticated multi-stage ranking but has a steeper schema and operations learning curve. Watch for: Putting expensive ranking expressions in the first phase can collapse throughput.}
\paragraph{MCQ-0224 [medium]} A design review is specifically concerned with this risk: Putting expensive ranking expressions in the first phase can collapse throughput. Which topic should be reviewed first?
\begin{enumerate}[label=\Alph*.]
\item Milvus and Zilliz
\item Elasticsearch vector search
\item Vespa
\item Weaviate
\end{enumerate}
\noindent\textbf{Answer.} Vespa
\par\textit{Vespa combines structured filtering, lexical search, tensors, ranking functions, and online serving. Tradeoff: It enables sophisticated multi-stage ranking but has a steeper schema and operations learning curve. Watch for: Putting expensive ranking expressions in the first phase can collapse throughput.}
\paragraph{MCQ-0225 [hard]} Which design note most accurately balances the benefits and costs of Vespa?
\begin{enumerate}[label=\Alph*.]
\item It enables sophisticated multi-stage ranking but has a steeper schema and operations learning curve.
\item It targets large-scale separation of compute and storage, with more operational concepts than embedded options.
\item It is compelling for existing Elastic estates but vector memory and shard design require care.
\item An integrated schema accelerates product work but may couple ingestion and inference choices.
\end{enumerate}
\noindent\textbf{Answer.} It enables sophisticated multi-stage ranking but has a steeper schema and operations learning curve.
\par\textit{Vespa combines structured filtering, lexical search, tensors, ranking functions, and online serving. Tradeoff: It enables sophisticated multi-stage ranking but has a steeper schema and operations learning curve. Watch for: Putting expensive ranking expressions in the first phase can collapse throughput.}
\paragraph{MCQ-0226 [medium]} Which explanation would be strongest in a system-design interview about Vespa?
\begin{enumerate}[label=\Alph*.]
\item Elasticsearch combines inverted indexes, filters, aggregations, and approximate kNN in one search platform.
\item Milvus is a distributed vector database; Zilliz provides its managed service with multiple index choices.
\item Vespa combines structured filtering, lexical search, tensors, ranking functions, and online serving.
\item A vector database with collections, hybrid search, filters, modules, and multi-tenancy features.
\end{enumerate}
\noindent\textbf{Answer.} Vespa combines structured filtering, lexical search, tensors, ranking functions, and online serving.
\par\textit{Vespa combines structured filtering, lexical search, tensors, ranking functions, and online serving. Tradeoff: It enables sophisticated multi-stage ranking but has a steeper schema and operations learning curve. Watch for: Putting expensive ranking expressions in the first phase can collapse throughput.}
\paragraph{MCQ-0227 [hard]} Which production warning is most directly associated with Vespa?
\begin{enumerate}[label=\Alph*.]
\item Selecting an index by popularity instead of dataset size and latency-recall tests wastes resources.
\item Putting expensive ranking expressions in the first phase can collapse throughput.
\item Oversharding fragments candidate pools and increases coordination overhead.
\item Automatic vectorization without explicit model version metadata makes reindexing unsafe.
\end{enumerate}
\noindent\textbf{Answer.} Putting expensive ranking expressions in the first phase can collapse throughput.
\par\textit{Vespa combines structured filtering, lexical search, tensors, ranking functions, and online serving. Tradeoff: It enables sophisticated multi-stage ranking but has a steeper schema and operations learning curve. Watch for: Putting expensive ranking expressions in the first phase can collapse throughput.}
\paragraph{FC-0228 [medium]} Define Vespa in interview-ready language.
\noindent\textbf{Answer.} Vespa combines structured filtering, lexical search, tensors, ranking functions, and online serving.
\par\textit{Vespa combines structured filtering, lexical search, tensors, ranking functions, and online serving.}
\paragraph{FC-0229 [hard]} State the main tradeoff and production pitfall for Vespa.
\noindent\textbf{Answer.} Tradeoff: It enables sophisticated multi-stage ranking but has a steeper schema and operations learning curve. Pitfall: Putting expensive ranking expressions in the first phase can collapse throughput.
\par\textit{Tradeoff: It enables sophisticated multi-stage ranking but has a steeper schema and operations learning curve. Pitfall: Putting expensive ranking expressions in the first phase can collapse throughput.}
\paragraph{LA-0230 [hard]} Explain Vespa from first principles, then show how you would evaluate and operate it in production.
\noindent\textbf{Answer.} Start with the mechanism: Vespa combines structured filtering, lexical search, tensors, ranking functions, and online serving. Make the tradeoff explicit: It enables sophisticated multi-stage ranking but has a steeper schema and operations learning curve. Define workload-specific offline metrics and a latency/cost budget, test important slices, instrument the critical path, and create a rollback criterion. The production review must explicitly guard against this failure: Putting expensive ranking expressions in the first phase can collapse throughput.
\par\textit{A strong answer connects mechanism, measurable tradeoffs, failure isolation, observability, and rollback rather than listing product features.}
\paragraph{MCQ-0231 [medium]} Which statement best describes HNSW?
\begin{enumerate}[label=\Alph*.]
\item LanceDB uses a columnar, versioned data format and supports vector, full-text, and multimodal data locally or in cloud storage.
\item Attribute constraints restrict eligible vectors before, during, or after ANN search.
\item A navigable multilayer proximity graph controlled by construction degree and search breadth.
\item OpenSearch supports lexical and vector workloads with pluggable kNN engines and search pipelines.
\end{enumerate}
\noindent\textbf{Answer.} A navigable multilayer proximity graph controlled by construction degree and search breadth.
\par\textit{A navigable multilayer proximity graph controlled by construction degree and search breadth. Tradeoff: It offers strong latency-recall behavior with significant memory and build cost. Watch for: Deleting or heavily filtering data without maintenance degrades graph quality.}
\paragraph{MCQ-0232 [hard]} What is the central engineering tradeoff of HNSW?
\begin{enumerate}[label=\Alph*.]
\item It offers strong latency-recall behavior with significant memory and build cost.
\item Open governance and integrated search are attractive, while engine-specific behavior complicates tuning.
\item Early filtering preserves relevance but can make graph traversal sparse; post-filtering may underfill results.
\item It fits data-lake and embedded workflows but differs operationally from a server-first database.
\end{enumerate}
\noindent\textbf{Answer.} It offers strong latency-recall behavior with significant memory and build cost.
\par\textit{A navigable multilayer proximity graph controlled by construction degree and search breadth. Tradeoff: It offers strong latency-recall behavior with significant memory and build cost. Watch for: Deleting or heavily filtering data without maintenance degrades graph quality.}
\paragraph{MCQ-0233 [hard]} Which failure mode should an interviewer expect you to identify for HNSW?
\begin{enumerate}[label=\Alph*.]
\item Ignoring compaction and version retention can degrade read performance and inflate storage.
\item Deleting or heavily filtering data without maintenance degrades graph quality.
\item Assuming identical parameters across FAISS, Lucene, and NMSLIB backends produces misleading tests.
\item Failing to test selective filters separately from unfiltered queries hides severe recall loss.
\end{enumerate}
\noindent\textbf{Answer.} Deleting or heavily filtering data without maintenance degrades graph quality.
\par\textit{A navigable multilayer proximity graph controlled by construction degree and search breadth. Tradeoff: It offers strong latency-recall behavior with significant memory and build cost. Watch for: Deleting or heavily filtering data without maintenance degrades graph quality.}
\paragraph{MCQ-0234 [medium]} A design review is specifically concerned with this risk: Deleting or heavily filtering data without maintenance degrades graph quality. Which topic should be reviewed first?
\begin{enumerate}[label=\Alph*.]
\item HNSW
\item OpenSearch vector search
\item Metadata filtering
\item LanceDB
\end{enumerate}
\noindent\textbf{Answer.} HNSW
\par\textit{A navigable multilayer proximity graph controlled by construction degree and search breadth. Tradeoff: It offers strong latency-recall behavior with significant memory and build cost. Watch for: Deleting or heavily filtering data without maintenance degrades graph quality.}
\paragraph{MCQ-0235 [hard]} Which design note most accurately balances the benefits and costs of HNSW?
\begin{enumerate}[label=\Alph*.]
\item Early filtering preserves relevance but can make graph traversal sparse; post-filtering may underfill results.
\item It fits data-lake and embedded workflows but differs operationally from a server-first database.
\item Open governance and integrated search are attractive, while engine-specific behavior complicates tuning.
\item It offers strong latency-recall behavior with significant memory and build cost.
\end{enumerate}
\noindent\textbf{Answer.} It offers strong latency-recall behavior with significant memory and build cost.
\par\textit{A navigable multilayer proximity graph controlled by construction degree and search breadth. Tradeoff: It offers strong latency-recall behavior with significant memory and build cost. Watch for: Deleting or heavily filtering data without maintenance degrades graph quality.}
\paragraph{MCQ-0236 [medium]} Which explanation would be strongest in a system-design interview about HNSW?
\begin{enumerate}[label=\Alph*.]
\item OpenSearch supports lexical and vector workloads with pluggable kNN engines and search pipelines.
\item A navigable multilayer proximity graph controlled by construction degree and search breadth.
\item LanceDB uses a columnar, versioned data format and supports vector, full-text, and multimodal data locally or in cloud storage.
\item Attribute constraints restrict eligible vectors before, during, or after ANN search.
\end{enumerate}
\noindent\textbf{Answer.} A navigable multilayer proximity graph controlled by construction degree and search breadth.
\par\textit{A navigable multilayer proximity graph controlled by construction degree and search breadth. Tradeoff: It offers strong latency-recall behavior with significant memory and build cost. Watch for: Deleting or heavily filtering data without maintenance degrades graph quality.}
\paragraph{MCQ-0237 [hard]} Which production warning is most directly associated with HNSW?
\begin{enumerate}[label=\Alph*.]
\item Assuming identical parameters across FAISS, Lucene, and NMSLIB backends produces misleading tests.
\item Ignoring compaction and version retention can degrade read performance and inflate storage.
\item Deleting or heavily filtering data without maintenance degrades graph quality.
\item Failing to test selective filters separately from unfiltered queries hides severe recall loss.
\end{enumerate}
\noindent\textbf{Answer.} Deleting or heavily filtering data without maintenance degrades graph quality.
\par\textit{A navigable multilayer proximity graph controlled by construction degree and search breadth. Tradeoff: It offers strong latency-recall behavior with significant memory and build cost. Watch for: Deleting or heavily filtering data without maintenance degrades graph quality.}
\paragraph{FC-0238 [medium]} Define HNSW in interview-ready language.
\noindent\textbf{Answer.} A navigable multilayer proximity graph controlled by construction degree and search breadth.
\par\textit{A navigable multilayer proximity graph controlled by construction degree and search breadth.}
\paragraph{FC-0239 [hard]} State the main tradeoff and production pitfall for HNSW.
\noindent\textbf{Answer.} Tradeoff: It offers strong latency-recall behavior with significant memory and build cost. Pitfall: Deleting or heavily filtering data without maintenance degrades graph quality.
\par\textit{Tradeoff: It offers strong latency-recall behavior with significant memory and build cost. Pitfall: Deleting or heavily filtering data without maintenance degrades graph quality.}
\paragraph{LA-0240 [hard]} Explain HNSW from first principles, then show how you would evaluate and operate it in production.
\noindent\textbf{Answer.} Start with the mechanism: A navigable multilayer proximity graph controlled by construction degree and search breadth. Make the tradeoff explicit: It offers strong latency-recall behavior with significant memory and build cost. Define workload-specific offline metrics and a latency/cost budget, test important slices, instrument the critical path, and create a rollback criterion. The production review must explicitly guard against this failure: Deleting or heavily filtering data without maintenance degrades graph quality.
\par\textit{A strong answer connects mechanism, measurable tradeoffs, failure isolation, observability, and rollback rather than listing product features.}
\paragraph{MCQ-0241 [medium]} Which statement best describes IVF and PQ?
\begin{enumerate}[label=\Alph*.]
\item LanceDB uses a columnar, versioned data format and supports vector, full-text, and multimodal data locally or in cloud storage.
\item OpenSearch supports lexical and vector workloads with pluggable kNN engines and search pipelines.
\item A navigable multilayer proximity graph controlled by construction degree and search breadth.
\item Inverted files narrow the search to clusters; product quantization compresses vectors for cheaper approximate scoring.
\end{enumerate}
\noindent\textbf{Answer.} Inverted files narrow the search to clusters; product quantization compresses vectors for cheaper approximate scoring.
\par\textit{Inverted files narrow the search to clusters; product quantization compresses vectors for cheaper approximate scoring. Tradeoff: Memory and latency fall, while training quality and quantization error reduce recall. Watch for: Training centroids or codebooks on an unrepresentative sample bakes bias into the index.}
\paragraph{MCQ-0242 [hard]} What is the central engineering tradeoff of IVF and PQ?
\begin{enumerate}[label=\Alph*.]
\item It fits data-lake and embedded workflows but differs operationally from a server-first database.
\item Memory and latency fall, while training quality and quantization error reduce recall.
\item It offers strong latency-recall behavior with significant memory and build cost.
\item Open governance and integrated search are attractive, while engine-specific behavior complicates tuning.
\end{enumerate}
\noindent\textbf{Answer.} Memory and latency fall, while training quality and quantization error reduce recall.
\par\textit{Inverted files narrow the search to clusters; product quantization compresses vectors for cheaper approximate scoring. Tradeoff: Memory and latency fall, while training quality and quantization error reduce recall. Watch for: Training centroids or codebooks on an unrepresentative sample bakes bias into the index.}
\paragraph{MCQ-0243 [hard]} Which failure mode should an interviewer expect you to identify for IVF and PQ?
\begin{enumerate}[label=\Alph*.]
\item Deleting or heavily filtering data without maintenance degrades graph quality.
\item Ignoring compaction and version retention can degrade read performance and inflate storage.
\item Assuming identical parameters across FAISS, Lucene, and NMSLIB backends produces misleading tests.
\item Training centroids or codebooks on an unrepresentative sample bakes bias into the index.
\end{enumerate}
\noindent\textbf{Answer.} Training centroids or codebooks on an unrepresentative sample bakes bias into the index.
\par\textit{Inverted files narrow the search to clusters; product quantization compresses vectors for cheaper approximate scoring. Tradeoff: Memory and latency fall, while training quality and quantization error reduce recall. Watch for: Training centroids or codebooks on an unrepresentative sample bakes bias into the index.}
\paragraph{MCQ-0244 [medium]} A design review is specifically concerned with this risk: Training centroids or codebooks on an unrepresentative sample bakes bias into the index. Which topic should be reviewed first?
\begin{enumerate}[label=\Alph*.]
\item HNSW
\item OpenSearch vector search
\item IVF and PQ
\item LanceDB
\end{enumerate}
\noindent\textbf{Answer.} IVF and PQ
\par\textit{Inverted files narrow the search to clusters; product quantization compresses vectors for cheaper approximate scoring. Tradeoff: Memory and latency fall, while training quality and quantization error reduce recall. Watch for: Training centroids or codebooks on an unrepresentative sample bakes bias into the index.}
\paragraph{MCQ-0245 [hard]} Which design note most accurately balances the benefits and costs of IVF and PQ?
\begin{enumerate}[label=\Alph*.]
\item Memory and latency fall, while training quality and quantization error reduce recall.
\item Open governance and integrated search are attractive, while engine-specific behavior complicates tuning.
\item It offers strong latency-recall behavior with significant memory and build cost.
\item It fits data-lake and embedded workflows but differs operationally from a server-first database.
\end{enumerate}
\noindent\textbf{Answer.} Memory and latency fall, while training quality and quantization error reduce recall.
\par\textit{Inverted files narrow the search to clusters; product quantization compresses vectors for cheaper approximate scoring. Tradeoff: Memory and latency fall, while training quality and quantization error reduce recall. Watch for: Training centroids or codebooks on an unrepresentative sample bakes bias into the index.}
\paragraph{MCQ-0246 [medium]} Which explanation would be strongest in a system-design interview about IVF and PQ?
\begin{enumerate}[label=\Alph*.]
\item Inverted files narrow the search to clusters; product quantization compresses vectors for cheaper approximate scoring.
\item OpenSearch supports lexical and vector workloads with pluggable kNN engines and search pipelines.
\item A navigable multilayer proximity graph controlled by construction degree and search breadth.
\item LanceDB uses a columnar, versioned data format and supports vector, full-text, and multimodal data locally or in cloud storage.
\end{enumerate}
\noindent\textbf{Answer.} Inverted files narrow the search to clusters; product quantization compresses vectors for cheaper approximate scoring.
\par\textit{Inverted files narrow the search to clusters; product quantization compresses vectors for cheaper approximate scoring. Tradeoff: Memory and latency fall, while training quality and quantization error reduce recall. Watch for: Training centroids or codebooks on an unrepresentative sample bakes bias into the index.}
\paragraph{MCQ-0247 [hard]} Which production warning is most directly associated with IVF and PQ?
\begin{enumerate}[label=\Alph*.]
\item Ignoring compaction and version retention can degrade read performance and inflate storage.
\item Assuming identical parameters across FAISS, Lucene, and NMSLIB backends produces misleading tests.
\item Deleting or heavily filtering data without maintenance degrades graph quality.
\item Training centroids or codebooks on an unrepresentative sample bakes bias into the index.
\end{enumerate}
\noindent\textbf{Answer.} Training centroids or codebooks on an unrepresentative sample bakes bias into the index.
\par\textit{Inverted files narrow the search to clusters; product quantization compresses vectors for cheaper approximate scoring. Tradeoff: Memory and latency fall, while training quality and quantization error reduce recall. Watch for: Training centroids or codebooks on an unrepresentative sample bakes bias into the index.}
\paragraph{FC-0248 [medium]} Define IVF and PQ in interview-ready language.
\noindent\textbf{Answer.} Inverted files narrow the search to clusters; product quantization compresses vectors for cheaper approximate scoring.
\par\textit{Inverted files narrow the search to clusters; product quantization compresses vectors for cheaper approximate scoring.}
\paragraph{FC-0249 [hard]} State the main tradeoff and production pitfall for IVF and PQ.
\noindent\textbf{Answer.} Tradeoff: Memory and latency fall, while training quality and quantization error reduce recall. Pitfall: Training centroids or codebooks on an unrepresentative sample bakes bias into the index.
\par\textit{Tradeoff: Memory and latency fall, while training quality and quantization error reduce recall. Pitfall: Training centroids or codebooks on an unrepresentative sample bakes bias into the index.}
\paragraph{LA-0250 [hard]} Explain IVF and PQ from first principles, then show how you would evaluate and operate it in production.
\noindent\textbf{Answer.} Start with the mechanism: Inverted files narrow the search to clusters; product quantization compresses vectors for cheaper approximate scoring. Make the tradeoff explicit: Memory and latency fall, while training quality and quantization error reduce recall. Define workload-specific offline metrics and a latency/cost budget, test important slices, instrument the critical path, and create a rollback criterion. The production review must explicitly guard against this failure: Training centroids or codebooks on an unrepresentative sample bakes bias into the index.
\par\textit{A strong answer connects mechanism, measurable tradeoffs, failure isolation, observability, and rollback rather than listing product features.}
\paragraph{MCQ-0251 [medium]} Which statement best describes Metadata filtering?
\begin{enumerate}[label=\Alph*.]
\item A navigable multilayer proximity graph controlled by construction degree and search breadth.
\item A PostgreSQL extension that stores vectors beside relational data and supports exact, HNSW, and IVFFlat search.
\item Attribute constraints restrict eligible vectors before, during, or after ANN search.
\item Chroma is a developer-oriented embedding store commonly used for local prototypes and small applications.
\end{enumerate}
\noindent\textbf{Answer.} Attribute constraints restrict eligible vectors before, during, or after ANN search.
\par\textit{Attribute constraints restrict eligible vectors before, during, or after ANN search. Tradeoff: Early filtering preserves relevance but can make graph traversal sparse; post-filtering may underfill results. Watch for: Failing to test selective filters separately from unfiltered queries hides severe recall loss.}
\paragraph{MCQ-0252 [hard]} What is the central engineering tradeoff of Metadata filtering?
\begin{enumerate}[label=\Alph*.]
\item It offers strong latency-recall behavior with significant memory and build cost.
\item It favors simplicity and iteration speed over complex distributed operations.
\item It simplifies transactions and joins, while specialized stores may scale independent vector workloads more easily.
\item Early filtering preserves relevance but can make graph traversal sparse; post-filtering may underfill results.
\end{enumerate}
\noindent\textbf{Answer.} Early filtering preserves relevance but can make graph traversal sparse; post-filtering may underfill results.
\par\textit{Attribute constraints restrict eligible vectors before, during, or after ANN search. Tradeoff: Early filtering preserves relevance but can make graph traversal sparse; post-filtering may underfill results. Watch for: Failing to test selective filters separately from unfiltered queries hides severe recall loss.}
\paragraph{MCQ-0253 [hard]} Which failure mode should an interviewer expect you to identify for Metadata filtering?
\begin{enumerate}[label=\Alph*.]
\item Adding an ANN index without tuning filtering and planner behavior can reduce recall or miss the index.
\item Promoting a local prototype topology directly to production skips backup, tenancy, and load tests.
\item Failing to test selective filters separately from unfiltered queries hides severe recall loss.
\item Deleting or heavily filtering data without maintenance degrades graph quality.
\end{enumerate}
\noindent\textbf{Answer.} Failing to test selective filters separately from unfiltered queries hides severe recall loss.
\par\textit{Attribute constraints restrict eligible vectors before, during, or after ANN search. Tradeoff: Early filtering preserves relevance but can make graph traversal sparse; post-filtering may underfill results. Watch for: Failing to test selective filters separately from unfiltered queries hides severe recall loss.}
\paragraph{MCQ-0254 [medium]} A design review is specifically concerned with this risk: Failing to test selective filters separately from unfiltered queries hides severe recall loss. Which topic should be reviewed first?
\begin{enumerate}[label=\Alph*.]
\item Chroma
\item pgvector
\item Metadata filtering
\item HNSW
\end{enumerate}
\noindent\textbf{Answer.} Metadata filtering
\par\textit{Attribute constraints restrict eligible vectors before, during, or after ANN search. Tradeoff: Early filtering preserves relevance but can make graph traversal sparse; post-filtering may underfill results. Watch for: Failing to test selective filters separately from unfiltered queries hides severe recall loss.}
\paragraph{MCQ-0255 [hard]} Which design note most accurately balances the benefits and costs of Metadata filtering?
\begin{enumerate}[label=\Alph*.]
\item It simplifies transactions and joins, while specialized stores may scale independent vector workloads more easily.
\item It offers strong latency-recall behavior with significant memory and build cost.
\item It favors simplicity and iteration speed over complex distributed operations.
\item Early filtering preserves relevance but can make graph traversal sparse; post-filtering may underfill results.
\end{enumerate}
\noindent\textbf{Answer.} Early filtering preserves relevance but can make graph traversal sparse; post-filtering may underfill results.
\par\textit{Attribute constraints restrict eligible vectors before, during, or after ANN search. Tradeoff: Early filtering preserves relevance but can make graph traversal sparse; post-filtering may underfill results. Watch for: Failing to test selective filters separately from unfiltered queries hides severe recall loss.}
\paragraph{MCQ-0256 [medium]} Which explanation would be strongest in a system-design interview about Metadata filtering?
\begin{enumerate}[label=\Alph*.]
\item Attribute constraints restrict eligible vectors before, during, or after ANN search.
\item A PostgreSQL extension that stores vectors beside relational data and supports exact, HNSW, and IVFFlat search.
\item Chroma is a developer-oriented embedding store commonly used for local prototypes and small applications.
\item A navigable multilayer proximity graph controlled by construction degree and search breadth.
\end{enumerate}
\noindent\textbf{Answer.} Attribute constraints restrict eligible vectors before, during, or after ANN search.
\par\textit{Attribute constraints restrict eligible vectors before, during, or after ANN search. Tradeoff: Early filtering preserves relevance but can make graph traversal sparse; post-filtering may underfill results. Watch for: Failing to test selective filters separately from unfiltered queries hides severe recall loss.}
\paragraph{MCQ-0257 [hard]} Which production warning is most directly associated with Metadata filtering?
\begin{enumerate}[label=\Alph*.]
\item Adding an ANN index without tuning filtering and planner behavior can reduce recall or miss the index.
\item Deleting or heavily filtering data without maintenance degrades graph quality.
\item Failing to test selective filters separately from unfiltered queries hides severe recall loss.
\item Promoting a local prototype topology directly to production skips backup, tenancy, and load tests.
\end{enumerate}
\noindent\textbf{Answer.} Failing to test selective filters separately from unfiltered queries hides severe recall loss.
\par\textit{Attribute constraints restrict eligible vectors before, during, or after ANN search. Tradeoff: Early filtering preserves relevance but can make graph traversal sparse; post-filtering may underfill results. Watch for: Failing to test selective filters separately from unfiltered queries hides severe recall loss.}
\paragraph{FC-0258 [medium]} Define Metadata filtering in interview-ready language.
\noindent\textbf{Answer.} Attribute constraints restrict eligible vectors before, during, or after ANN search.
\par\textit{Attribute constraints restrict eligible vectors before, during, or after ANN search.}
\paragraph{FC-0259 [hard]} State the main tradeoff and production pitfall for Metadata filtering.
\noindent\textbf{Answer.} Tradeoff: Early filtering preserves relevance but can make graph traversal sparse; post-filtering may underfill results. Pitfall: Failing to test selective filters separately from unfiltered queries hides severe recall loss.
\par\textit{Tradeoff: Early filtering preserves relevance but can make graph traversal sparse; post-filtering may underfill results. Pitfall: Failing to test selective filters separately from unfiltered queries hides severe recall loss.}
\paragraph{LA-0260 [hard]} Explain Metadata filtering from first principles, then show how you would evaluate and operate it in production.
\noindent\textbf{Answer.} Start with the mechanism: Attribute constraints restrict eligible vectors before, during, or after ANN search. Make the tradeoff explicit: Early filtering preserves relevance but can make graph traversal sparse; post-filtering may underfill results. Define workload-specific offline metrics and a latency/cost budget, test important slices, instrument the critical path, and create a rollback criterion. The production review must explicitly guard against this failure: Failing to test selective filters separately from unfiltered queries hides severe recall loss.
\par\textit{A strong answer connects mechanism, measurable tradeoffs, failure isolation, observability, and rollback rather than listing product features.}
\subsection{RAG Architecture and Evaluation}
\paragraph{MCQ-0261 [medium]} Which statement best describes Document ingestion?
\begin{enumerate}[label=\Alph*.]
\item Models generate questions, expected evidence, and variants from a source corpus to expand evaluation coverage.
\item The system generates several query formulations and merges their result sets.
\item A post-retrieval step extracts or rewrites only query-relevant evidence from candidates.
\item Ingestion parses, normalizes, classifies, deduplicates, and records provenance before indexing.
\end{enumerate}
\noindent\textbf{Answer.} Ingestion parses, normalizes, classifies, deduplicates, and records provenance before indexing.
\par\textit{Ingestion parses, normalizes, classifies, deduplicates, and records provenance before indexing. Tradeoff: Rich preprocessing improves retrieval but increases pipeline cost and reprocessing complexity. Watch for: Losing stable document and chunk identifiers makes updates create duplicates.}
\paragraph{MCQ-0262 [hard]} What is the central engineering tradeoff of Document ingestion?
\begin{enumerate}[label=\Alph*.]
\item Rich preprocessing improves retrieval but increases pipeline cost and reprocessing complexity.
\item Coverage grows cheaply, while generator bias and leakage can make the benchmark too easy.
\item Recall rises at the cost of additional searches, duplicates, and possible intent drift.
\item It saves context tokens but can delete qualifiers or introduce summarization errors.
\end{enumerate}
\noindent\textbf{Answer.} Rich preprocessing improves retrieval but increases pipeline cost and reprocessing complexity.
\par\textit{Ingestion parses, normalizes, classifies, deduplicates, and records provenance before indexing. Tradeoff: Rich preprocessing improves retrieval but increases pipeline cost and reprocessing complexity. Watch for: Losing stable document and chunk identifiers makes updates create duplicates.}
\paragraph{MCQ-0263 [hard]} Which failure mode should an interviewer expect you to identify for Document ingestion?
\begin{enumerate}[label=\Alph*.]
\item Testing with the same model family that authored the cases can overstate quality.
\item Losing stable document and chunk identifiers makes updates create duplicates.
\item Using the same weak model for compression and answering creates correlated failures.
\item Treating all generated queries as equally trustworthy invites off-topic evidence.
\end{enumerate}
\noindent\textbf{Answer.} Losing stable document and chunk identifiers makes updates create duplicates.
\par\textit{Ingestion parses, normalizes, classifies, deduplicates, and records provenance before indexing. Tradeoff: Rich preprocessing improves retrieval but increases pipeline cost and reprocessing complexity. Watch for: Losing stable document and chunk identifiers makes updates create duplicates.}
\paragraph{MCQ-0264 [medium]} A design review is specifically concerned with this risk: Losing stable document and chunk identifiers makes updates create duplicates. Which topic should be reviewed first?
\begin{enumerate}[label=\Alph*.]
\item Synthetic test generation
\item Multi-query retrieval
\item Contextual compression
\item Document ingestion
\end{enumerate}
\noindent\textbf{Answer.} Document ingestion
\par\textit{Ingestion parses, normalizes, classifies, deduplicates, and records provenance before indexing. Tradeoff: Rich preprocessing improves retrieval but increases pipeline cost and reprocessing complexity. Watch for: Losing stable document and chunk identifiers makes updates create duplicates.}
\paragraph{MCQ-0265 [hard]} Which design note most accurately balances the benefits and costs of Document ingestion?
\begin{enumerate}[label=\Alph*.]
\item Coverage grows cheaply, while generator bias and leakage can make the benchmark too easy.
\item Recall rises at the cost of additional searches, duplicates, and possible intent drift.
\item It saves context tokens but can delete qualifiers or introduce summarization errors.
\item Rich preprocessing improves retrieval but increases pipeline cost and reprocessing complexity.
\end{enumerate}
\noindent\textbf{Answer.} Rich preprocessing improves retrieval but increases pipeline cost and reprocessing complexity.
\par\textit{Ingestion parses, normalizes, classifies, deduplicates, and records provenance before indexing. Tradeoff: Rich preprocessing improves retrieval but increases pipeline cost and reprocessing complexity. Watch for: Losing stable document and chunk identifiers makes updates create duplicates.}
\paragraph{MCQ-0266 [medium]} Which explanation would be strongest in a system-design interview about Document ingestion?
\begin{enumerate}[label=\Alph*.]
\item The system generates several query formulations and merges their result sets.
\item Models generate questions, expected evidence, and variants from a source corpus to expand evaluation coverage.
\item A post-retrieval step extracts or rewrites only query-relevant evidence from candidates.
\item Ingestion parses, normalizes, classifies, deduplicates, and records provenance before indexing.
\end{enumerate}
\noindent\textbf{Answer.} Ingestion parses, normalizes, classifies, deduplicates, and records provenance before indexing.
\par\textit{Ingestion parses, normalizes, classifies, deduplicates, and records provenance before indexing. Tradeoff: Rich preprocessing improves retrieval but increases pipeline cost and reprocessing complexity. Watch for: Losing stable document and chunk identifiers makes updates create duplicates.}
\paragraph{MCQ-0267 [hard]} Which production warning is most directly associated with Document ingestion?
\begin{enumerate}[label=\Alph*.]
\item Losing stable document and chunk identifiers makes updates create duplicates.
\item Treating all generated queries as equally trustworthy invites off-topic evidence.
\item Testing with the same model family that authored the cases can overstate quality.
\item Using the same weak model for compression and answering creates correlated failures.
\end{enumerate}
\noindent\textbf{Answer.} Losing stable document and chunk identifiers makes updates create duplicates.
\par\textit{Ingestion parses, normalizes, classifies, deduplicates, and records provenance before indexing. Tradeoff: Rich preprocessing improves retrieval but increases pipeline cost and reprocessing complexity. Watch for: Losing stable document and chunk identifiers makes updates create duplicates.}
\paragraph{FC-0268 [medium]} Define Document ingestion in interview-ready language.
\noindent\textbf{Answer.} Ingestion parses, normalizes, classifies, deduplicates, and records provenance before indexing.
\par\textit{Ingestion parses, normalizes, classifies, deduplicates, and records provenance before indexing.}
\paragraph{FC-0269 [hard]} State the main tradeoff and production pitfall for Document ingestion.
\noindent\textbf{Answer.} Tradeoff: Rich preprocessing improves retrieval but increases pipeline cost and reprocessing complexity. Pitfall: Losing stable document and chunk identifiers makes updates create duplicates.
\par\textit{Tradeoff: Rich preprocessing improves retrieval but increases pipeline cost and reprocessing complexity. Pitfall: Losing stable document and chunk identifiers makes updates create duplicates.}
\paragraph{LA-0270 [hard]} Explain Document ingestion from first principles, then show how you would evaluate and operate it in production.
\noindent\textbf{Answer.} Start with the mechanism: Ingestion parses, normalizes, classifies, deduplicates, and records provenance before indexing. Make the tradeoff explicit: Rich preprocessing improves retrieval but increases pipeline cost and reprocessing complexity. Define workload-specific offline metrics and a latency/cost budget, test important slices, instrument the critical path, and create a rollback criterion. The production review must explicitly guard against this failure: Losing stable document and chunk identifiers makes updates create duplicates.
\par\textit{A strong answer connects mechanism, measurable tradeoffs, failure isolation, observability, and rollback rather than listing product features.}
\paragraph{MCQ-0271 [medium]} Which statement best describes Chunking?
\begin{enumerate}[label=\Alph*.]
\item Small child chunks retrieve precisely while larger parent sections are supplied to the model.
\item Faithfulness asks whether claims in an answer are supported by the supplied context.
\item Chunking chooses semantic units, size, overlap, and metadata boundaries for retrieval and context assembly.
\item Hypothetical Document Embeddings generate a likely answer passage and embed it as the retrieval query.
\end{enumerate}
\noindent\textbf{Answer.} Chunking chooses semantic units, size, overlap, and metadata boundaries for retrieval and context assembly.
\par\textit{Chunking chooses semantic units, size, overlap, and metadata boundaries for retrieval and context assembly. Tradeoff: Smaller chunks improve precision; larger chunks preserve context and reduce fragmentation. Watch for: Fixed-size splitting across tables, code, or headings destroys meaning.}
\paragraph{MCQ-0272 [hard]} What is the central engineering tradeoff of Chunking?
\begin{enumerate}[label=\Alph*.]
\item It can bridge terse questions to document style but may anchor retrieval on invented details.
\item It is narrower and more testable than global factuality but depends on context quality.
\item Smaller chunks improve precision; larger chunks preserve context and reduce fragmentation.
\item It separates retrieval granularity from generation context at extra storage and mapping cost.
\end{enumerate}
\noindent\textbf{Answer.} Smaller chunks improve precision; larger chunks preserve context and reduce fragmentation.
\par\textit{Chunking chooses semantic units, size, overlap, and metadata boundaries for retrieval and context assembly. Tradeoff: Smaller chunks improve precision; larger chunks preserve context and reduce fragmentation. Watch for: Fixed-size splitting across tables, code, or headings destroys meaning.}
\paragraph{MCQ-0273 [hard]} Which failure mode should an interviewer expect you to identify for Chunking?
\begin{enumerate}[label=\Alph*.]
\item Fixed-size splitting across tables, code, or headings destroys meaning.
\item Marking unsupported common knowledge as grounded inflates results.
\item Returning many overlapping parents wastes context and amplifies one source.
\item Passing the hypothetical text to generation as evidence confuses a retrieval aid with a source.
\end{enumerate}
\noindent\textbf{Answer.} Fixed-size splitting across tables, code, or headings destroys meaning.
\par\textit{Chunking chooses semantic units, size, overlap, and metadata boundaries for retrieval and context assembly. Tradeoff: Smaller chunks improve precision; larger chunks preserve context and reduce fragmentation. Watch for: Fixed-size splitting across tables, code, or headings destroys meaning.}
\paragraph{MCQ-0274 [medium]} A design review is specifically concerned with this risk: Fixed-size splitting across tables, code, or headings destroys meaning. Which topic should be reviewed first?
\begin{enumerate}[label=\Alph*.]
\item HyDE
\item Chunking
\item Parent-child retrieval
\item Faithfulness
\end{enumerate}
\noindent\textbf{Answer.} Chunking
\par\textit{Chunking chooses semantic units, size, overlap, and metadata boundaries for retrieval and context assembly. Tradeoff: Smaller chunks improve precision; larger chunks preserve context and reduce fragmentation. Watch for: Fixed-size splitting across tables, code, or headings destroys meaning.}
\paragraph{MCQ-0275 [hard]} Which design note most accurately balances the benefits and costs of Chunking?
\begin{enumerate}[label=\Alph*.]
\item Smaller chunks improve precision; larger chunks preserve context and reduce fragmentation.
\item It separates retrieval granularity from generation context at extra storage and mapping cost.
\item It is narrower and more testable than global factuality but depends on context quality.
\item It can bridge terse questions to document style but may anchor retrieval on invented details.
\end{enumerate}
\noindent\textbf{Answer.} Smaller chunks improve precision; larger chunks preserve context and reduce fragmentation.
\par\textit{Chunking chooses semantic units, size, overlap, and metadata boundaries for retrieval and context assembly. Tradeoff: Smaller chunks improve precision; larger chunks preserve context and reduce fragmentation. Watch for: Fixed-size splitting across tables, code, or headings destroys meaning.}
\paragraph{MCQ-0276 [medium]} Which explanation would be strongest in a system-design interview about Chunking?
\begin{enumerate}[label=\Alph*.]
\item Chunking chooses semantic units, size, overlap, and metadata boundaries for retrieval and context assembly.
\item Hypothetical Document Embeddings generate a likely answer passage and embed it as the retrieval query.
\item Small child chunks retrieve precisely while larger parent sections are supplied to the model.
\item Faithfulness asks whether claims in an answer are supported by the supplied context.
\end{enumerate}
\noindent\textbf{Answer.} Chunking chooses semantic units, size, overlap, and metadata boundaries for retrieval and context assembly.
\par\textit{Chunking chooses semantic units, size, overlap, and metadata boundaries for retrieval and context assembly. Tradeoff: Smaller chunks improve precision; larger chunks preserve context and reduce fragmentation. Watch for: Fixed-size splitting across tables, code, or headings destroys meaning.}
\paragraph{MCQ-0277 [hard]} Which production warning is most directly associated with Chunking?
\begin{enumerate}[label=\Alph*.]
\item Returning many overlapping parents wastes context and amplifies one source.
\item Marking unsupported common knowledge as grounded inflates results.
\item Fixed-size splitting across tables, code, or headings destroys meaning.
\item Passing the hypothetical text to generation as evidence confuses a retrieval aid with a source.
\end{enumerate}
\noindent\textbf{Answer.} Fixed-size splitting across tables, code, or headings destroys meaning.
\par\textit{Chunking chooses semantic units, size, overlap, and metadata boundaries for retrieval and context assembly. Tradeoff: Smaller chunks improve precision; larger chunks preserve context and reduce fragmentation. Watch for: Fixed-size splitting across tables, code, or headings destroys meaning.}
\paragraph{FC-0278 [medium]} Define Chunking in interview-ready language.
\noindent\textbf{Answer.} Chunking chooses semantic units, size, overlap, and metadata boundaries for retrieval and context assembly.
\par\textit{Chunking chooses semantic units, size, overlap, and metadata boundaries for retrieval and context assembly.}
\paragraph{FC-0279 [hard]} State the main tradeoff and production pitfall for Chunking.
\noindent\textbf{Answer.} Tradeoff: Smaller chunks improve precision; larger chunks preserve context and reduce fragmentation. Pitfall: Fixed-size splitting across tables, code, or headings destroys meaning.
\par\textit{Tradeoff: Smaller chunks improve precision; larger chunks preserve context and reduce fragmentation. Pitfall: Fixed-size splitting across tables, code, or headings destroys meaning.}
\paragraph{LA-0280 [hard]} Explain Chunking from first principles, then show how you would evaluate and operate it in production.
\noindent\textbf{Answer.} Start with the mechanism: Chunking chooses semantic units, size, overlap, and metadata boundaries for retrieval and context assembly. Make the tradeoff explicit: Smaller chunks improve precision; larger chunks preserve context and reduce fragmentation. Define workload-specific offline metrics and a latency/cost budget, test important slices, instrument the critical path, and create a rollback criterion. The production review must explicitly guard against this failure: Fixed-size splitting across tables, code, or headings destroys meaning.
\par\textit{A strong answer connects mechanism, measurable tradeoffs, failure isolation, observability, and rollback rather than listing product features.}
\paragraph{MCQ-0281 [medium]} Which statement best describes Parent-child retrieval?
\begin{enumerate}[label=\Alph*.]
\item Small child chunks retrieve precisely while larger parent sections are supplied to the model.
\item The system generates several query formulations and merges their result sets.
\item Packing selects, orders, deduplicates, and labels evidence within a token budget.
\item A router chooses retrievers, indexes, tools, or no-retrieval paths based on intent and risk.
\end{enumerate}
\noindent\textbf{Answer.} Small child chunks retrieve precisely while larger parent sections are supplied to the model.
\par\textit{Small child chunks retrieve precisely while larger parent sections are supplied to the model. Tradeoff: It separates retrieval granularity from generation context at extra storage and mapping cost. Watch for: Returning many overlapping parents wastes context and amplifies one source.}
\paragraph{MCQ-0282 [hard]} What is the central engineering tradeoff of Parent-child retrieval?
\begin{enumerate}[label=\Alph*.]
\item Specialization improves quality and cost but adds classification errors and operational branches.
\item Recall rises at the cost of additional searches, duplicates, and possible intent drift.
\item It separates retrieval granularity from generation context at extra storage and mapping cost.
\item Diversity and source coverage compete with local relevance and continuity.
\end{enumerate}
\noindent\textbf{Answer.} It separates retrieval granularity from generation context at extra storage and mapping cost.
\par\textit{Small child chunks retrieve precisely while larger parent sections are supplied to the model. Tradeoff: It separates retrieval granularity from generation context at extra storage and mapping cost. Watch for: Returning many overlapping parents wastes context and amplifies one source.}
\paragraph{MCQ-0283 [hard]} Which failure mode should an interviewer expect you to identify for Parent-child retrieval?
\begin{enumerate}[label=\Alph*.]
\item A route without confidence or fallback turns ambiguous queries into hard failures.
\item Treating all generated queries as equally trustworthy invites off-topic evidence.
\item Returning many overlapping parents wastes context and amplifies one source.
\item Greedy top-score packing often repeats one document and crowds out complementary facts.
\end{enumerate}
\noindent\textbf{Answer.} Returning many overlapping parents wastes context and amplifies one source.
\par\textit{Small child chunks retrieve precisely while larger parent sections are supplied to the model. Tradeoff: It separates retrieval granularity from generation context at extra storage and mapping cost. Watch for: Returning many overlapping parents wastes context and amplifies one source.}
\paragraph{MCQ-0284 [medium]} A design review is specifically concerned with this risk: Returning many overlapping parents wastes context and amplifies one source. Which topic should be reviewed first?
\begin{enumerate}[label=\Alph*.]
\item Parent-child retrieval
\item Context packing
\item Query routing
\item Multi-query retrieval
\end{enumerate}
\noindent\textbf{Answer.} Parent-child retrieval
\par\textit{Small child chunks retrieve precisely while larger parent sections are supplied to the model. Tradeoff: It separates retrieval granularity from generation context at extra storage and mapping cost. Watch for: Returning many overlapping parents wastes context and amplifies one source.}
\paragraph{MCQ-0285 [hard]} Which design note most accurately balances the benefits and costs of Parent-child retrieval?
\begin{enumerate}[label=\Alph*.]
\item It separates retrieval granularity from generation context at extra storage and mapping cost.
\item Recall rises at the cost of additional searches, duplicates, and possible intent drift.
\item Specialization improves quality and cost but adds classification errors and operational branches.
\item Diversity and source coverage compete with local relevance and continuity.
\end{enumerate}
\noindent\textbf{Answer.} It separates retrieval granularity from generation context at extra storage and mapping cost.
\par\textit{Small child chunks retrieve precisely while larger parent sections are supplied to the model. Tradeoff: It separates retrieval granularity from generation context at extra storage and mapping cost. Watch for: Returning many overlapping parents wastes context and amplifies one source.}
\paragraph{MCQ-0286 [medium]} Which explanation would be strongest in a system-design interview about Parent-child retrieval?
\begin{enumerate}[label=\Alph*.]
\item A router chooses retrievers, indexes, tools, or no-retrieval paths based on intent and risk.
\item The system generates several query formulations and merges their result sets.
\item Packing selects, orders, deduplicates, and labels evidence within a token budget.
\item Small child chunks retrieve precisely while larger parent sections are supplied to the model.
\end{enumerate}
\noindent\textbf{Answer.} Small child chunks retrieve precisely while larger parent sections are supplied to the model.
\par\textit{Small child chunks retrieve precisely while larger parent sections are supplied to the model. Tradeoff: It separates retrieval granularity from generation context at extra storage and mapping cost. Watch for: Returning many overlapping parents wastes context and amplifies one source.}
\paragraph{MCQ-0287 [hard]} Which production warning is most directly associated with Parent-child retrieval?
\begin{enumerate}[label=\Alph*.]
\item Greedy top-score packing often repeats one document and crowds out complementary facts.
\item A route without confidence or fallback turns ambiguous queries into hard failures.
\item Treating all generated queries as equally trustworthy invites off-topic evidence.
\item Returning many overlapping parents wastes context and amplifies one source.
\end{enumerate}
\noindent\textbf{Answer.} Returning many overlapping parents wastes context and amplifies one source.
